\documentclass[11pt,pdflatex,sn-mathphys-ay]{sn-jnl}% Math and Physical Sciences Numbered Reference Style
\usepackage{graphicx}%
\usepackage{multirow}%
\usepackage{lineno}
\usepackage{amsthm}
\theoremstyle{definition}
\usepackage{amsmath,amssymb,amsfonts}%
\usepackage{amsthm}%
\usepackage{tikz}
\usetikzlibrary{arrows.meta, positioning, calc}
% Submission History: 
% Dec 20, 2025,  Mathematical Methods of Operations Research
%% as per the requirement new theorem styles can be included as shown below
\theoremstyle{thmstyleone}%
\newtheorem{theorem}{Theorem}%  meant for continuous numbers
%%\newtheorem{theorem}{Theorem}[section]% meant for sectionwise numbers
%% optional argument [theorem] produces theorem numbering sequence instead of independent numbers for Proposition
\newtheorem{proposition}[theorem]{Proposition}% 
%%\newtheorem{proposition}{Proposition}% to get separate numbers for theorem and proposition etc.
\theoremstyle{thmstyletwo}%
\newtheorem{example}{Example}%
\newtheorem{lemma}{Lemma}%
\newtheorem{remark}{Remark}%
\theoremstyle{thmstylethree}%
\newtheorem{definition}{Definition}%
\raggedbottom
%%\unnumbered% uncomment this for unnumbered level heads
\usepackage{mathrsfs}%
\usepackage[title]{appendix}%
\usepackage{xcolor}%
\usepackage{textcomp}%
\usepackage{manyfoot}%
\usepackage{booktabs}%
\usepackage[table]{xcolor}
\usepackage{colortbl}
\newcommand{\fobj}[1]{f_{#1}}
\newcommand{\weight}[1]{w_{#1}}
\newcommand{\alert}[1]{\textcolor{red}{#1}}
\usepackage{algorithm}
\usepackage{algorithmic}
\usepackage{listings}%
\usepackage{subfigure}   % provides \subfloat
\usepackage{epsfig}
\usepackage{epstopdf}
\usepackage{hyperref} 
\usepackage{optidef}
\usepackage{enumitem}
\usepackage{bm}
% % Custom commands for math notation matching your python code
\newcommand{\R}{\mathbb{R}}
\newcommand{\norm}[1]{\left\lVert#1\right\rVert}
\newcommand{\omegaVec}{\boldsymbol{\omega}} % Combined vector [x, lam, u, v, mu]
\newcommand{\tb}{\textcolor{blue}}
\newcommand{\bb}[1]{\mathbf{#1}}
\newcommand{\tr}{\textcolor{red}}
\newcounter{cases}
\newcounter{subcases}[cases]
\newenvironment{mycases}
{%
	\setcounter{cases}{0}%
	\setcounter{subcases}{0}%
	\def\case
	{%
		\par\noindent
		\refstepcounter{cases}%
		\textbf{Case (\thecases).}
	}%
	\def\subcase
	{%
		\par\noindent
		\refstepcounter{subcases}%
		\textit{Subcase \thesubcases:}
	}%
}
{%
	\par
}
\renewcommand*\thecases{\alph{cases}}
\renewcommand*\thesubcases{\arabic{subcases}}
%%%%%%%%%%%%%%%%%%%%%%%%%%%%%%%%%%%%%%%%%%%%%%%%%%%%%%%%%%%%%%%%%%%
% Mark changes in revised manuscript
\usepackage[normalem]{ulem}
\usepackage{xcolor}
\definecolor{revisionblue}{RGB}{0,70,180}
\newcommand{\del}[1]{\textcolor{red}{\sout{#1}}}
\newcommand{\add}[1]{\textcolor{revisionblue}{#1}}
%%%%%%%%%%%%%%%%%%%%%%%%%%%%%%%%%%%%%%%%%%%%%%%%%%%%%%%%%%%%%%%%%%%

\geometry{left=1.8cm,right=1.8cm,top=2.2cm,bottom=2.2cm}
\begin{document}
	\linenumbers
	\title{A Homotopy Continuation Framework for Constrained Multiobjective Optimization}
	%\subtitle{Using  the  LaTex Template}
	\author[1]{Olaoluwa Ogunleye} 
	\author*[1]{Guangming Yao}\email{gyao@clarkson.edu}
	\author[2]{Jianhua Zhang} 
	\author[2]{Manh Bui} 
	\affil[1]{Department of Mathematics, Clarkson University, Potsdam, NY USA}
	\affil[2]{ Department of Electrical and Computer Engineering, Clarkson University, Potsdam, NY USA}
	
	
	\abstract{We develop a homotopy continuation framework for computing Karush--Kuhn--Tucker (KKT) points of constrained multiobjective optimization problems. The proposed homotopy map continuously deforms an easily solvable system into the KKT system associated with the original multiobjective problem, yielding a deterministic continuation trajectory that preserves the structure of the constrained optimization problem. 
		\del{Under mild regularity assumptions, global convergence of the homotopy path to Pareto-critical solutions is established for initial points chosen in the interior of the feasible region.}
		\add{Under the stated regularity and reference-point cone avoidance
			assumptions, global convergence of the homotopy path to
			Pareto-critical solutions is established for admissible initial
			points in the interior of the feasible region.} The framework is further extended with adaptive predictor--corrector continuation strategies to improve robustness for nonconvex, degenerate, and disconnected solution structures. Numerical experiments demonstrate that the proposed method remains effective even when initialized from nonfeasible points, indicating practical robustness beyond the theoretical guarantees. Comparative studies on benchmark multiobjective optimization problems evaluate the proposed framework against classical scalarization approaches and the evolutionary algorithm NSGA-II. Additional experiments on disconnected nonconvex problems and a multiobjective IEEE 5-bus optimal power flow problem illustrate the capability of the continuation framework to trace complex KKT-consistent structures arising in realistic applications. The results demonstrate that the proposed homotopy continuation framework provides an effective and interpretable deterministic alternative for constrained multiobjective optimization.}
	\keywords{Multiobjective Optimization, Homotopy Method, Global Convergence}
	\maketitle
	%%%%%%%%%%%%%%%%%%%%%%%%%%%%%%%%%%%%%%%%%%%%%%%%%%%%%%%%%%%%%
	\section{Multiobjective Optimization with Equality and Inequality Constraints}
	\label{sec:moo_formulation}
	Multiobjective optimization (MOO) provides a fundamental framework for decision-making problems involving multiple, often conflicting, objectives. Applications arise in a wide range of areas including engineering design~\cite{rao2009engineering}, economics~\cite{miettinen1999nonlinear}, machine learning~\cite{emmerich2018tutorial}, management, and operations research~\cite{cohon2013multiobjective}. Representative examples include investment project selection, production scheduling, and resource allocation problems. Comprehensive treatments of MOO and its economic interpretations can be found in~\cite{Charne1961,Fandel1980} and the references therein. The conceptual foundations of the field trace back to the pioneering works of~\cite{pareto1964cours,neumann1944,koopmans1951}, which established the basis of modern decision and optimization theory.
	%%%%%%%%%%%%%%%%%%%%%%%%%%%%%%%%%%%%%%%%%%%%%%%%%%%%%%%%%%%
	
	Let $\mathbb{R}^n$ be the $n$-dimensional Euclidean space, $\mathbb{R}^n_+ := \{ \bm{x} \in \mathbb{R}^n : x_i \ge 0, ~ i=1,\ldots,n \}$, $\mathbb{R}^n_{++} := \{ \bm{x} \in \mathbb{R}^n : x_i > 0, ~ i=1,\ldots,n \}$
	be the nonnegative and positive orthants of $\mathbb{R}^n$, respectively. 
	For any two vectors $\bm{y}=(y_1, y_2, \dots, y_n)$ and $\bm{z}=(z_1, z_2, \dots, z_n)$ in
	$\mathbb{R}^n,$ we adopt the standard component-wise notation: 
	\begin{center}
		$
		\begin{cases}
			\bm{y} = \bm{z} ~\text{ if and only if }~ y_i=z_i, \forall i=1, 2, \dots, n\\
			\bm{y} < \bm{z} ~\text{ if and only if }~  y_i<z_i, \forall i=1, 2, \dots, n\\
			\bm{y} \leqq \bm{z} ~\text{ if and only if }~  y_i \leq z_i, \forall i=1, 2, \dots, n\\
			\bm{y} \leq \bm{z} ~\text{ if and only if }~  y_i\leq z_i, \forall i=1, 2, \dots, n \text{ and } \bm{y}\neq \bm{z}.
		\end{cases}
		$
	\end{center}
	A MOO seeks to determine a vector of decision variables that simultaneously optimizes a set of conflicting objective functions. Conventionally, the canonical MOO is formulated as a minimization problem:
	%
	\begin{mini!}{\bm{x} \in \mathbb{R}^n}{\bm{f}(\bm{x})}{}{}\nonumber
		\addConstraint{\bm{g}(\bm{x}) \leq 0\tag{1}
			\label{MOO}}
		\addConstraint{\bm{h}(\bm{x}) = 0,}\nonumber
	\end{mini!}
	%
	\noindent  where $\bm{x}$ is the decision variable in the $n$-dimensional space, and $\add{\bm{f}} =(f_1, f_2, \dots, f_p)^T $ maps the decision space to the $p$-dimensional objective space. The feasible region $\Omega = \left\{ \bm{x} \in \mathbb{R}^n : \bm{g}(\bm{x}) \leq 0, \bm{h}(\bm{x}) = 0 \right\}$, is defined by a set of inequality constraint functions $g=(g_1, g_2, \dots, g_m)^T$ and equality constraint functions $h=(h_1,h_2,\dots,h_s)^T$. We assume that all \( g_i(\bm{x}) \), \( h_j(\bm{x}) \), and \( f_k(\bm{x}) \) are twice continuously differentiable functions, for indices \( i \in \{1, \ldots, m\} \),  \( j \in J:=\{1, \ldots, s\} \), and \( k \in \{1, \ldots, p\} \). \del{The interior of the feasible region} \add{The strictly inequality-feasible subset} is $\Omega^0 = \left\{ \bm{x} \in \mathbb{R}^n : g(\bm{x}) < 0, \ h(\bm{x}) = 0 \right\}$, and \del{the boundary of $\Omega$ } \add{we write $\partial\Omega=\Omega\setminus\Omega^0$
		for its relative boundary in this formulation.} 
	For any point $\bm{x}\in\Omega$, the active index set of inequality constraints is $I(\bm{x}) := \left\{ i \in \{1, \ldots, m\} : g_i(\bm{x}) = 0 \right\}.$ The set of strictly positive convex weights is defined as $W^{++} = \left\{ \bb{w} \in \mathbb{R}_{++}^p : \sum\limits_{i=1}^p \text{w}_i = 1 \right\}$. 
	
	In contrast to single-objective optimization, which typically yields a unique optimal solution,  a unique global optimum generally does not exist in a MOO due to the conflicting nature of the objectives. Instead, the solution concept is based on Pareto optimality~\cite{pareto1964cours, miettinen1999nonlinear}. \add{ For a sufficiently regular continuous MOP with $p$ objectives, the Pareto set can generically be expected to form a
		$(p-1)$-dimensional object, with a corresponding $(p-1)$-dimensional
		Pareto front in objective space; see, e.g., Hillermeier
		\cite{hillermeier2012nonlinear}. Degeneracy occurs when this expected regular
		structure breaks down, for example when the Pareto set or front has
		lower local dimension or contains singular or intersecting
		structures. Such degenerate configurations are particularly relevant
		for continuation methods because the regularity assumptions underlying
		standard path-following arguments may fail.}
	
	%
	\begin{definition}
		A feasible point $\bm{x}\in \Omega$ is said to be an efficient solution (or Pareto optimal) to problem \eqref{MOO} if there is no $\bm{y}\in \Omega$ such that 
		$\bm{f}(\bm{y}) \leq \bm{f}(\bm{x}).$
		The set of all Pareto optimal solutions in the decision space is called the Pareto set $\mathcal{P}$, and the mapping of the Pareto set into the objective space, $\bm{f}(\mathcal{P})$, is termed the Pareto front. 
	\end{definition}
	%
	In other words, $\bm{x}$ is an efficient solution if no objective function \( f_k \) can be improved without causing degradation in at least one other objective.
	It is well-known that if $\bm{x}$ is an efficient solution of \eqref{MOO} and some constraint qualifications, such as the Linear Independence Constraint Qualification or the Abadie Constraint Qualification holds, then the necessary conditions for Pareto optimality at $\bm{x}$ are given by the Karush-Kuhn-Tucker (KKT) conditions\,\cite{kuhn1951nonlinear}:
	%
	\begin{subequations} \label{KKT}
		\begin{align}
			\bb{J}_{\bm{f}}(\bm{x})^T \bb{w} + \bb{J}_{\bm{g}}(\bm{x})^T \bm{u} + \bb{J}_{\bm{h}}(\bm{x})^T \bm{v} &= 0, \label{eq:KKT_stationarity} \\
			U \bm{g}(\bm{x}) &= 0, \label{eq:KKT_complementarity} \\
			\bm{h}(\bm{x}) &= 0, \label{eq:KKT_feasibility}\\
			\sum_{i=1}^p \text{w}_i & = 1, \label{eq:KKT_c}
		\end{align}
	\end{subequations}
	%
	\noindent where \( \bb{w} \in \mathbb{R}_+^p \setminus \{0\} \) is the weight vector, \( \bm{u} \in \mathbb{R}_+^m \), \( \bm{v} \in \mathbb{R}^s \) are the Lagrange multipliers, \( U = \operatorname{diag}(u_1, \ldots, u_m) \) is the diagonal matrix formed from multipliers of the inequality constraints, $\bb{J}_{\bm{f}}(\bm{x})$ is the Jacobian matrix of ${\bm{f}}$ w.r.t $\bm{x}$, and the same applies to $\bb{J}_{\bm{g}}(\bm{x})$ and $\bb{J}_{\bm{h}}(\bm{x})$:
	\[
	\begin{aligned}
		\bb{J}_{\bm{f}}(\bm{x}) &= [\,\nabla f_1(\bm{x}), \dots, \nabla f_p(\bm{x})\,]^T 
		&&\in \mathbb{R}^{p \times n}, \\[4pt]
		\bb{J}_{\bm{g}}(\bm{x}) &= [\,\nabla g_1(\bm{x}), \dots, \nabla g_m(\bm{x})\,]^T 
		&&\in \mathbb{R}^{m \times n}, \\[4pt]
		\bb{J}_{\bm{h}}(\bm{x}) &= [\,\nabla h_1(\bm{x}), \dots, \nabla h_s(\bm{x})\,]^T 
		&&\in \mathbb{R}^{s \times n}.
	\end{aligned}
	\]
	A point \( \bm{x} \in \Omega \) that satisfies the KKT conditions is referred to as a \textit{KKT point} of the MOO problem. 
	The computation of KKT points, which characterize Pareto-critical trade-off solutions in constrained multiobjective optimization, therefore represents a central problem in this area. 
	Observe that system \eqref{KKT} encodes three conditions: \eqref{eq:KKT_stationarity} is the stationarity condition; \eqref{eq:KKT_complementarity} represents the complementarity condition for inequality constraints; and \eqref{eq:KKT_feasibility} enforces equality constraint satisfaction. These KKT conditions serve as the foundation for the homotopy-based framework developed in subsequent sections. Feasible solutions satisfying the KKT conditions are referred to as Pareto-critical points throughout this work.
	
	% %%%%%%%%%%%%%%%%%%%%%%%%%%%%%%%%%%%%%%%%%%%%%%%%%%%%%%%%%
	% The primary motivation of this work is to develop a homotopy continuation framework for MOO problems capable of robustly tracing Pareto-critical solution manifolds and exploring connected branches of Pareto-critical solutions in constrained multiobjective optimization problems.
	
	% The theoretical foundation of homotopy methods was established through the seminal works of~\cite{kellogg1976,smale1976, chow1978}, which laid the groundwork for their global convergence analysis. Since then, homotopy techniques have evolved into indispensable numerical tools for solving complementarity problems, variational inequalities, and nonlinear mathematical programming models~\cite{gowda1996, facchinei2003, shang2011, mccormick1989, monteroro1990, lin1997}. While classical algorithms such as Interior Point Methods (IPMs) and Sequential Quadratic Programming (SQP) remain the predominant techniques for solving KKT systems~\cite{nocedal2006numerical}, homotopy-based methods provide an attractive alternative for nonlinear and multiobjective formulations by continuously deforming a simple, well-posed system into the original KKT system~\eqref{KKT} through a smooth transformation known as a homotopy path~\cite{allgower1990numerical, watson2002applied}. In contrast to IPMs, homotopy approaches do not require strict feasibility conditions and are capable of tracing multiple solution branches~\cite{mehta2016numerical,nocedal2006numerical} without imposing strict interiority assumptions~\cite{mehta2016numerical, hauenstein2012polynomial}. This property makes them particularly suitable for addressing nonconvex and MOO problems~\cite{hauenstein2012polynomial, kruger2019continuation}. Furthermore, homotopy formulations naturally mitigate numerical issues related to degeneracy and ill-conditioning that often occur in complex constrained optimization problems~\cite{allgower2003introduction}. 
	% Their intrinsic path-tracking mechanism enables the exploration of complex solution structures and provides improved insight into feasible and Pareto-critical regions arising in nonlinear constrained optimization problems\,\cite{hauenstein2012polynomial, allgower1990numerical}, thereby offering enhanced robustness for nonconvex and multiobjective problems.
	
	% The modern development of homotopy methods in optimization began with the work of Garcia and Zangwill~\cite{Garcia1981}, who demonstrated the applicability of homotopy techniques to convex programming problems. Since then, homotopy and continuation methods have been widely studied in nonlinear programming, interior-point methods, and constrained optimization~\cite{Megiddo1988,Kojima1988,dong2023spline,zhou2018global}. In~\cite{lin1997}, the Combined Homotopy Interior Point (CHIP) method was introduced for computing KKT points of convex nonlinear programs by integrating Newton and linear homotopies. The CHIP framework was later extended to convex MOO problems~\cite{Lin2003,Maeda2004}. Subsequent work further generalized homotopy constructions under weakened constraint qualifications and normal cone conditions, establishing smooth and bounded continuation paths together with global convergence properties~\cite{song2008,YaoSong2013_AM,Zhao2012_JAM}. More recently, homotopy continuation techniques have also been incorporated into algorithms for challenging nonconvex optimization problems, including mixed-integer nonlinear programming and branch-and-bound frameworks~\cite{MaLi2022_AICHE}. Additional surveys and applied studies can be found in~\cite{ZhangEtAl2016_AIMS}.
	% Despite significant advances in homotopy-based optimization methods, many existing approaches rely on restrictive assumptions such as strong constraint qualifications, which may fail in complex constrained multiobjective optimization problems. 
	
	% To address these limitations, this work develops a generalized homotopy continuation framework for constrained multiobjective optimization under reference-point cone avoidance condition. The proposed framework computes KKT points through a continuation strategy that traces Pareto-critical solution branches and emonstrates practical convergence even when initialized from infeasible points.
	% Systematic numerical experiments are performed on benchmark, disconnected nonconvex, and power-system optimization problems. The proposed method is compared with classical scalarization techniques, including the weighted sum, $\epsilon$-constraint, global criterion, and lexicographic methods, as well as the evolutionary algorithm NSGA-II. The comparisons examine convergence behavior, computational cost, and Pareto-front coverage.
	% The main contributions of this work are summarized as follows:
	% \vspace{-0.3cm}
	% \begin{enumerate}
		% \item A generalized homotopy continuation framework for constrained MOO problems is developed.
		% \item Global convergence to KKT points is established under weakened assumptions. \add{Here, ``global convergence'' refers to the ability of the homotopy
			% continuation process to follow the solution path from an admissible
			% initial point toward a Pareto-critical (KKT) solution. It does not
			% mean convergence to a globally Pareto-optimal solution; the limiting
			% solution may be only locally Pareto optimal.}
		% \item Two numerical continuation algorithms for tracing Pareto-critical solution branches are developed and implemented.
		% \item The framework is demonstrated on benchmark, disconnected nonconvex, and power-system optimization problems.
		% \item Comparative studies with scalarization methods and NSGA-II are performed.
		% \end{enumerate}
	
	%%%%%%%%%%%%%%%%%%%%%%%%%%%%%%%%%%%%%%%%%%%%%%%%%%%%%%%%%
	The primary motivation of this work is to develop a homotopy continuation
	framework for MOO problems capable of
	\del{robustly tracing Pareto-critical solution manifolds and exploring
		connected branches of Pareto-critical solutions}
	\add{constructing continuation paths toward Pareto-critical (KKT)
		solutions}
	in constrained multiobjective optimization problems.
	
	The theoretical foundation of homotopy methods was established through
	the seminal works of~\cite{kellogg1976,smale1976, chow1978}, which laid
	the groundwork for their global convergence analysis. Since then,
	homotopy techniques have evolved into indispensable numerical tools for
	solving complementarity problems, variational inequalities, and nonlinear
	mathematical programming models~\cite{gowda1996, facchinei2003,
		shang2011, mccormick1989, monteroro1990, lin1997}. While classical
	algorithms such as Interior Point Methods (IPMs) and Sequential Quadratic
	Programming (SQP) remain the predominant techniques for solving KKT
	systems~\cite{nocedal2006numerical}, homotopy-based methods provide an
	attractive alternative for nonlinear and multiobjective formulations by
	continuously deforming a simple, well-posed system into the original KKT
	system~\eqref{KKT} through a smooth transformation known as a homotopy
	path~\cite{allgower1990numerical, watson2002applied}. In contrast to
	IPMs, homotopy approaches do not require strict feasibility conditions
	and are capable of tracing multiple solution branches
	~\cite{mehta2016numerical,nocedal2006numerical} without imposing strict
	interiority assumptions~\cite{mehta2016numerical,
		hauenstein2012polynomial}. This property makes them particularly suitable
	for addressing nonconvex and MOO problems
	~\cite{hauenstein2012polynomial, kruger2019continuation}. Furthermore,
	homotopy formulations naturally mitigate numerical issues related to
	degeneracy and ill-conditioning that often occur in complex constrained
	optimization problems~\cite{allgower2003introduction}.
	Their intrinsic path-tracking mechanism enables the exploration of
	complex solution structures and provides improved insight into feasible
	and Pareto-critical regions arising in nonlinear constrained optimization
	problems\,\cite{hauenstein2012polynomial, allgower1990numerical}, thereby
	offering enhanced robustness for nonconvex and multiobjective problems.
	
	\add{Continuation has also been used in MOO in a different, but closely
		related, manner: to move \emph{along} an already identified set of
		Pareto-critical points. Hillermeier's continuation framework exploits
		the local differential-geometric structure of the Pareto-critical set
		to generate neighboring Pareto-critical solutions. Related
		predictor--corrector approaches, including Pareto Tracer and Pareto
		Explorer\,\cite{martin2018pareto,schutze2020explorer}, likewise aim to explore the Pareto-critical set/front by
		continuing from known Pareto-critical solutions. 
		\add{Other continuation-based approaches include
			the path-following method of Recchioni
			\cite{recchioni2003path} for box-constrained
			multiobjective problems, the numerical integration
			method of Bolten et al.~\cite{bolten2021tracing}
			for tracing locally Pareto-optimal points, and
			the certified interval-based continuation approach
			of Martin et al.~\cite{martin2016continuation}
			for constrained bi-objective optimization.}
		Such approaches have
		been developed for constrained problems and for problems with more than
		two objectives, and subsequent extensions have also considered
		degenerate Pareto sets/fronts. In the latter case, the usual regular
		$(p-1)$-dimensional structure may break down, requiring special treatment
		of singular or lower-dimensional continuation directions\,\cite{schutze2025degenerate}.
		These methods should be distinguished from the primary use of homotopy
		in the present work: our homotopy parameter defines a deformation from
		an auxiliary system associated with an initial point \emph{toward} the
		KKT system of the original MOO problem. Thus, the continuation path is
		used primarily as a globalization mechanism for reaching a
		Pareto-critical point, rather than assuming a known Pareto-critical
		point and moving along the Pareto-critical set. The two viewpoints are
		therefore complementary.}
	
	The modern development of homotopy methods in optimization began with
	the work of Garcia and Zangwill~\cite{Garcia1981}, who demonstrated the
	applicability of homotopy techniques to convex programming problems.
	Since then, homotopy and continuation methods have been widely studied
	in nonlinear programming, interior-point methods, and constrained
	optimization~\cite{Megiddo1988,Kojima1988,dong2023spline,zhou2018global}.
	In~\cite{lin1997}, the Combined Homotopy Interior Point (CHIP) method was
	introduced for computing KKT points of convex nonlinear programs by
	integrating Newton and linear homotopies. The CHIP framework was later
	extended to convex MOO problems~\cite{Lin2003,Maeda2004}. Subsequent work
	further generalized homotopy constructions under weakened constraint
	qualifications and normal cone conditions, establishing smooth and
	bounded continuation paths together with global convergence properties
	~\cite{song2008,YaoSong2013_AM,Zhao2012_JAM}. More recently, homotopy
	continuation techniques have also been incorporated into algorithms for
	challenging nonconvex optimization problems, including mixed-integer
	nonlinear programming and branch-and-bound frameworks
	~\cite{MaLi2022_AICHE}. Additional surveys and applied studies can be
	found in~\cite{???????????????}.
	
	Despite significant advances in homotopy-based optimization methods,
	many existing approaches rely on restrictive assumptions such as strong
	constraint qualifications, which may fail in complex constrained
	multiobjective optimization problems.
	\add{At the same time, existing Pareto-set continuation methods provide
		powerful mechanisms for exploring a Pareto-critical set once a suitable
		point and regular local continuation structure are available. The
		present work addresses a different aspect of the problem by constructing
		a homotopy path from an auxiliary system toward the KKT system under the
		reference-point cone avoidance condition.}
	
	To address these limitations, this work develops a generalized homotopy
	continuation framework for constrained multiobjective optimization under
	\add{the} reference-point cone avoidance condition. The proposed
	framework computes KKT points through a continuation strategy that
	\del{traces Pareto-critical solution branches and demonstrates}
	\add{deforms an auxiliary system toward the KKT system and demonstrates}
	practical convergence even when initialized from infeasible points.
	Systematic numerical experiments are performed on benchmark, disconnected
	nonconvex, and power-system optimization problems. The proposed method
	is compared with classical scalarization techniques, including the
	weighted sum, $\epsilon$-constraint, global criterion, and lexicographic
	methods, as well as the evolutionary algorithm NSGA-II\,\cite{deb2002fast}. The comparisons
	examine convergence behavior, computational cost, and Pareto-front
	coverage.
	
	The main contributions of this work are summarized as follows:
	\vspace{-0.3cm}
	\begin{enumerate}
		\item A generalized homotopy continuation framework for constrained MOO
		problems is developed.
		
		\item Global convergence to KKT points is established under weakened
		assumptions. \add{Here, ``global convergence'' refers to the ability of
			the homotopy continuation process to follow the solution path from an
			admissible initial point toward a Pareto-critical (KKT) solution. It does
			not mean convergence to a globally Pareto-optimal solution; the limiting
			solution may be only locally Pareto optimal.}
		
		\item
		\del{Two numerical continuation algorithms for tracing Pareto-critical
			solution branches are developed and implemented.}
		\add{Two numerical continuation algorithms for following the proposed
			homotopy paths toward Pareto-critical solutions are developed and
			implemented.}
		
		\item The framework is demonstrated on benchmark, disconnected nonconvex,
		and power-system optimization problems.
		
		\item Comparative studies with scalarization methods and NSGA-II are
		performed.
	\end{enumerate}
	
	
	The remainder of the paper is organized as follows. Section~\ref{sec:related} presents the proposed homotopy continuation framework for constrained MOO, including the construction of the homotopy map and theoretical results on regularity, path boundedness, and global convergence to KKT points. Section~\ref{sec:alg} introduces the numerical continuation algorithms and implementation details underlying the proposed homotopy framework. In particular, Algorithm~1 presents a standard predictor--corrector continuation strategy for relatively smooth and connected Pareto-critical branches, while Algorithm~2 develops an enhanced pseudo-arclength continuation framework with adaptive step-size control, multi-step Newton correction, and branch-restarting mechanisms for handling nonconvex, degenerate, and disconnected solution structures. Section\,\ref{sec:numerical} presents computational experiments on benchmark, disconnected nonconvex, and multiobjective optimal power flow problems, together with comparisons against classical scalarization methods and the evolutionary algorithm NSGA-II. Finally, Section~\ref{sec:conclusion} summarizes the main findings and discusses directions for future research.
	
	
	%%%%%%%%%%%%%%%%%%%%%%%%%%%%%%%%%%%%%%%%%%%%%%%%%%%%%%%%%%%%%
	\section{Homotopy Theory and Continuation Framework for Multiobjective Optimization} 
	\label{sec:related}
	
	\subsection{Classical Multiobjective Optimization Approaches}
	MOO methods are commonly categorized as \textit{a priori}, \textit{a posteriori}, or interactive approaches according to how preference information is incorporated during the optimization process~\cite{miettinen1999nonlinear,coello2006evolutionary}. Classical scalarization techniques such as the weighted sum method\,\cite{marler2004survey, van2012evaluation,gass1955computational}, the $\epsilon$-constraint method\,\cite{haimes1971bicriterion}, the global criterion method\,\cite{marler2004survey}, and the lexicographic method\,\cite{isermann1982linear, miettinen1999nonlinear, ehrgott2005multicriteria} belong to the \textit{a priori} category, since they require preference parameters to be specified in advance. 
	\del{Early continuation and homotopy approaches for multiobjective optimization were developed by Hillermeier. 
		The present work extends this general direction to constrained many-objective settings with disconnected and degenerate Pareto-critical structures.} 
	
	\add{Continuation-based approaches to multiobjective optimization have also been developed for tracing Pareto-critical sets. In particular, Hillermeier~\cite{hillermeier2012nonlinear} interprets parameters of a scalarized multiobjective optimization problem, such as objective weights, as continuation parameters and follows the corresponding
		family of stationary solutions. In this setting, continuation proceeds
		along the set of KKT or Pareto-critical points as the scalarization
		parameters vary.}
	
	\add{The homotopy considered in the present work has a different purpose.
		Rather than using an objective weight as the homotopy parameter, we
		introduce an auxiliary parameter $t$ that continuously deforms an
		initially solvable system into the KKT system of the original
		constrained multiobjective problem. Thus, for a fixed initialization,
		the homotopy path is followed toward the KKT set as $t\to0$, rather
		than along the KKT set as an objective weight varies.} 
	
	In contrast, evolutionary multiobjective algorithms such as NSGA-II and SPEA2 approximate trade-off structures through population-based search without explicit scalarization~\cite{deb2002nsga2,zitzler2001spea2}. Representative interactive methods include the Step Method\,\cite{benayoun1971linear}, the Interactive Nondifferentiable Multiobjective Bundle-Based System\,\cite{miettinen2006synchronous}, Pareto Race\,\cite{korhonen2008pareto}, and reference-point or Interactive Weighted Chebyshev methods \cite{wierzbicki1980reference}. 
	
	%%%%%%%%%%%%%%%%%%%%%%%%%%%%%%%%%%%%%%%%%%%%%%%%%%%%%%%%%%%%%%
	\subsection{Assumptions and Regularity Conditions}
	For each $\bm x\in\Omega$, define the generalized normal cone\,\cite{rockafellar1998variational}
	\[
	\mathcal N(\bm{x})
	=
	\left\{
	\sum_{i\in I(\bm{x})}u_i\nabla g_i(\bm{x})
	+
	\sum_{j\in J}v_j\nabla h_j(\bm{x})
	:
	u_i\ge0,\ v_j\in\mathbb R
	\right\}.
	\]
	
	The following three assumptions are used throughout the rest of the paper:
	\begin{description}
		\item[(A)] The feasible set $\Omega$ is nonempty and bounded.
		\item[(B)] The active constraint gradients satisfy the linear independence constraint qualification; that is, for every $\bm{x}\in\Omega$, the set
		\[
		\{\nabla h_j(\bm{x}),\nabla g_i(\bm{x}): i\in I(\bm{x}),\ j\in J\}
		\]
		is linearly independent.
		\item[(C)] Reference-Point Cone Avoidance Condition:
		For the chosen initial point $\bm{x}^0\in\Omega^0$, the following hold.
		
		\begin{enumerate}
			\item[(i)] For every $\bm{x}\in\partial\Omega$,
			\[
			\bm{x}^0\notin \bm{x}+\mathcal N(\bm{x}).
			\]
			
			\item[(ii)] For every $\bm{x}\in\Omega^0$,
			\[
			\bm{x}^0\in \bm{x}+\mathcal N(\bm{x})
			\quad\Longrightarrow\quad
			\bm{x}=\bm{x}^0.
			\]
		\end{enumerate}
		
	\end{description}
	Assumption (C) is introduced to exclude spurious intersections between the chosen initialization point $\bm{x}^0\in\Omega^0$ and the affine cones generated by active constraint gradients. In particular, it guarantees that $\bm{x}^0$ cannot lie on a translated cone generated at a boundary point of the feasible region, while ensuring uniqueness of such an intersection inside the strictly feasible region. 
	\add{Assumption (C) is a geometric condition on the reference point
		$\bm{x}^0$ relative to the normal cones of the feasible set and is
		used as a sufficient condition in the global continuation analysis.
		It is not assumed to hold for arbitrary constraint representations or
		arbitrary nonconvex feasible sets. In particular, if $\Omega$ is an
		equality-constrained manifold, the equality multipliers are
		unrestricted in sign and $\mathcal N(\bm{x})$ contains the
		corresponding normal space. For example, if $\Omega$ is the unit
		circle represented by an equality constraint, a normal line through
		one point can intersect the circle at another point, so condition
		(C)(ii) need not hold.}
	
	\add{
		\paragraph{Example illustrating Assumption (C).}
		Assumption (C) can hold for problems involving both inequality and
		equality constraints. Consider, for example,
		\[
		\Omega
		=
		\left\{
		\bm{x}=(x_1,x_2,x_3)\in\mathbb{R}^3:
		g(\bm{x})=x_1^2+x_2^2-1\leq 0,\quad
		h(\bm{x})=x_3=0
		\right\},
		\]
		and choose the reference point $\bm{x}^0=(0,0,0).$
		Thus, $\Omega$ is the closed unit disk in the plane $x_3=0$.}
	\add{For an interior point $\bm{x}=(x_1,x_2,0)$, $x_1^2+x_2^2<1,$
		the inequality constraint is inactive. Since
		$\nabla h(\bm{x})=(0,0,1)^\top,$
		the normal cone reduces to the normal space generated by the equality
		constraint,
		\[
		\mathcal N(\bm{x})
		=
		\left\{
		v(0,0,1)^\top:v\in\mathbb{R}
		\right\}.
		\]
		Consequently,
		\[
		\bm{x}+\mathcal N(\bm{x})
		=
		\left\{
		(x_1,x_2,v)^\top:v\in\mathbb{R}
		\right\}.
		\]
		Hence
		\[
		\bm{x}^0\in\bm{x}+\mathcal N(\bm{x})
		\quad\Longrightarrow\quad
		x_1=x_2=0,
		\]
		and therefore $\bm{x}=\bm{x}^0$. This verifies condition (C)(ii)
		for interior points.}
	
	\add{Now let $\bm{x}\in\partial\Omega$, so that $x_1^2+x_2^2=1, x_3=0.$
		Both constraints are active, with 
		$\nabla g(\bm{x})=(2x_1,2x_2,0)^\top$, $\nabla h(\bm{x})=(0,0,1)^\top$.
		Therefore,
		\[
		\mathcal N(\bm{x})
		=
		\left\{
		u(2x_1,2x_2,0)^\top
		+
		v(0,0,1)^\top:
		u\geq0,\;v\in\mathbb{R}
		\right\}.
		\]
		If $\bm{x}^0\in\bm{x}+\mathcal N(\bm{x})$, then for some
		$u\geq0$ and $v\in\mathbb{R}$,
		$(0,0,0)^\top = \bigl((1+2u)x_1,(1+2u)x_2,v\bigr)^\top$.
		Since $x_1^2+x_2^2=1$, this would require
		$1+2u=0$, or equivalently $u=-\frac12$, 
		which contradicts $u\geq0$. Hence
		\[
		\bm{x}^0\notin\bm{x}+\mathcal N(\bm{x})
		\qquad
		\text{for every }\bm{x}\in\partial\Omega,
		\]
		and condition (C)(i) also holds.}
	
	\add{This example also illustrates why the presence of equality constraints
		does not by itself cause Assumption (C) to fail. The equality constraint
		generates the unrestricted normal direction $(0,0,v)^\top$, but this
		direction cannot cancel the radial displacement in the
		$(x_1,x_2)$-plane. At a boundary point, such cancellation would require
		a negative multiplier for the active inequality constraint, which is
		excluded by $u\geq0$.
	}
	
	\add{
		In contrast, if the feasible set is the unit circle itself,
		\[
		\Omega
		=
		\left\{
		(x_1,x_2)\in\mathbb{R}^2:
		h(x_1,x_2)=x_1^2+x_2^2-1=0
		\right\},
		\]
		then the multiplier associated with the equality constraint is
		unrestricted in sign. Thus, at a point $\bm{x}\in\Omega$,
		\[
		\mathcal N(\bm{x})
		=
		\{\,2v\bm{x}:v\in\mathbb{R}\,\},
		\]
		which is the full normal line. In particular, this line also passes
		through $-\bm{x}$, since
		\[
		-\bm{x}
		=
		\bm{x}+2(-1)\bm{x}.
		\]
		Hence a translated normal space can intersect the feasible set at a
		second point, and condition (C)(ii) may fail. This contrasts with the
		mixed inequality--equality example above and shows that Assumption (C)
		depends on the geometry of the combined normal cone rather than simply
		on the presence or absence of equality constraints.
	}
	
	\add{For inequality-defined regions, the contribution of an active
		inequality constraint is one-sided because $u_i\geq0$. For a fixed
		reference point $\bm{x}^0$, condition (C) can be checked by determining
		whether there exist prohibited feasible points $\bm{x}$ and
		multipliers satisfying
		\[
		\bm{x}^0-\bm{x}
		=
		\sum_{i\in I(\bm{x})}u_i\nabla g_i(\bm{x})
		+
		\sum_{j\in J}v_j\nabla h_j(\bm{x}),
		\qquad u_i\geq0.
		\]
		For simple constraint sets this condition can be checked analytically
		from the geometry of the active normal cones; for more complicated
		sets it can be formulated as an auxiliary feasibility problem over
		the possible active sets. Consequently, the convergence result below
		is conditional on (C) and does not constitute a convergence guarantee
		for arbitrary nonconvex or equality-constrained feasible sets.}
	
	\del{Geometrically, Assumption~(C) prevents affine combinations
		of active inequality constraint gradients and equality
		constraint gradients from re-entering a protected interior
		subset of the feasible region.}
	\add{Geometrically, Assumption~(C) imposes a restriction
		on combinations of inequality and equality constraint
		gradients, excluding certain configurations that would
		violate the prescribed normal-cone condition.} The condition is used only in the convergence analysis to exclude degenerate cone intersections arising in the homotopy equations and plays a key role in establishing boundedness and uniqueness properties of the continuation trajectories. While Assumption~(C) is imposed for the theoretical convergence analysis, the numerical experiments demonstrate that the proposed continuation algorithms remain effective for several degenerate and disconnected multiobjective optimization problems beyond the theoretical setting.
	
	\subsection{Homotopy Continuation Formulation}
	To solve constrained MOO problems, the proposed framework constructs a homotopy continuation structure whose solution trajectories can be followed continuously as the homotopy parameter varies. This continuation perspective enables the systematic tracing of Pareto-critical solutions and provides a deterministic framework for exploring complex multiobjective trade-off structures.
	To solve the KKT system~\eqref{eq:KKT_stationarity}-\eqref{eq:KKT_c} of \eqref{MOO}, 
	we construct the following homotopy map:
	\begin{equation} \label{homotopy} 
		H(\bm{\omega}^0, \bm{\omega}, t) =
		\begin{bmatrix}
			(1-t)\big(\bb{J}_{\bm{f}}(\bm{x})^T \bb{w} 
			+ \bb{J}_{\bm{g}}(\bm{x})^T \bm{u}\big)
			+ \bb{J}_{\bm{h}}(\bm{x})^T \bm{v}
			+ t(\bm{x}-\bm{x}^0) \\[2pt]
			\bm{h}(\bm{x}) \\[2pt]
			U\, \bm{g}(\bm{x}) - t\,U^0 \bm{g}(\bm{x}^0) \\[2pt]
			(1-t)\!\left(1-\sum\limits_{i=1}^p \text{w}_i\right)e 
			- t\!\left(\bb{w}^{3/2} - (\bb{w}^0)^{3/2}\right)
		\end{bmatrix}
		= 0,
	\end{equation}
	%
	\noindent
	where $\bm{\omega}^0 := (\bm{x}^0, \bb{w}^0, \bm{u}^0, \bm{v}^0=0) \in N$, $\bm{\omega} = (\bm{x}, \bb{w}, \bm{u}, \bm{v}) \in M$,  and $t \in [0,1]$, in which 
	\[
	N = \Omega^0 \times W^{++} \times \mathbb{R}^m_{++} \times \mathbb{R}^s, 
	\qquad 
	M = \Omega \times \mathbb{R}^{p+m}_{+} \times \mathbb{R}^s,
	\qquad
	H: N \times M \times (0,1] \to \mathbb{R}^{n+p+m+s}.
	\]
	%
	\noindent 
	The unit vector $\bm{e} = (1, 1, \dots, 1)^T \in \mathbb{R}^p$, 
	and the fractional power vector $\bb{w}^{3/2} = (\text{w}_1^{3/2}, \dots, \text{w}_p^{3/2})^T$  are introduced to avoid trivial solutions.  
	%%%%%%%%%%%%%%%%%%%%%%%%%%%%%%%%%%%%%%%%%%%%%%%%%%%%%%%%%%
	\subsection{Initialization and Limiting Behavior}
	At \(t=1\), the homotopy system reduces to
	\begin{subequations}
		\label{eq:t1_system}
		\begin{align}
			\mathbf{J}_{\bm h}(\bm{x})^T \bm{v} + \bm{x}-\bm{x}^0 &= 0, \label{eq:t1a}\\
			\bm{h}(\bm{x}) &= 0, \label{eq:t1b}\\
			U\bm g(\bm{x}) - U^{0}\bm g(\bm{x}^{0}) &= 0, \label{eq:t1c}\\
			\bb{w}^{3/2} &= (\bb{w}^{0})^{3/2}. \label{eq:t1d}
		\end{align}
	\end{subequations}
	From \eqref{eq:t1a}, we have
	\[
	\bm{x}^0=\bm{x}+\mathbf J_{\bm h}(\bm{x})^T\bm v.
	\]
	Since $\bm{x}\in\Omega$, there are two cases. If $\bm{x}\in\partial\Omega$, then $\bm{x}^0$ lies in the translated cone generated by the equality gradients and active inequality gradients at a boundary point, contradicting Assumption~(C)(i). Hence $\bm{x}\in\Omega^0$.
	
	For $\bm{x}\in\Omega^0$, we have $I(\bm{x})=\emptyset$, and therefore
	\[
	\bm{x}^0\in \bm{x}+\mathcal N(\bm{x}).
	\]
	By Assumption~(C)(ii), $\bm{x}=\bm{x}^0$. Substituting this into \eqref{eq:t1a} gives
	\[
	\mathbf J_{\bm h}(\bm{x})^T\bm v=0.
	\]
	By Assumption~(B), $\bm v=\bm 0$.
	
	As $t \to 0$, the system continuously deforms into the KKT conditions\,\eqref{eq:KKT_stationarity}-\eqref{eq:KKT_c} of \eqref{MOO}. 
	Thus, following the continuation trajectory from \(t=1\) to \(t=0\) yields a KKT point of the original constrained MOO problem.
	
	\subsection{Global Convergence Analysis}
	Next, we introduce $H_{\bm{\omega}^{0}}(\bm{\omega}, t)$ and zero set $H_{\bm{\omega}^{0}}^{-1}(0)$ and establish a collection of lemmas and theorems demonstrating that the proposed homotopy map guarantees global convergence under Assumptions (A), (B), and (C). For a given initial point \(\bm{\omega}^{0}\), we define
	\[
	H_{\bm{\omega}^{0}}(\bm{\omega}, t) := H(\add{\bm{\omega}^{0}, \bm{\omega},} t),
	\]
	and the associated zero set
	\[
	H_{\bm{\omega}^{0}}^{-1}(0) 
	= 
	\Big\{ 
	(\bm{\omega}, t) \in 
	\Omega \times \mathbb{R}_{+}^{p+m} \times \mathbb{R}^s \times (0,1] 
	~\big|~ 
	H(\add{\bm{\omega}^{0}, \bm{\omega},} t) = 0 
	\Big\}.
	\]
	
	%%%%%%%%%%%%%%%%%%%%%%%%%%%%%%%%%%%%%%%%%%%%%%%%%
	%%%%%%%%%%%%%%%%%%%%%%%%%%%%%
	\begin{definition} 
		\label{def: regularvalue}\cite{LeeSmoothManifolds2013}
		Let $M, N$ be differential manifolds with $\text{dim} N=p$ and let
		$H : M \rightarrow N$ be a differentiable mapping. If 
		$$\text{rank} \left[\frac{\partial H(\bm{x})}{\partial x}\right]=p,\qquad\qquad\forall x\in H^{-1}(y),$$
		we say that $y\in N$ is a regular value of $H$ and $x\in M$ is a
		regular point. Given a curve $\Gamma\subset H^{-1}(y),$ if every
		$x\in \Gamma$ is a regular point, then we say that $\Gamma$ is a
		regular path.
	\end{definition}
	%%%%%%%%%%%%%%%%%%%%%%%%%%%%%
	\begin{definition}[Diffeomorphism] \label{def_diffeomo}\cite{LeeSmoothManifolds2013}
		Given two manifolds $M$ and $N$, a differentiable map $ {\displaystyle f\colon M\rightarrow N}$ is called a diffeomorphism if it is a bijection and its inverse ${\displaystyle f^{-1}\colon N\rightarrow M}$ is differentiable as well. If these functions are  $r$ times continuously differentiable, $f$ is called a $C^{r}-$ diffeomorphism.
	\end{definition}
	In other words, a diffeomorphism is an isomorphism of smooth manifolds. It is an invertible function that maps one differentiable manifold to another such that both the function and its inverse are continuously differentiable.
	%%%%%%%%%%%%%%%%%%%%%%%%%%
	% \begin{definition} [Archlength] \label{def_archlength}
		% The definition of arclength of a curve
		% \end{definition}
	%%%%%%%%%%%%%%%%%%%%%%%%%%%%%%
	\begin{lemma}[The Parametric Form
		of the Sard Theorem on a Manifold with Boundary]\label{Sard}\cite{sard1942measure,LeeSmoothManifolds2013} Let
		$ W, \text{and}~ N$ be differential manifolds of dimension $q$
		and $p$, respectively, and let $M$ be a $m-$dimensional differential
		manifold with boundary. Suppose   $F :  W\times M\rightarrow N$
		is a $C^r$ mapping, where $r> max \{0, m-p\}.$ If $ 0\in N$ is a
		regular value of $F$ and $\partial F,$ then for almost all
		$\text{w}\in  W, 0$ is a regular value of
		$F_{\text{w}}=F(\text{w}, \cdot)$ and $\partial{F_{\text{w}}}$, where
		$\partial F, ~\partial{F_{\text{w}}}$ denote the restriction of $F$
		and $F_{\text{w}}$ to $ W\times\partial M,$ and $
		\partial M,$ respectively.
	\end{lemma}
	
	Lemma\,\ref{Sard} follows from the parametric transversality theorem, since Sard’s theorem is obtained as the special case in which transversality to a point implies that the set of parameters for which the slice map fails to have that point as a regular value has measure zero; see\,\cite[Theorem~6.35]{LeeSmoothManifolds2013}.
	%%%%%%%%%%%%%%%%%%%%%%%%%%%
	\begin{lemma}[The Inverse Image Theorem]\label{inverse}\cite{LeeSmoothManifolds2013} Suppose $M$ is an
		$m-$dimensional $C^r$ differential manifold with boundary, $N$ is a
		$p-$dimensional $C^r$ differential manifold, $r\geqslant 1,$ and $ F
		: M\rightarrow N$ is a $C^r$ map. If $q\in N$ is a regular value of
		$F$ and $\partial F,$ then either $S=F^{-1}(q)$ is empty or a
		$(m-p)-$dimensional submanifold, and $\partial S=S\cap
		\partial M.$
	\end{lemma}
	The identity $\partial F^{-1}(q)=F^{-1}(q)\cap\partial M$ follows from the regularity of $q$ as a value of $F$. In particular, surjectivity of $DF_x$ implies that $F$ is transverse to $\{q\}$ along $\partial M$, so the preimage intersects the boundary cleanly and its boundary is precisely the intersection with $\partial M$; see \cite[Theorem~5.12]{LeeSmoothManifolds2013}.
	
	%%%%%%%%%%%%%%%%%%%%%%%%%%%%%%%%%%%%%%%%%%%%%%%%%
	\begin{lemma}[Classification of One-Dimensional Manifolds with Boundary]
		\label{fenglei}
		\cite[Theorem~1.8]{LeeSmoothManifolds2013}
		\del{Each connected component of a one-dimensional
			manifold with boundary is homeomorphic to
			either the unit circle $S^1$ or a unit interval.}
		\add{Let $M$ be a connected, Hausdorff, second-countable
			one-dimensional smooth manifold with boundary.
			Then $M$ is diffeomorphic to one of the following:
			\[
			S^1,\qquad
			\mathbb{R},\qquad
			[0,\infty),\qquad
			[0,1].
			\]
			In particular, if $M$ is compact, then it is
			diffeomorphic to either $S^1$ or $[0,1]$.}
	\end{lemma}
	\begin{proof}
		\del{This follows from the classification theorem for
			one-dimensional manifolds;}
		\add{The result follows from the classification of connected
			one-dimensional smooth manifolds with boundary.
			The compact cases are $S^1$ and $[0,1]$,
			whereas the noncompact cases are $\mathbb{R}$
			and $[0,\infty)$;} see
		\cite[Theorem~1.8]{LeeSmoothManifolds2013}.
	\end{proof}
	%%%%%%%%%%%%%%%%%%%%%%%%%%%%%%%%%%%%%%%%%%%%%%%%%
	
	
	%%%%%%%%%%%%%%%%%%%%%%%%
	% \begin{lemma}[Closure of $\Omega$ and existence of $\hat{\Omega}$]
		% \label{lem:closed}
		% The set $\Omega$ is closed, and there exists an open set $\hat{\Omega}$ such that $\hat{\Omega} \subset \Omega^{0}.$
		% \end{lemma}
	
	% \begin{proof}
		% Since $\Omega$ is defined by closed conditions, it is closed. Because $\Omega^{0}$ is open by definition, it contains an open neighborhood $\hat{\Omega}$.
		% \end{proof}
	
	%%%%%%%%%%%%%%%%%%%%%%
	\begin{lemma}
		\label{lem:closed}
		Suppose that $\bm g$ and $\bm h$ are continuous. Then the feasible set
		$\Omega=\left\{\bm x\in\mathbb R^n:
		\bm g(\bm x)\le0,\ \bm h(\bm x)=0\right\}$
		is closed. Moreover, the strictly feasible set
		$\Omega^0=
		\left\{\bm x\in\mathbb R^n:
		\bm g(\bm x)<0,\ \bm h(\bm x)=0\right\}$
		is relatively open in the equality-constraint manifold
		$\mathcal M=\{\bm x\in\mathbb R^n:\bm h(\bm x)=0\}.$
	\end{lemma}
	
	\begin{proof}
		Since each component of $\bm g$ and $\bm h$ is continuous, the sets
		\[
		\{\bm x:g_i(\bm x)\le0\}
		\quad\text{and}\quad
		\{\bm x:h_j(\bm x)=0\}
		\]
		are closed. Hence $\Omega$, being a finite intersection of closed sets, is closed.
		The set $\{\bm x:g_i(\bm x)<0\}$
		is open for each $i$. Therefore, $\Omega^0 = \mathcal M \cap \bigcap_{i=1}^m \{\bm x:g_i(\bm x)<0\}$, 
		which is relatively open in $\mathcal M$.
	\end{proof}
	%%%%%%%%%%%%%%%%%%%%%%%%%%%%%%%%%%%%%%%%%%%%%%%%%%%%%
	\begin{lemma}\label{zhengze}
		Suppose that $\Omega \ne \emptyset$ and conditions (A), (B)
		hold. Then for almost all initial points $\bm{\omega}^0\in
		\Omega^0\times W^{++}\times
		\mathbb{R}^{m}_{++}\times\{0\}, ~0$ is a regular value of $H_{\bm{\omega}^0},$ and $H_{\bm{\omega}^0}^{-1}(0)$
		consists of some smooth curves. Among them, a smooth curve, say
		$\Gamma_{\bm{\omega}^0},$  starts from $(\bm{\omega}^0, 1).$
	\end{lemma}
	\vspace{-0.5cm}
	%----------------------------------------------------
	\begin{proof}
		Observe that \begin{eqnarray*}
			&& H^{-1}_{\bm{\omega}^0}(0)  =  \{(\bm{\omega},t) \in \Omega \times
			\mathbb{R}^{p+m}_{+} \times \mathbb{R}^s \times (0, 1] \colon \quad H(\bm{\omega}^0, \bm{\omega},t) = 0\}.
		\end{eqnarray*}
		
		For any $ (\bm{\omega}^0, \bm{\omega},  t)\in H^{-1}(0),$
		$$\frac{\partial H( \bm{\omega}^0, \bm{\omega}, t)}{\partial(\bm{x}^0, \textbf{w}^0, \bm{u}^0, \bm{x})}=
		\begin{matrix}
			\begin{bmatrix}
				-t I & 0 & 0 & Q(\bm{x}) \\
				0 & 0 & 0 & \textbf{J}_ h(\bm{x})\\
				\add{-t U^0\textbf{J}_{\bm{g}}(\bm{x}^0)}
				& 0 & -t \text{diag} (\bm{g}(\bm{x}^0)) & U \textbf{J}_ g(\bm{x})\\
				0 & \frac{3}{2}t \text{diag}((\textbf{w}^0)^{\frac{1}{2}}) & 0 & 0 
			\end{bmatrix}
		\end{matrix},$$
		where 
		\[Q(\bm{x})=(1-t)(\sum^p\limits_{i=1}\text{w}_i\nabla^2f_i(\bm{x})
		+\sum^m\limits_{j=1}u_j\nabla^2g_j(\bm{x}))+\sum\limits_{k=1}^sv_k\nabla
		^2 h_k(\bm{x})+t I.\]
		Because $\bm{x}\in \Omega, ~ \bm{x}^0 \in \Omega^0 $ and $t\in (0, 1],$ by Assumption (B), we obtain that 
		\[
		\text{rank}\left[\frac{\partial
			H(\bm{\omega}^0, ~\bm{\omega}, ~t)}{\partial(\bm{x}^0, \textbf{w}^0,
			\bm{u}^0, x)}\right]=n+p+m+s.
		\]
		So that 
		\[\text{rank}\left[\frac{\partial
			H(\bm{\omega}^0, ~\bm{\omega}, ~t)}{\partial(\bm{\omega}^0, ~\bm{\omega}, ~ t)}\right] = \text{rank}\left[\frac{\partial
			H(\bm{\omega}^0, ~\bm{\omega},~ , 1)}{\partial(\bm{\omega}^0, ~\bm{\omega})}\right]=n+p+m+s.
		\]
		Thus both Jacobian matrices are full-row rank. This means that, $0$ is a regular value of $H$ and
		$\partial H$ from Definition \ref{def: regularvalue}, where $\partial H$ denotes the restriction of $H$ to $N \times \partial (M \times (0,1]).$ By Lemma \ref{Sard}, for almost all $\bm{\omega}^0 \in
		\Omega^0\times W^{++}\times
		\mathbb{R}^{m}_{++}\times\{0\}=N, ~0$ is a regular value of $~H_{\bm{\omega}^0}$ and $~\partial H_{\bm{\omega}^0}$, where $~\partial H_{\bm{\omega}^0}$ denotes the restriction of $H_{\bm{\omega}^0}$ to $\partial (M \times (0,1]).$
		
		\noindent Since $~H(\bm{\omega}^0, \bm{\omega}^0, 1)=0$, then by Lemma\,\ref{inverse} we have that $H_{\bm{\omega}^0}^{-1}(0)$ cannot be empty, that is, it is a $1$-dimensional submanifold and by Lemma\,\ref{fenglei} consists of
		some smooth curves. Also, since $~H(\bm{\omega}^0, \bm{\omega}^0,
		1)=0,$ there must be a smooth curve,  denoted by $\Gamma_{\bm{\omega}^0},$ that starts from $(\bm{\omega}^0, 1) $.
	\end{proof}
	
	%%%%%%%%%%%%%%%%%%%%%%%%%%%%%%%%%%%%%%%%%%%%%%%%%%%%%%%
	%%%%%%%%%%%%%%%%%%%%%%%%%%%%%%%%%%%%%%%%%%%%%%%%%
	\begin{lemma}\label{fenliangyoujie}
		Suppose that $\Omega \ne \emptyset$ and conditions (A), (B)
		hold. For a given $\bm{\omega}^0\in \bm{\omega}^0\times  W^{++} \times
		\mathbb{R}^{m}_{+}\times \mathbb{R}^s,$ if ~$  0 $ is a regular
		value of $H_{\bm{\omega}^0}$, then the $\text{w}$ component of the smooth curve
		$\Gamma_{\bm{\omega}^0}$ is bounded.
	\end{lemma}
	\vspace{-0.5cm}
	%%%%%%%%%%%%%%%%%%%%%%%%%%%%%%%%%%%%%%%%%%%%%%%%%
	\begin{proof} Suppose that the conclusion does not hold, that is, \textit{the $\text{w}$ component of the smooth curve
			$\Gamma_{\bm{\omega}^0}$ is unbounded}. Since (0,1]
		is bounded, there exists a sequence
		$\{(\bm{\omega}^k,t_k)\}\subset\Gamma_{\bm{\omega}^0}$
		such that,
		$$  t_k\rightarrow t_*, ~
		\text{and}~ \|\textbf{w} ^k\|\rightarrow +\infty,~~k\rightarrow
		\infty.$$
		
		\noindent From the last row of the homotopy equation (\ref{homotopy}), we have
		\begin{equation}
			(1-t_k)(1-\sum_{i=1}^p \text{w}_i^k)e-t_k [(\textbf{w}^{k})^{\frac{3}{2}} -(\textbf{w}^0)^{\frac{3}{2}}] =0 . \label{4d}
		\end{equation}
		That is
		\begin{equation}\label{w}
			\begin{split}
				\begin{matrix}
					\begin{bmatrix}
						1-t_k\\
						1-t_k\\
						\vdots\\
						1-t_k
					\end{bmatrix}
				\end{matrix}-
				\begin{matrix}
					\begin{bmatrix}
						(1-t_k)\text{w}_1^{k}+(1-t_k)\sum\limits_{i \neq 1}\text{w}_i^{k} + t_k (\text{w}_1^{k})^\frac{3}{2}\\
						(1-t_k)\text{w}_2^{k}+(1-t_k)\sum\limits_{i\neq2}\text{w}_i^{k} + t_k (\text{w}_2^{k})^\frac{3}{2}\\
						\vdots\\
						(1-t_k)\text{w}_p^k+(1-t_k)\sum\limits_{i\neq p}\text{w}_i^{k} + t_k (\text{w}_p^{k})^\frac{3}{2}
					\end{bmatrix}
				\end{matrix}+t_k\begin{matrix}
					\begin{bmatrix}
						(\text{w}^0_1)^\frac{3}{2}\\
						(\text{w}^0_2)^\frac{3}{2}\\
						\vdots\\
						(\text{w}^0_p)^\frac{3}{2}
					\end{bmatrix}
				\end{matrix}
				=0. \end{split} \end{equation}
		
		\noindent \del{Let
			$I=\{j\in\{1,2,\dots,p\}:
			\lim\limits_{k\rightarrow\infty}\text{w}_j^k=\infty\}.$
			The assumption $\|\textbf{w}^k\|\rightarrow+\infty$
			implies $I\neq\emptyset$.
			Since $t_k\rightarrow t_*\in[0,1]$
			and $\textbf{w}^k\geq0$, it follows that
			the second part in the left-hand side of
			some equation of (\ref{w}) tends to infinity
			as $k\rightarrow+\infty$, but the first
			and third parts are bounded.
			This is a contradiction. Thus the projection
			of the smooth curve $\Gamma_{\bm{\omega}^0}$
			on the $\text{w}$ component is bounded.
		}
		
		\noindent 
		\add{
			\noindent Summing the $p$ equations in
			(\ref{w}), we obtain
			\[
			p(1-t_k)\sum_{i=1}^p\text{w}_i^k
			+t_k\sum_{i=1}^p(\text{w}_i^k)^{3/2}
			=
			p(1-t_k)
			+t_k\sum_{i=1}^p(\text{w}_i^0)^{3/2}.
			\]
			Since $\|\textbf{w}^k\|\rightarrow+\infty$
			and $\textbf{w}^k\geq0$, both
			$\sum_{i=1}^p\text{w}_i^k$ and
			$\sum_{i=1}^p(\text{w}_i^k)^{3/2}$
			tend to $+\infty$.
			As $t_k\rightarrow t_*\in[0,1]$,
			the left-hand side tends to $+\infty$,
			whereas the right-hand side remains bounded.
			This is a contradiction. Thus the projection
			of the smooth curve $\Gamma_{\bm{\omega}^0}$
			onto the $\textbf{w}$ component is bounded.
		}
		
	\end{proof}  
	
	We next prove that $\Gamma_{\bm{\omega}^0}$ is a bounded curve.
	
	%%%%%%%%%%%%%%%%%%%%%%%%%%%%%%%%%%%%%%%%%%%%%%%%%%%%%%%%
	\begin{lemma}\label{youjie} Suppose that the conditions (A), (B), and (C)
		hold. Then for a given $\bm{\omega}^0\in \hat{\Omega}\times  W^{++} \times
		\mathbb{R}^{m}_{+}\times \mathbb{R}^s,$ if $0$ is a regular value
		of $H_{\bm{\omega}^0}$, then $\Gamma_{\bm{\omega}^0}$ is a bounded curve.
	\end{lemma}
	\vspace{-0.5cm}
	%%%%%%%%%%%%%%%%%%%%%%%%%%%%%%%%%%%%%%%%%%%%%%%%%%%%%%%%%%%%%%
	\begin{proof}   Assume that $\Gamma_{\bm{\omega}^0}$ is an unbounded curve. Since $\Omega$ is bounded then by  Bolzano–Weierstrass theorem and \add{Lemma}\,\ref{fenliangyoujie}, there exists a sequence
		$\{(\bm{\omega}^k,t_k)\}\subset \Gamma_{\bm{\omega}^0}$ such that as $k\rightarrow\infty$, 
		\begin{equation*}
			\bm{x}^k\rightarrow \bm{x}^* \in \Omega,\quad t_k\rightarrow t_* \in [0, 1], \quad \textbf{w} ^k\rightarrow \textbf{w}^* \in \mathbb{R}^{p}_{+},\quad
			\text{and} \quad \|(\bm{u}^k,\bm{v}^k)\|\rightarrow +\infty.
		\end{equation*}
		
		\noindent From the homotopy equation (\ref{homotopy}), we have
		\begin{eqnarray}
			(1-t_k)(\textbf{J}^T_ f(\bm{x}^k) \textbf{w}^k +\textbf{J}^T_ g(\bm{x}^k)\bm{u}^k)+\textbf{J}^T_
			h(\bm{x}^k)\bm{v}^k+t_k (\bm{x}^k-\bm{x}^0) & =0, \label{4a}\\
			h(\bm{x}^k) & =0, \label{4b}\\
			\bm{u}^k g(\bm{x}^k)-t_k \bm{u}^0 g(\bm{x}^0) &=0. \label{4c}
		\end{eqnarray}
		%
		\noindent Notice that $\bm{x}^k\rightarrow \bm{x}^* \in  \Omega$ means that $\bm{x}^k\rightarrow \bm{x}^* \in \Omega^0$ or $\bm{x}^k\rightarrow \bm{x}^* \in \partial \Omega$. First, suppose that $\bm{x}^k\rightarrow \bm{x}^* \in \partial \Omega$. Let 
		\begin{equation}
			J^*=\{j\in J : \lim\limits_{k\rightarrow \infty} v_j^k=\infty\}, 
			I^*=\{i\in I : \lim\limits_{k\rightarrow\infty} u_i^k=\infty\}.
		\end{equation}
		
		\noindent Recall that $I(\bm{x}) = \{i\in I :g_i(\bm{x}) =0\add{\}}$. From (\ref{4c}) since  $\lim\limits_{k\rightarrow \infty} u_i^k g_i(\bm{x}^k) < \infty$, if $\lim\limits_{k\rightarrow \infty} u_i^k = \infty$ then $\lim\limits_{k\rightarrow \infty} g_i(\bm{x}^k) =  g_i(\bm{x}^*) = 0$. This means that $ I^*\subset I(\bm{x}^*).$ This fact will be used consequently below.
		We are going to establish a contradiction that these two sets are empty, that is $J^* = \emptyset$ and $I^* = \emptyset$.
		
		\noindent Suppose $J^*\ne \emptyset.$ Rewrite (\ref{4a}) as
		\begin{eqnarray}
			(1-t_k)\left[\textbf{J}^T_ f(\bm{x}^k)\textbf{w}^k+\sum_{i\not\in I^*}u_i^k\textbf{J}_ {g_i}(\bm{x}^k)\right]  + t_k \left(\bm{x}^k-\bm{x}^0\right) + \left[\textbf{J}^T_ h(\bm{x}^k)\bm{v}^k+
			(1-t_k)\sum_{i\in I^*}u_i^k\textbf{J}_ {g_i}(\bm{x}^k) \right] =0.
		\end{eqnarray}
		
		\noindent Since
		$J^*\ne \emptyset$ and Assumption (B) holds, as $k\rightarrow \infty,$
		\begin{eqnarray}
			\underbrace{(1-t_k)\left[\textbf{J}^T_ f(\bm{x}^k)\textbf{w}^k+\sum_{i\not\in I^*}u_i^k\textbf{J}_ {g_i}(\bm{x}^k)\right]}_\text{bounded} + \underbrace{t_k \left(\bm{x}^k-\bm{x}^0\right)}_\text{bounded} + \underbrace{\left[\textbf{J}^T_ h(\bm{x}^k)\bm{v}^k+
				(1-t_k)\sum_{i\in I^*}u_i^k\textbf{J}_ {g_i}(\bm{x}^k) \right]}_\text{unbounded} =0.
		\end{eqnarray} 
		This is a contradiction. Thus $J^* = \emptyset$ and we can assume that $
		\bm{v}^k\rightarrow \bm{v}^*$ as $k\rightarrow\infty$. 
		
		\noindent Now suppose that $
		I^* \ne \emptyset.$ We establish a contradiction this time in two cases.
		\begin{mycases}
			When $t_*=1.$ Rewrite $(\ref{4a})$ as 
			%
			\begin{align} \label{eq:right_zero}
				\textbf{J}^T_ h(\bm{x}^k)\bm{v}^k & + \bm{x}^k-\bm{x}^0 + \sum_{i\in I(\bm{x}^*)}(1-t_k)u_i^k\textbf{J}_ {g_i}(\bm{x}^k) =(1-t_k)\left[\bm{x}^k-\bm{x}^0 -\textbf{J}^T_ f(\bm{x}^k)\textbf{w}^k - \sum_{i\not\in I(\bm{x}^*)}u_i^k\textbf{J}_ {g_i}(\bm{x}^k)\right].
			\end{align}
			
			\noindent Observe that for all $i \notin I(\bm{x}^*)$, since $ I^*\subset I(\bm{x}^*)$ then $i \notin I^*$ and $\lim\limits_{k\rightarrow
				\infty} u_i^k<\infty$. So that as $k\rightarrow \infty,$ since Assumption (B) holds,
			\begin{align*}
				\lim\limits_{k\rightarrow \infty} ~ (1-t_k) &\left[\bm{x}^k-\bm{x}^0 -\textbf{J}^T_ f(\bm{x}^k)\textbf{w}^k - \sum_{i\not\in I(\bm{x}^*)}u_i^k\textbf{J}_ {g_i}(\bm{x}^k)\right] \\
				&= (1-t_*)\left[\bm{x}^*-\bm{x}^0 -\textbf{J}^T_ f(\bm{x}^*)\text{w}^* - \sum_{i\not\in I(\bm{x}^*)} \left(\lim\limits_{k\rightarrow \infty} u_i^k\right) \textbf{J}_ {g_i}(\bm{x}^*)\right] = 0 
			\end{align*}
			%
			\noindent The above equation $(\ref{eq:right_zero})$ becomes\\
			\begin{equation}\label{cone}
				\bm{x}^* +\textbf{J}^T_h(\bm{x}^*)\bm{v}^* +\sum_{i\in I(\bm{x}^*)}\lim_{k\rightarrow
					\infty}[(1-t_k)u_i^k]\textbf{J}_{g_i}(\bm{x}^*)=\bm{x}^0\in \hat{\Omega}.
			\end{equation}
			Equation~\eqref{cone} implies
			\[
			\bm{x}^0\in \bm{x}^*+\mathcal N(\bm{x}^*).
			\]
			If $I^*\neq\emptyset$, then $\bm{x}^*\in\partial\Omega$, which contradicts
			Assumption (C)(i). Hence $I^*=\emptyset$.
		\end{mycases}
		
		\begin{mycases}
			When $t_*\in [0,1).$ Rewrite $(\ref{4a})$ as
			\begin{align}
				t_k(\bm{x}^k-\bm{x}^0) + \textbf{J}^T_h(\bm{x}^k)\bm{v}^k + (1-t_k) \left [\textbf{J}^T_f(\bm{x}^k)\textbf{w}^k +\sum_{i\not\in I^*}u_i^k\textbf{J}_ {g_i}(\bm{x}^k) \right] +(1-t_k)\sum_{i\in I^*}u_i^k\textbf{J}_ {g_i}(\bm{x}^k)=0.
			\end{align}
			%
			\noindent Using the information we have proved so far that $\textbf{w}^k,~ \bm{v}^k, ~u_i^k ~ (i\not\in I^*)$ are bounded, $t_*\in [0,1),$ and by Assumption (B),
			\begin{align}
				\underbrace{\left[t_k(\bm{x}^k-\bm{x}^0) + \textbf{J}^T_ h(\bm{x}^k)\bm{v}^k \right]}_\text{bounded}  + \underbrace{(1-t_k) \left [\textbf{J}^T_f(\bm{x}^k)\textbf{w}^k +\sum_{i\not\in I^*}u_i^k\textbf{J}_ {g_i}(\bm{x}^k) \right]}_\text{bounded} + \underbrace{\left[(1-t_k)\sum_{i\in I^*}u_i^k\textbf{J}_ {g_i}(\bm{x}^k) \right]}_\text{unbounded} =0.
			\end{align}
			\noindent as $k \to \infty$. This is a contradiction, hence $I^* = \emptyset$.
		\end{mycases}
		
		On the other hand, suppose that $\bm{x}^k\rightarrow \bm{x}^* \in \Omega^0 $. We will show that $\{\bm{u}^k\}$ and $\{\bm{v}^k\}$ have to be bounded. We know that $g(\bm{x}^*)<0$ since $\bm{x}^* \in \Omega^0$. Without loss of generality, let $\{\bm{x}^k\}$ be a subsequence such that $g(\bm{x}^k)\neq0$, then from $(\ref{4c})$, 
		\begin{equation}
			\lim\limits_{k\rightarrow \infty} \bm{u}^k  = \lim\limits_{k\rightarrow \infty} \dfrac{t_k \bm{u}^0 g(\bm{x}^0)}{g(\bm{x}^k)} = \dfrac{t_* \bm{u}^0 g(\bm{x}^0)}{g(\bm{x}^*)} < \infty
		\end{equation}
		That is, $\bm{u}^k$ is bounded. Also, from $(\ref{4a})$, since $\bm{x}^k, t_k, \textbf{w}^k, \bm{u}^k$ all are bounded and Assumption (B) holds, then $\bm{v}^k$ has to be bounded. This completes the proof. 
		Therefore, $\Gamma_{\bm{\omega}^0}$ is bounded. Furthermore, there exists a smooth continuation curve originating from $(\bm{\omega}^0,1)$.
	\end{proof}
	
	
	\vspace{0.3cm}
	%%%%%%%%%%%%%%%%%%%%%%%%%%%%%%%%%%%%%%%%%%%%%%%
	\begin{theorem}\label{conver} Suppose that the conditions (A),
		$(B)$, and (C) hold. Then for almost all $\bm{\omega}^0\in
		\Omega^0\times  W^{++} \times \mathbb{R}^{m}_{++}\times
		\mathbb{R}^s,$ the zero-point set $H_{\bm{\omega}^0}^{-1}(0)$ of the
		homotopy map (\ref{homotopy}) contains a smooth curve
		\del{$\Gamma_{\bm{\omega}^0}\subset \Omega \times
			\mathbb{R}^{p+m}_{+}\times (0,1],$} 
		\add{$\Gamma_{\bm{\omega}^0}\subset \Omega \times
			\mathbb{R}^{p+m}_{+}\times\mathbb{R}^s\times (0,1]$,} which starts from $(\bm{\omega}^0,1).$
		\del{As $t\rightarrow 0,$ the limit set}
		\add{The set of accumulation points of
			$\Gamma_{\bm{\omega}^0}$ with $t\rightarrow0$,}
		\del{$T\times \{0\}\subset
			\Omega \times \mathbb{R}^{p+m}_{+}\times \{0\}$} 
		%
		\add{$T\times \{0\}\subset
			\Omega \times \mathbb{R}^{p+m}_{+}
			\times\mathbb{R}^s\times \{0\}$} of
		$\Gamma_{\bm{\omega}^0}$ is nonempty, and every point in $T$ is a
		solution of (\ref{KKT}).
	\end{theorem}
	\vspace{-0.5cm}
	%%%%%%%%%%%%%%%%%%%%%%%%%%%%%%%%%%%%%%%%%%%%%%%%%
	\begin{proof} By \add{Lemma}\,\ref{zhengze} and Lemma \ref{youjie}, the existence of $\Gamma_{\bm{\omega}^0}$ is obtained. It remains to show
		that the limit set \del{$T\times \{0\}\subset \Omega \times
			\mathbb{R}^{p+m}_{+}\times \{0\}$} 
		\add{$T\times \{0\}\subset
			\Omega \times \mathbb{R}^{p+m}_{+}
			\times\mathbb{R}^s\times \{0\}$ }of $\Gamma_{\bm{\omega}^0}$ is nonempty as
		$t\rightarrow 0.$  By the classification theorem of
		one-dimensional smooth manifolds (Lemma \ref{fenglei}), $\Gamma_{\bm{\omega}^0}$ is
		diffeomorphic to a unit circle or the unit interval $(0,1]$ (See Definition \ref{def_diffeomo}). 
		
		\begin{align*}
			\frac{\partial H_{\bm{\omega}^0}( \bm{\omega} ^0, 1)}{\partial \bm{\omega}} &=
			\begin{matrix}
				\begin{bmatrix}
					I & 0 & 0 & \textbf{J}_ h(\bm{x}^0) \\
					\textbf{J}_ h(\bm{x}^0)^T & 0 & 0 & 0 \\
					\textbf{J}_{\bm{g}}(\bm{x}^0)^T\bm{u}^0 & 0 & \text{diag}(\bm{g}(\bm{x}^0)) & 0\\
					0 & -\frac{3}{2} \text{diag}((\textbf{w}^{0})^{\frac{1}{2}}) & 0 & 0 
				\end{bmatrix}
			\end{matrix}\\
			&= \begin{matrix}
				\begin{bmatrix}
					I_{n} & 0_{n\times p} & 0_{n\times m} & \textbf{J}_ h(\bm{x}^0)_{n\times s} \\
					\textbf{J}_h(\bm{x}^0)^T_{s\times n} & 0_{s\times p} & 0_{s\times m} & 0_{s\times s} \\
					\textbf{J}_g(\bm{x}^0)^T\bm{u}^0_{m\times n} & 0_{m\times p} & \text{diag}(\bm{g}(\bm{x}^0))_{m\times m} & 0_{m\times s}\\
					0_{p\times n} & -\frac{3}{2} \text{diag}((\textbf{w}^{0})^{\frac{1}{2}})_{p\times p} & 0_{p\times m} & 0_{p\times s}
				\end{bmatrix}
			\end{matrix}
		\end{align*}
		
		Since $g(\bm{x}^0)<0, \textbf{w}^0 \in \mathbb{R}_{++}^{p}$, and with Assumption (B),  we have that  $\partial H_{\bm{\omega}^0}( \bm{\omega} ^0, 1)/\partial \bm{\omega}$ is nonsingular. Therefore, the smooth curve
		$\Gamma_{\bm{\omega}^0}$, which starts from $(\bm{\omega}^0, 1)$ is diffeomorphic to $(0, 1]$. Let
		$(\bar{\bm{\omega}}, \bar t)=
		(\bar{\bm{x}}, \bar{\textbf{w}}, \bar{\bm{u}}, \bar{\bm{v}}, \bar t)$ be the limit point of $\Gamma_{\bm{\omega}^0}$, then only three cases are possible:
		\begin{description}
			\item (a)~$(\bar{\bm{\omega}}, \bar t) \in \Omega \times
			\mathbb{R}^{p+m}_{+}\times \mathbb{R}^s\times \{1\};$
			\item (b)~$(\bar{\bm{\omega}}, \bar t) \in \partial (\Omega^0\times
			\mathbb{R}^{p+m}_{++})\times \mathbb{R}^s\times (0,1);$
			\item (c)~$(\bar{\bm{\omega}}, \bar t)\in \Omega \times
			\mathbb{R}^{p+m}_{+}\times \mathbb{R}^s\times \{0\}.$
		\end{description}
		We prove that Case (c) is the only possible case.
		If Case (a) is true, the limit point satisfies $\bar{t} = 1$. However, the equation $H(\bm{\omega}, \bm{\omega}^0, 1) = 0$ has only one solution: $\bm{\omega} = \bm{\omega}^0$. Since $\Gamma_{\bm{\omega}^0}$ starts at $(\bm{\omega}^0, 1)$, it cannot return to this point as $t \to 0$. Thus, this case is impossible.
		If Case (b) is true, that means that $(\bar{\bm{\omega}}, \bar t) \in \partial (\Omega^0\times \mathbb{R}^{p+m}_{++})\times \mathbb{R}^s\times (0,1).$ 
		Thus, at least one of the following holds:
		\begin{description}
			\item (1) $\bar{\bm{u}} \in \partial \mathbb{R}^m_{++}~$ (i.e., $\bar{u}_j = 0$ for some $j$),
			\item (2) $\bar{\bm{x}} \in \partial \Omega^0~$ (i.e., $g_j(\bar{\bm{x}}) = 0$ for some $j$),
			\item (3) $\bar{\textbf{w}} \in \partial \mathbb{R}^p_{++}~$ (i.e., $\bar{\text{w}}_j = 0$ for some $j$).
		\end{description}
		
		\noindent 
		If Case (1) is true, $\bar{\bm{u}} \in \partial \mathbb{R}^m_{++}$. Suppose $\bar{u}_{j_0} = 0$ for some $j_0$. Then, there exists a sequence $\{(\bm{\omega}^k, t_k)\} \subset \Gamma_{\bm{\omega}^0}$ such that $u^k_{j_0} \to 0$. From the homotopy equation (\ref{4c}):
		\[
		u^k_{j_0} g_{j_0}(\bm{x}^k) = t_k u^0_{j_0} g_{j_0}(\bm{x}^0).
		\]
		Taking the limit as $k \to \infty$, we have:
		\begin{align*}
			0 = \lim_{k\rightarrow \infty} \left[u^k_{j_0}g_{j_0}(\bm{x}^k) \right] &= \lim_{k\rightarrow \infty} \left[t_ku^0_{j_0}g_{j_0}(\bm{x}^0) \right] = 
			\bar t u_{j_0}^0g_{j_0}(\bm{x}^0) <0.\\
			0 = \bar{u}_{j_0} g_{j_0}(\bar{\bm{x}}) &= \bar{t} u^0_{j_0} g_{j_0}(\bm{x}^0) < 0.
		\end{align*}
		% \[
		% \bar{u}_{j_0} g_{j_0}(\bar{\bm{x}}) = \bar{t} u^0_{j_0} g_{j_0}(\bm{x}^0).
		% \]
		Since $\bar{u}_{j_0} = 0, ~ g_{j_0}(\bm{x}^0) < 0, ~\bar t \in (0,1)$ and $\bm{u}^0 \in \mathbb{R}_{++}^{m}$, the right-hand side is negative, but the left-hand side is zero. This is a contradiction.
		
		\noindent If Case (2) is true, $\bar{\bm{x}} \in \partial \Omega^0.$ Suppose $g_j(\bar{\bm{x}}) = 0$ for some $j$. Then, there exists a sequence $\{(\bm{\omega}^k, t_k)\} \subset \Gamma_{\bm{\omega}^0}$ such that $g_j(\bm{x}^k) \to 0$. From the homotopy equation (\ref{4c}):
		\[
		u^k_j g_j(\bm{x}^k) = t_k u^0_j g_j(\bm{x}^0).
		\]
		Taking the limit as $k \to \infty$, we have:
		\begin{align*}
			0 = \lim_{k\rightarrow \infty} \left[u^k_{j}g_{j}(\bm{x}^k) \right] &= \lim_{k\rightarrow \infty} \left[t_ku^0_{j}g_{j}(\bm{x}^0) \right] = 
			\bar t u_{j}^0g_{j}(\bm{x}^0) <0.\\
			0 = \bar{u}_j g_j(\bar{\bm{x}}) &= \bar{t} u^0_j g_j(\bm{x}^0) < 0.
		\end{align*}
		% \[
		% \bar{u}_j g_j(\bar{\bm{x}}) = \bar{t} u^0_j g_j(\bm{x}^0).
		% \]
		Since $g_j(\bar{\bm{x}}) = 0, ~ g_{j_0}(\bm{x}^0) < 0, ~\bar t \in (0,1)$ and $\bm{u}^0 \in \mathbb{R}_{++}^{m}$, the right-hand side is negative, but the left-hand side is zero. This is a contradiction.
		
		
		\noindent If Case 3 is true, $\bar{\textbf{w}} \in \partial \mathbb{R}^p_{++}.$ Suppose $\bar{\text{w}}_{j_0} = 0$ for some $j_0$. Then, there exists a sequence $\{(\bm{\omega}^k, t_k)\} \subset \Gamma_{\bm{\omega}^0}$ such that $\textbf{w}^k_{j_0} \to 0$. From the homotopy equation (\ref{4d}):
		\[\displaystyle
		(1 - t_k)\left(1 - \sum_{i=1}^p \text{w}^k_i\right) - t_k \left[ (\text{w}^k_{j_0})^{3/2} - (\text{w}^0_{j_0})^{3/2} \right] = 0.
		\]
		Taking the limit as $k \to \infty$, we have:
		\begin{equation} \label{eq: sum_alternative}
			(1 - \bar{t})\left(1 - \sum_{i=1}^p \bar{\text{w}}_i\right) + \bar{t} (\text{w}^0_{j_0})^{3/2} = 0.
		\end{equation}
		
		
		\noindent Since $\bar{t} \in (0, 1)$ and $\text{w}^0_{j_0} > 0$, this implies that $1 - \sum_{i=1}^p \bar{\text{w}}_i < 0$, or equivalently:
		\begin{equation} \label{eq: sum_big_1}
			\displaystyle \sum_{i=1}^p \bar{\text{w}}_i > 1.
		\end{equation}
		
		\noindent Notice that, inequality (\ref{eq: sum_big_1}) poses a problem in itself as we generally expect $\displaystyle \sum_{i=1}^p \text{w}^k_i \to 1$ and not exceed $1.$ However, let's go further and generate more contradictions. Observe that from equation (\ref{eq: sum_alternative}), 
		\begin{align}
			% (1 - \bar{t})\left(1 - \sum_{i=1}^p \bar{\text{w}}_i\right) &+ \bar{t} (\text{w}^0_{j_0})^{3/2} = 0. \nonumber \\
			(1 - \bar{t})\left(1 - \sum_{i=1}^p \bar{\text{w}}_i\right) & = -\bar{t} (\text{w}^0_{j_0})^{3/2}. \label{eq_sum_substitute}
		\end{align}
		
		\noindent Pick $j \neq j_0$ such that $\text{w}_j^0$ is the smallest among the components of $\text{w}^0$, the homotopy equation (\ref{4d}) in the limit becomes:
		\begin{align}
			(1 - t_k)\left(1 - \sum_{i=1}^p \text{w}^k_i\right) - t_k \left[ (\text{w}^k_{j})^{3/2} - (\text{w}^0_{j})^{3/2} \right] = 0.
		\end{align}
		\noindent Substituting $\displaystyle (1 - \bar{t})\left(1 - \sum_{i=1}^p \bar{\text{w}}_i\right)  = -\bar{t} (\text{w}^0_{j_0})^{3/2}$ from equation (\ref{eq_sum_substitute}), we obtain:
		%
		\begin{align}
			-\bar{t} \left(\text{w}^0_{j_0}\right)^{3/2} &- \bar t \left[ (\bar {\text{w}}_{j})^{3/2} - (\text{w}^0_{j})^{3/2} \right] = 0. \nonumber \\
			(\bar{\text{w}}_j)^{3/2} &= (\text{w}_j^0)^{3/2} - (\text{w}_{j_0}^0)^{3/2}. \label{eq:w_depending}
		\end{align}
		
		\noindent Equation (\ref{eq:w_depending}) implies that $\bar {\text{w}}_j$ 
		depends on the difference between $(\text{w}_j^0)^{3/2}$ and $(\text{w}_{j_0}^0)^{3/2}$. Since $\text{w}_j^0$ is the smallest component, then $(\text{w}_j^0)^{3/2} \leq (\text{w}_{j_0}^0)^{3/2}$. 
		If $(\text{w}_j^0)^{3/2} < (\text{w}_{j_0}^0)^{3/2}$ i.e., $\text{w}_j^0 < \text{w}_{j_0}^0,$ the right-hand side of equation (\ref{eq:w_depending}) becomes negative, which is impossible because the left side $\bar{\text{w}}_j \geq 0$. This is a contradiction. Thus, $(\text{w}_j^0)^{3/2} = (\text{w}_{j_0}^0)^{3/2}$, i.e., $\text{w}_j^0 = \text{w}_{j_0}^0$. This leads to a deeper contradiction as follows:
		%       
		\begin{enumerate}
			\item If $\text{w}_j^0 = \text{w}_{j_0}^0$ for all $j$, this implies all components of $\text{w}^0$ are equal. Then $\bar{\text{w}}_j = 0$ from  equation (\ref{eq:w_depending}). This is impossible since there is at least one $\bar{\text{w}}_j > 0$ for some $j$ because equation (\ref{eq: sum_big_1}) forces at least one $\bar{\text{w}}_j$ to ``large enough'' to make the sum exceed $1$. Thus, $\text{w}_j^0 \neq \text{w}_{j_0}^0$ for some $j$, restoring the contradiction.
			
			\item If $\text{w}_j^0 = \text{w}_{j_0}^0$ for specific $j\neq j_0$, we have $\text{w}_j^0 = \text{w}_{j_0}^0 =: \text{w}_{\min}^0 > 0.$
			Then equation \((\ref{eq:w_depending})\) yields
			\[
			(\bar{\text{w}}_j)^{3/2} = (\text{w}_j^0)^{3/2} - (\text{w}_{j_0}^0)^{3/2} = (\text{w}_{\min}^0)^{3/2} - (\text{w}_{\min}^0)^{3/2} = 0,
			\]
			so that
			\[
			\bar{\text{w}}_j = 0.
			\]
			%
			\noindent However, by assumption we already had \(\bar{\text{w}}_{j_0}=0\). Hence, at least two components of the multiplier vector \(\bar{\text{w}}\) vanish. Recall that equation from \((\ref{eq: sum_big_1})\) we have $\sum_{i=1}^p \bar{\text{w}}_i > 1.$
			On the other hand, the original homotopy construction is set up 
			so that the multiplier components satisfy
			$\sum_{i=1}^p \text{w}_i^0 = 1$,
			and the homotopy is designed so that along the smooth curve the multipliers remain in a neighborhood of \(\text{w}^0\). In the limit, if two (or more) components become zero (i.e. \(\bar{\text{w}}_{j_0} = \bar{\text{w}}_j = 0\)), then the remaining components (being perturbations of their original values) cannot compensate enough to push the sum beyond \(1\). In fact, if several components vanish, the limit sum \(\sum_{i=1}^p \bar{\text{w}}_i\) must be strictly less than \(1\) (or at most equal to \(1\) under a very special and nongeneric scenario), contradicting the inequality $\sum_{i=1}^p \bar{\text{w}}_i > 1.$
			Thus, the assumption that \(\text{w}_j^0 = \text{w}_{j_0}^0\) for some \(j\neq j_0\) leads to a contradiction.
		\end{enumerate} 
		
		\noindent Hence \textbf{Case (b)} is impossible.
		
		Since \textbf{Case (a)} and \textbf{Case (b)} are impossible, the only possible limit is \textbf{Case (c)}: $(\bar{\bm{\omega}}, \bar{t}) \in \Omega \times \mathbb{R}^{p+m}_{+} \times \mathbb{R}^s \times \{0\}$. By continuity, $\bar{\bm{\omega}}$ satisfies the KKT system at $t = 0$. Thus, the limit set $T \times \{0\}$ is nonempty, and every point in $T$ is a solution of the KKT system. 
		\medskip
		By Theorem \ref{conver}, for almost all $\bm{\omega}^0\in \Omega\times
		\mathbb{R}^{p+m}_{++}\times \mathbb{R}^s\times\{0\},$ the homotopy equation
		(\ref{homotopy}) generates a smooth curve $\Gamma_{\bm{\omega}^0}$ which
		we call the homotopy path, and 
		\del{the $\bm{\omega}$ component of
			$(\bm{\omega}(s), t(s))$ in the homotopy path,
			is the solution of (\ref{KKT}) as
			$t(s)\rightarrow 0$.}
		\add{every accumulation point of the
			$\bm{\omega}$ component along the homotopy
			path with $t(s)\rightarrow0$ satisfies
			(\ref{KKT}).}
		
	\end{proof}
	%
	Theorem~\ref{conver} establishes that, for almost all initial points \(\bm{\omega}^0\) satisfying the stated regularity assumptions, the proposed homotopy map generates a smooth continuation path \(\Gamma_{\bm{\omega}^0}\) originating from \((\bm{\omega}^0,1)\). As the homotopy parameter approaches zero, the continuation trajectory converges to a nonempty limit set whose \(\bm{x}\)-components satisfy the KKT conditions of the constrained MOO problem~\eqref{MOO}. The result highlights the importance of constructing an appropriate homotopy map that smoothly interpolates between an initial well-posed system and the original KKT system, thereby ensuring stable continuation trajectories and robust convergence behavior.
	
	The proposed framework can be interpreted as a constraint-based predictor--corrector homotopy continuation method for constrained multiobjective nonlinear optimization. The homotopy formulation directly embeds the KKT conditions, including stationarity, feasibility, complementarity, and normalization conditions, into a unified continuation system $H(\bm{x},\textbf{w},\bm{u},\bm{v},t)=0,$
	where \(t\in[0,1]\) denotes the homotopy parameter. Unlike scalarization-based approaches that repeatedly solve independent optimization problems for different parameter selections, the proposed framework continuously follows branches of KKT points as \(t\to0\). The complementarity structure
	\[
	u_i g_i(\bm{x})-t u_i^0 g_i(\bm{x}^0)=0
	\]
	establishes strong connections with complementarity homotopies and interior-point continuation methods. However, unlike classical interior-point algorithms that typically converge to a single minimizer, the proposed continuation framework is designed to explore complex solution structures, including disconnected and nonconvex branches.
	
	%%%%%%%%%%%%%%%%%%%%%%%%%%%%%%%%%%%%%%%%%%%%%%%%%%%%%%%%%%%%%%%
	\section{Numerical Continuation Algorithms}
	\label{sec:alg}
	
	This section presents two numerical continuation algorithms used in the proposed homotopy framework. 
	Algorithm~\ref{alg:homotopy} gives the basic predictor--corrector continuation method, which is suitable for relatively smooth and connected solution branches. Algorithm~\ref{alg:robust_homotopy} extends this method to more difficult MOO problems involving nonconvex feasible regions, degeneracy, near-singular Jacobians, and disconnected Pareto-critical structures. 
	The main differences between the two algorithms are the tangent-orientation strategy, the correction mechanism, the use of pseudo-arclength continuation, adaptive step-size control, and branch restarting.
	\add{For $p$ objectives, the weight vectors are selected from a finite
		discretization of the positive weight simplex $W^{++}$ defined above,
		\[
		W^{++}
		=
		\left\{
		\mathbf{w}\in\mathbb{R}_{++}^{p}:
		\sum_{i=1}^{p}w_i=1
		\right\}.
		\]
		A simplex-lattice discretization is used to obtain the finite set
		$W=\{\mathbf{w}_1,\ldots,\mathbf{w}_P\}\subset W^{++}$ employed in
		Algorithms~1 and~2.}
	
	\add{A regular distribution of weights in $W^{++}$ does not, in general,
		guarantee a uniform distribution of the corresponding Pareto-critical
		solutions. The mapping from weights to Pareto-critical points may be
		nonlinear and non-injective, so different weights may lead to the same
		or nearby solutions. This effect can be particularly pronounced for
		nonconvex or degenerate MOPs. Therefore, the weight discretization is
		used as a sampling strategy rather than as a guarantee of uniform
		Pareto-front coverage. For disconnected or degenerate problems,
		Algorithm~2 additionally employs multiple branch/restart initializations
		to improve the exploration of distinct Pareto-critical components.}
	
	\add{The use of positive weights in the homotopy formulation should not be
		confused with classical weighted-sum scalarization. Here, the weights
		enter the Pareto-criticality/KKT system and the associated homotopy
		construction; the algorithm does not independently minimize a weighted
		sum of the objectives for each weight. Consequently, the KKT formulation
		is not restricted to supported Pareto points, and Pareto-critical
		solutions on nonconvex portions of the Pareto front are not excluded by
		the formulation. Nevertheless, this does not guarantee that every
		nonconvex portion or every connected component of the Pareto front will
		be reached from a finite set of weights and initializations.}
	
	
	\subsection{Basic Predictor--Corrector Homotopy Continuation}
	
	Although the homotopy framework is theoretically designed to converge to KKT points under the assumptions stated earlier, its numerical implementation requires careful path tracking. 
	The basic continuation algorithm follows a standard predictor--corrector strategy. 
	At each step, a tangent direction is computed from the null space of the homotopy Jacobian, and an Euler-type prediction is used to advance along the solution path. \add{The determinant in the orientation test is not associated with the number of objectives. If
		$J_H\in\mathbb{R}^{N\times(N+1)}$ is the Jacobian of the homotopy
		system and $\boldsymbol{\xi}\in\mathbb{R}^{N+1}$ is its normalized null vector, then
		\[
		\begin{pmatrix}
			J_H\\
			\boldsymbol{\xi}^{T}
		\end{pmatrix}
		\in\mathbb{R}^{(N+1)\times(N+1)}.
		\]
		Thus, the determinant-based orientation test applies for an arbitrary
		number $p$ of objectives; the dimension $N$ changes with the
		corresponding homotopy system.} 
	A Moore--Penrose Newton correction is then applied to project the predicted point back onto the homotopy manifold. 
	This approach is efficient for well-behaved problems where the Pareto-critical branch is smooth, connected, and not strongly affected by folds or near-singular Jacobians.
	
	\begin{algorithm}[htb]
		\caption{Homotopy Continuation Method (Predictor-Corrector)}
		\label{alg:homotopy}
		\begin{algorithmic}[1]
			\STATE Objectives $\bm{f} \in \R^p$, Constraints $\bm{g} \in \R^m,~ \bm{h} \in \R^s$
			\STATE Initial guess $\bm{x_0} \in \R^n$, step size $\alpha_0$, tolerance $\epsilon$, and parameters $u_0, v_0, t_0$
			\STATE \del{Weights $\bb{W} = \{\bb{w_1}, \bb{w_2}, \dots, \bb{w_p}\}$, Pareto optimal points $P_{set} \gets \emptyset$}
			\add{Generate $\bb{W} = \{\bb{w_1}, \bb{w_2}, \dots, \bb{w_P}\} \subset W^{++}$ using the prescribed simplex-lattice discretization}
			\FOR{\textbf{each} $\bb{w} \in \bb{W}$}
			% \STATE Initialize parameters: $u_0, v_0, \mu_0$
			% \STATE $\lambda_0 \gets [w, 1-w]$
			% \STATE $u_0, v_0, \mu_0 \gets 1.0, 1.0, 1.0$
			\STATE Define state vector $\omegaVec \gets [\bm{x_0}, \bb{w}, \bm{u_0}, \bm{v_0}, t_0]^T$
			\STATE $k \gets 1$ 
			\add{\STATE $\alpha \gets \alpha_0$}
			\WHILE{$|t| > \epsilon$ \textbf{and} $k \leq k_{max}$}
			\STATE Evaluate the homotopy and its Jacobian:
			$H \gets H(\omegaVec),\quad J_H \gets D H(\omegaVec).$
			% \STATE Compute Homotopy Jacobian $J_H(\omegaVec)$
			% \STATE Compute Homotopy Jacobian $J_H(\omegaVec) \in \R^{m \times n}$    
			\STATE \textbf{Phase 1: Prediction (Tangent Step)}
			\STATE Compute tangent vector $\xi \in \text{null}(J_H)$ such that $J_H \cdot \bm{\xi} = 0$
			\STATE Normalize: $\xi \gets \xi / \norm{\xi}$
			
			\STATE \textit{Orientation Check:}
			\IF{$\det\!\left( \begin{bmatrix} J_H \\ \bm{\xi}^\top \end{bmatrix} \right) < 0$}
			\STATE $\bm{\xi} \gets -\bm{\xi}$
			\ENDIF       
			\STATE Predict next point: $\omegaVec_{pred} \gets \omegaVec + \alpha \cdot \bm{\xi}$
			\STATE Adapt step size:
			\[
			\alpha =
			\begin{cases}
				\min(\alpha_0, 2\alpha), & \text{if } \|H\| < 0.01 \\
				\add{\max(\alpha_{\min}, \alpha/2)}, & \text{if } \|H\| > 1
			\end{cases}
			\]
			\add{where $\alpha_{\min}$ is a minimum allowable step size. Alternatively, use simply $\alpha/2$ if no minimum step size is imposed.}
			\STATE \textbf{Phase 2: Correction (Newton-Raphson)}
			\STATE Evaluate $H(\omegaVec_{pred})$ and Jacobian $J_{pred} = DH(\omegaVec_{pred})$       
			\STATE Solve minimum norm correction (Moore-Penrose):
			\STATE \quad $\bm{y} \gets (J_{pred} J_{pred}^T)^{-1} H(\omegaVec_{pred})$
			\STATE \quad $\Delta \omegaVec \gets -J_{pred}^T \bm{y}$
			
			\STATE Update state: $\omegaVec \gets \omegaVec_{pred} + \Delta \omegaVec$
			\STATE Extract $t$ from $\omegaVec$
			\STATE $k \gets k + 1$
			\ENDWHILE   
			\IF{$|t| \leq \epsilon$}
			\STATE Extract $ \bm{x}^*$ from $\omegaVec$
			\STATE $P_{set} \gets P_{set} \cup \{ \bm{x}^*, \bm{f}(\bm{x}^*)\}$
			\ENDIF
			\ENDFOR
			\STATE $P_{set}$
		\end{algorithmic}
	\end{algorithm}
	
	\subsection{Enhanced Continuation for Nonconvex and Disconnected Problems}
	
	For more challenging MOO problems, the basic predictor--corrector method may become unstable or may fail to follow disconnected or highly curved solution structures. 
	Therefore, Algorithm~\ref{alg:robust_homotopy} introduces several modifications.
	
	First, the determinant-based orientation check in Algorithm~\ref{alg:homotopy} is replaced by a continuous tangent-orientation rule. 
	Instead of relying on the sign of a determinant, the new method compares the current tangent direction with the previous tangent direction and reverses the orientation whenever their inner product is negative. 
	This reduces abrupt changes in direction and improves path-tracking stability.
	
	Second, the one-step Moore--Penrose correction is replaced by a multi-step Newton corrector applied to an augmented pseudo-arclength system. These continuation enhancements are related to modern predictor--corrector continuation approaches for multiobjective optimization, including the Pareto Tracer framework developed by~\cite{martin2018pareto}.
	The additional condition
	\[
	\xi^T(\bm{\omega}-\bm{\omega}_{\mathrm{pred}})=0
	\]
	keeps the corrected iterate on a hyperplane normal to the tangent direction at the predicted point. 
	This improves robustness near folds, turning points, near-singular Jacobians, and highly nonlinear regions where the homotopy parameter alone may not be a reliable local continuation parameter.
	
	Third, Algorithm~\ref{alg:robust_homotopy} incorporates adaptive step-size control. 
	The step size is increased after fast Newton convergence and decreased when the corrector converges slowly or fails. 
	This allows the method to take larger steps on smooth portions of the branch while using smaller steps near difficult regions.
	
	Finally, branch restarts from multiple feasible initial points are included to improve coverage of disconnected feasible regions and disconnected Pareto-critical sets. 
	Since a local continuation method cannot automatically jump between disconnected branches, the restart strategy allows the algorithm to explore multiple components of the solution structure.
	\begin{algorithm}[htbp]
		\caption{Robust Pseudo-Arclength Homotopy Continuation Method}
		\label{alg:robust_homotopy}
		\begin{algorithmic}[1]
			\STATE Given objectives $\bm f\in\mathbb{R}^p$, inequality constraints $\bm g\in\mathbb{R}^m$, equality constraints $\bm h\in\mathbb{R}^s$.
			\STATE Choose branch restart points $\mathcal X_0=\{\bm x_0^{(1)},\dots,\bm x_0^{(R)}\}$, \add{Generate} weights $\mathcal W=\{\bm w^{(1)},\dots,\bm w^{(N_w)}\}$\add{$\subset
				W^{++}$} , initial step size $\alpha_0$, tolerance $\epsilon$, and maximum iterations $k_{\max}$.
			\STATE Initialize the Pareto-critical solution set $P_{\mathrm{set}}\gets\emptyset$.
			
			\FOR{each restart point $\bm x_0^{(r)}\in\mathcal X_0$}
			\FOR{each weight vector $\bm w\in\mathcal W$}
			\STATE Initialize
			$\bm{\omega}_0=\left[\bm x_0^{(r)},\bm w,\bm u_0,\bm v_0,\mu_0\right]^T,
			\qquad \mu_0=1.$    
			\STATE Set $\bm{\omega}\gets\bm{\omega}_0$, $\alpha\gets\alpha_0$, and compute an initial tangent $\xi_{\mathrm{old}}$.
			\STATE Set $k\gets0$.
			
			\WHILE{$|\mu|>\epsilon$ and $k<k_{\max}$}
			\STATE Evaluate the homotopy system and Jacobian: $H(\bm{\omega}), J_H(\bm{\omega})=DH(\bm{\omega}).$
			\STATE \textbf{Prediction step.}
			\STATE Compute a tangent vector
			$\xi\in\operatorname{null}(J_H(\bm{\omega})),
			\qquad \|\xi\|=1.$
			\STATE Orient the tangent continuously:
			$\text{if } \xi^T\xi_{\mathrm{old}}<0,\quad \xi\gets-\xi.$
			\del{If $\mu>\epsilon$ and $\xi_\mu>0$, set $\xi\gets-\xi$ to encourage progress toward $\mu=0$.}
			\STATE Predict $\bm{\omega}_{\mathrm{pred}}=\bm{\omega}+\alpha\xi.$
			
			\STATE \textbf{Correction step: pseudo-arclength Newton iteration.}
			\STATE Set $\bm{\omega}^{(0)}\gets\bm{\omega}_{\mathrm{pred}}$.
			\FOR{$j=0,1,\dots,j_{\max}$}
			\STATE Form the augmented residual
			$R(\bm{\omega}^{(j)})=
			\begin{bmatrix}
				H(\bm{\omega}^{(j)})\\
				\xi^T(\bm{\omega}^{(j)}-\bm{\omega}_{\mathrm{pred}})
			\end{bmatrix}.$
			\STATE Form the augmented Jacobian
			$J_{\mathrm{aug}}(\bm{\omega}^{(j)})=
			\begin{bmatrix}
				DH(\bm{\omega}^{(j)})\\
				\xi^T
			\end{bmatrix}.$
			\STATE Solve $J_{\mathrm{aug}}(\bm{\omega}^{(j)})\Delta\bm{\omega}
			=
			-R(\bm{\omega}^{(j)}).$
			\STATE Update
			$\bm{\omega}^{(j+1)}=\bm{\omega}^{(j)}+\Delta\bm{\omega}.$     \IF{$\|R(\bm{\omega}^{(j+1)})\|\leq\epsilon_{\mathrm{Newton}}$}
			\STATE break
			\ENDIF
			\ENDFOR
			\STATE Accept the corrected point $\bm{\omega}_{\mathrm{new}}\gets\bm{\omega}^{(j+1)}$ if it is finite, satisfies the residual tolerance, and makes sufficient progress toward $\mu=0$.
			\STATE \add{ If the correction is rejected, reduce $\alpha$ and retry the prediction--correction step without updating $\bm{\omega}$. }      
			\STATE Adapt the step size $\alpha$: increase $\alpha$ after fast Newton convergence and decrease $\alpha$ after slow convergence or failed correction.
			\STATE Set $\xi_{\mathrm{old}}\gets\xi$, $\bm{\omega}\gets\bm{\omega}_{\mathrm{new}}$, and $k\gets k+1$.
			\ENDWHILE
			\IF{$|\mu|\leq\epsilon$}
			\STATE Extract $\bm{x}^*$ from $\bm{\omega}$.
			\STATE Add the Pareto-critical point:
			$P_{\mathrm{set}}\gets P_{\mathrm{set}}\cup\{(\bm{x}^*,\bm f(\bm{x}^*))\}.$
			\ENDIF
			\ENDFOR
			\ENDFOR
			\STATE Return $P_{\mathrm{set}}$.
		\end{algorithmic}
	\end{algorithm}
	Overall, Algorithm~\ref{alg:homotopy} is intended for standard continuation on smooth connected branches, while Algorithm~\ref{alg:robust_homotopy} is designed for more difficult many-objective, nonconvex, degenerate, and disconnected constrained optimization problems. 
	Both algorithms compute Pareto-critical points rather than guaranteed globally Pareto-optimal points, since the homotopy system traces KKT stationary solutions.
	
	%%%%%%%%%%%%%%%%%%%%%%%%%%%%%%%%%%%%%%%%%%%%%%%%%%%%%%%%%%%%%%%%%
	\section{Computational Experiments}
	\label{sec:numerical}
	%%%%%%%%%%%%%%%%%%%%%%%%%%%%%%%%%%%%%%%%%%%%%%%%%%%%%%%%%%%%%%%%%
	This section presents several computational experiments designed to evaluate the proposed homotopy continuation framework on constrained multiobjective optimization problems with varying levels of complexity. The experiments include benchmark problems, nonconvex and disconnected many-objective problems, and an IEEE 5-bus optimal power flow application.
	
	
	\del{All numerical experiments were conducted in Google Colab
		using a 2021 MacBook Pro equipped with an Apple M1 Pro
		processor and 16 GB of unified memory.}
	\add{All numerical experiments were performed locally
		on a Windows laptop equipped with a 13th-generation
		Intel Core i7-1355U processor (1.70 GHz) and 16 GB
		of RAM.}
	The implementations were developed in Python~3.12 using standard scientific computing libraries, including NumPy and SciPy. All computations were performed in double precision, and consistent termination tolerances were applied across all methods to ensure fair numerical comparison. The implementation code and computational scripts are publicly available  \del{on Dropbox to support reproducibility and further research:
		\url{https://tinyurl.com/2z52n284}.}
	\add{\url{https://xxxxxxxxxxxxxxxxxxxxx}.}
	
	%%%%%%%%%%%%%%%%%%%%%%%%%%%%%%%%%%%%%%%%%%%%%%%%%%%%%%%%%%
	To evaluate the proposed homotopy continuation framework, we compare it with several widely used MOO approaches, including the Weighted Sum Method (WSM), the $\epsilon$-Constraint Method (ECM), the Global Criterion Method (GCM), the Lexicographic Method, and the evolutionary algorithm NSGA-II. These methods remain popular in engineering and scientific applications due to their simplicity, interpretability, and established computational frameworks.
	
	All six methods are tested on the benchmark problem using consistent stopping tolerances, feasibility checks, and evaluation metrics to ensure fair comparison. For the deterministic scalarization methods (Weighted Sum, $\epsilon-$\add{Constraint}, Global Criterion, and Lexicographic optimization), constrained single-objective subproblems are solved to obtain Pareto-critical KKT solutions. The proposed homotopy framework is implemented as a continuation scheme applied directly to the KKT system. Although the lexicographic method does not generate a full Pareto set, it is included here as a classical deterministic baseline for comparison of representative Pareto-critical solutions and tradeoff behavior. NSGA-II is included as a stochastic \textit{a posteriori} baseline that generates a finite nondominated approximation of the Pareto front. Across all examples, we report Pareto-critical solutions, computational cost, and objective trade-off behavior.
	%%%%%%%%%%%%%%%%%%%%%%%%%%%%%%%%%%%%%%%%%%%%%%%%%%%%
	%-------------------------------------------------------------
	%\medskip
	%\add{\noindent\textbf{Pareto-front Quality Measures}}
	%-------------------------------------------------------------
	
	\add{To quantify the quality of the computed nondominated sets, let
		$A=\{\bm{x}_1,\ldots,\bm{x}_{N_A}\}$ denote the normalized objective
		vectors obtained by a given method and let
		$R=\{\bm{r}_1,\ldots,\bm{r}_{N_R}\}$ denote the normalized empirical
		reference front. For a point $\bm{y}$ and a finite set $S$, define
		\[d(\bm{y},S)=\min_{\bm{s}\in S}\|\bm{y}-\bm{s}\|_2.\]
		The following Pareto-front quality measures, including the inverted generational distance (IGD), the hausdorff distance, the spacing coefficient or variation (Spacing CV), the maximum gap, and the coverage, are used:}
	%
	\begin{itemize}
		\item\add{The number of nondominated solutions produced by a method after filtering the computed feasible solutions for Pareto dominance is denoted by $N_{\mathrm{ND}}$.} 
		\item The maximum equality-constraint violation
		\[
		\max_{\bm{x}\in A_{\mathrm{dec}}}
		\max_{1\leq i\leq m}
		\max\{g_i(\bm{x}),0\},
		\]
		
		\item The maximum inequality-constraint violation 
		\[\max_{\bm{x}}[g(\bm{x})]_+:=\max\{g(\bm{x}),0\},
		\]
		denote the maximum equality-constraint residual and maximum
		inequality-constraint violation, respectively, over the reported
		nondominated solutions.
		
		\item
		\add{ \textbf{The inverted generational distance (IGD)}\,\cite{van1999multiobjective,15892}: }
		\[\operatorname{IGD}(A,R) =
		\frac{1}{N_R}
		\sum_{\bm{r}\in R}d(\bm{r},A)\]
		measures how well the computed set represents the reference front.
		Smaller values indicate better overall approximation and coverage.
		\smallskip 
		\item
		\add{The \textbf{Hausdorff distance\,\cite{6151115}:}}
		\[d_H(A,R)
		=
		\max\left\{
		\max_{\bm{x}\in A}d(\bm{x},R),
		\;
		\max_{\bm{r}\in R}d(\bm{r},A)
		\right\}\]
		measures the largest discrepancy between the computed and reference
		sets. Smaller values are preferable.
		\smallskip 
		\item
		\add{The \textbf{spacing coefficient of variation (Spacing CV)}:}
		\[
		\operatorname{CV}_{\mathrm{sp}}
		=
		\frac{\operatorname{std}(d_1,\ldots,d_{N_A})}
		{\operatorname{mean}(d_1,\ldots,d_{N_A})},
		\]
		where $d_i$ is the distance from $\bm{x}_i$ to its nearest neighboring
		point in $A$. Smaller values indicate a more uniform distribution of
		computed points.
		\smallskip 
		\item
		\add{The \textbf{maximum gap} is the largest separation between consecutive
			computed points after ordering them along the front. Smaller values
			indicate fewer large gaps in the approximation.}
		\smallskip 
		\item
		\add{The \textbf{coverage} measures the fraction of the objective range of
			the reference front spanned by the computed nondominated set. Values
			closer to $1$ indicate better coverage of the reference-front range.}
		
		\item Hypervolume\,\cite{zitzler2007hypervolume,797969}
	\end{itemize}
	%%%%%%%%%%%%%%%%%%%%%%%%%%%%%%%%%%%%%%%%%%%%%%%%%%%%
	
	\begin{table}[htbp]
		\centering
		\caption{Comparison of MOO methods.}
		\label{tab:comparison_methods}
		\renewcommand{\arraystretch}{1.2}
		
		\begin{tabular}{p{2.2cm} p{2cm} p{4.8cm} p{4.4cm}}
			\hline
			\textbf{Method}
			& \textbf{Category}
			& \textbf{Mechanism}
			& \textbf{Solution-set characteristics} \\
			\hline
			
			Homotopy
			& Continuation method
			& Continuation applied to the KKT system
			& Computes KKT-consistent Pareto-critical solutions through homotopy
			continuation \\
			\hline
			
			Weighted Sum
			& Deterministic scalarization
			& Linear weighted aggregation with $\mathbf{w}$
			& Varying positive weights recovers supported Pareto-optimal solutions;
			non-supported portions of a nonconvex Pareto front may be missed \\
			\hline
			
			$\epsilon$-Constraint
			& Deterministic scalarization
			& Optimize one objective while bounding the remaining objectives
			& Can recover both supported and non-supported Pareto-optimal solutions
			by varying the objective bounds \\
			\hline
			
			Global Criterion
			& Deterministic scalarization
			& Distance minimization to an ideal point $\bm{f}^{\ast}$
			& Produces compromise solutions according to the selected metric and
			scaling \\
			\hline
			
			Lexicographic
			& Deterministic scalarization
			& Sequential priority-based optimization
			& Produces solutions determined by the prescribed objective-priority
			ordering \\
			\hline
			
			NSGA-II
			& Evolutionary algorithm
			& Population-based nondominated sorting and evolution
			& Produces a population-based approximation of the
			nondominated solution set \\
			\hline
		\end{tabular}
	\end{table}
	
	\del{Table~\ref{tab:comparison_methods} summarizes representative MOO approaches and highlights their differences in formulation, solution mechanisms, and typical Pareto coverage. Classical scalarization methods transform a multiobjective problem into one or more single-objective optimization problems whose KKT systems generate Pareto-critical solutions. The Pareto regions recovered depend strongly on the scalarization strategy, with weighted-sum methods typically restricted to convex regions of the Pareto front. In contrast, homotopy continuation methods operate as deterministic path-following procedures applied directly to the KKT system, while NSGA-II produces stochastic population-based approximations of the Pareto front.}
	\add{Table~\ref{tab:comparison_methods} summarizes representative MOO
		approaches and distinguishes their solution mechanisms and the types of
		Pareto solutions they are designed to recover. In the weighted-sum
		method, varying positive objective weights generates supported
		Pareto-optimal solutions; consequently, non-supported portions of a
		nonconvex Pareto front may not be recovered. The
		$\epsilon$-constraint method instead optimizes one objective while
		placing bounds on the others and, by varying these bounds, can recover
		both supported and non-supported Pareto-optimal solutions. The global
		criterion and lexicographic methods generate compromise or
		priority-dependent solutions according to their respective
		scalarizations. The proposed homotopy method applies continuation to the KKT system to compute Pareto-critical solutions, whereas NSGA-II uses a population-based stochastic search to approximate the nondominated solution set. } 
	
	Additional implementation details for the comparison algorithms are
	provided in Appendices~\ref{appendix:A}--\ref{appendix:B}.
	
	\add{The following four examples are designed to evaluate complementary
		aspects of the proposed homotopy continuation method, progressing from
		a constrained benchmark problem to degenerate and disconnected
		many-objective problems, a practical engineering application, and
		controlled problems with analytically known Pareto fronts. In
		particular,}
	
	\begin{itemize}
		\item \textbf{Example~1:} Constrained bi-objective benchmark problem.
		
		This example compares the proposed method with existing deterministic
		approaches and NSGA-II. It examines the quality and distribution of
		the computed Pareto-front approximation, constraint satisfaction,
		and sensitivity to the equality-constraint tolerance and initial
		point.
		
		\item \textbf{Example~2:} Degenerate and disconnected four-objective
		problem.
		
		This example tests the method on a more challenging many-objective
		problem having disconnected feasible components and a degenerate
		Pareto structure. It is used to assess the continuation and
		branch-restarting strategy, Pareto-front coverage, and sensitivity
		to the construction of the numerical reference set.
		
		\item \textbf{Example~3:} Three-objective IEEE 5-bus optimal power
		flow problem.
		
		This example evaluates the method on a constrained engineering
		application with multiple competing objectives. It provides a test
		of the proposed approach on a problem for which an analytical Pareto
		front is not available and permits comparison with modern
		many-objective evolutionary algorithms.
		
		\item \textbf{Example~4:} Three-objective problems with analytically
		known Pareto fronts.
		
		\add{Two closely related problems with convex and concave Pareto-front
			geometries are considered. Because their Pareto sets, Pareto fronts,
			and interior Pareto-critical points are known analytically, this
			example permits direct accuracy assessment. The concave case also
			provides a controlled investigation of unsupported Pareto-critical
			points and the effect of initialization on homotopy branch selection.}
	\end{itemize}
	
	\add{The numerical comparisons are tailored to the purpose of each example.
		Where an analytical Pareto front is available, the computed solutions
		are compared directly with the exact front; otherwise, an empirical
		reference set is used. Feasibility and KKT or homotopy residuals are
		reported together with appropriate measures of convergence, coverage,
		and distribution. Modern evolutionary multiobjective algorithms,
		including NSGA-III and MOEA/D/D, are included where appropriate. The
		purpose of these experiments is not to establish universal superiority
		of one method, but to assess the accuracy, feasibility, coverage,
		robustness, and limitations of the proposed continuation approach under
		different objective, constraint, and Pareto-front structures.}
	
	\begin{example}[\textbf{Constrained Bi-Objective Benchmark Problem}] 
		\label{EX2}
		
		This example serves as a benchmark validation problem for the proposed framework. The experiments investigate feasible-region structure, comparison with classical scalarization and evolutionary methods, constraint feasibility of reported literature solutions, computational efficiency, and robustness with respect to initialization.
		Consider the MOO problem:
		\[
		\begin{aligned}
			\min_{x \in \mathbb{R}^5} \quad 
			\begin{bmatrix}
				f_1(\bm{x}) \\
				f_2(\bm{x})
			\end{bmatrix}  
			&=
			\begin{bmatrix}
				x_1^2 + x_2^2 + x_3^2 + x_4^2 + x_5^2 \\
				3x_1 + 2x_2 - \tfrac{1}{3}x_3 + 0.01(x_4 - x_5)^3
			\end{bmatrix}, \\\\
			\text{s.t.} \quad
			\begin{bmatrix}
				h_1(\bm{x}) \\
				h_2(\bm{x})
			\end{bmatrix}  &= 
			\begin{bmatrix}
				4x_1 - 2x_2 + 0.8x_3 + 0.6x_4 + 0.5x_5^2 \\
				x_1 + 2x_2 - x_3 - 0.5x_4 + x_5 - 2
			\end{bmatrix}
			= 0.\\
			g(\bm{x}) &= x_1^2 + x_2^2 + x_3^2 + x_4^2 - 10 \leq 0.
		\end{aligned}
		\]
	\end{example}
	
	%-----------------------------------------------------------
	\medskip
	\noindent \textbf{Feasible Region Approximation}
	%-----------------------------------------------------------
	
	
	
	\del{Due to the sparsity of feasible points under direct random
		sampling, a projection-based strategy is employed to
		approximate the feasible decision space.}
	\add{Since the equality constraints restrict the feasible
		set to a lower-dimensional subset of the decision space,
		uniform random sampling is unlikely to produce points
		satisfying these constraints exactly. \add{This difficulty is also highlighted
			in the benchmark study of equality-constrained multiobjective
			optimization by Cuate et al.~\cite{cuate2020benchmark}.} Therefore, a
		projection-based strategy is employed to approximate
		the feasible decision space.}
	Uniform samples $\bm{x}_0 \in [-5,5]^5$ are generated and projected onto the feasible set by solving
	\begin{equation}
		\label{eq:projection-problem}
		\begin{aligned}
			\min_{\bm{z} \in \mathbb{R}^5} \quad & \|\bm{z} - \bm{x}_0\|_2^2 \\
			\text{s.t.} \quad & g(\bm{z}) \le 0, \\
			& h_1(\bm{z}) = 0, \\
			& h_2(\bm{z}) = 0 .
		\end{aligned}
	\end{equation}
	
	\begin{figure}[h]
		\hspace*{-0.5cm}
		\centering
		\includegraphics[scale=0.35]{Code/Example2/sampled_proj_feas_homotopy_new.png}
		\caption{Example~\ref{EX2}: \add{Approximation of the feasible
				objective-space region.} Blue points \add{in the left panel} represent $200{,}000$ uniformly sampled diagnostic points from the hypercube $[-5,5]^5$; \add{while } golden points denote $20{,}000$ feasible candidate solutions obtained by projecting onto the constraint set. The \del{left panel displays all diagnostic samples (faint blue) together with the projected feasible candidates (gold), while the} right panel shows only the projected feasible candidates \del{and the corresponding homotopy-based Pareto front.} \add{Pareto-critical
				solutions obtained by the different MOO methods are compared separately
				in Fig.~\ref{fig:pareto_plots_ex1}.}}
		\label{Fig:points}
	\end{figure}
	The resulting optimization problem identifies the nearest feasible point to $\bm{x}_0$ in the Euclidean sense and is solved using the Sequential Least Squares Programming (SLSQP) algorithm. The projected feasible points are then used to approximate the feasible objective-space structure and provide a reference for subsequent Pareto-critical solution comparisons. 
	Repeating this procedure over a large number of samples yields an empirical approximation of the feasible \del{set}\add{objective-space region}, which serves as a reference for subsequent \del{Pareto-structure identification}\add{comparison of Pareto-critical
		solutions}. In Fig.~\ref{Fig:points}, \del{blue points represent uniformly sampled decision variables, while golden points denote feasible solutions obtained through projection onto the constraint set. The left panel illustrates the relationship between the sampled and projected points, whereas the right panel isolates the feasible solutions together with the homotopy-generated Pareto-critical structures.} \add{the left panel shows the diagnostic samples together with the projected feasible candidates, whereas the right panel isolates the projected feasible candidates. The
		Pareto-critical solutions produced by the different MOO methods are
		shown and compared separately in Fig.~\ref{fig:pareto_plots_ex1}.} 
	
	%--------------------------------------------------------
	\medskip \noindent \textbf{Pareto-critical Solution Comparison}
	%--------------------------------------------------------
	
	Fig.\,\ref{fig:pareto_plots_ex1} \add{illustrates} the Pareto structure and numerical results for Example\,\ref{EX2}. This compares the Pareto solutions obtained by the proposed homotopy,  weighted sum, global criterion, NSGA-II, lexicographic, and $\epsilon-$constraint methods. \del{The weighted sum and homotopy approaches recover identical Pareto-critical points, confirming that the homotopy method accurately solves the KKT system of the weighted-sum formulation.} 
	\add{For a prescribed objective-weight vector $\bm{w}$, the endpoint of the proposed homotopy satisfies the corresponding weighted KKT
		system of the original constrained MOO problem. The direct weighted-sum
		method and the proposed homotopy therefore target solutions satisfying
		the same weighted stationarity conditions when the same weights are
		used. Their computed points need not, however, coincide numerically:
		in a nonconvex problem the weighted KKT system may admit multiple
		solutions, and the direct optimization and continuation procedures may
		converge to different KKT points. Accordingly, Fig.~\ref{fig:pareto_plots_ex1}
		should be interpreted as a comparison of the Pareto-critical structures
		recovered by the different numerical procedures, rather than as evidence
		that the weighted-sum and homotopy methods produce identical points.
	} 
	The global criterion and $\epsilon-$constraint methods generate \del{continuous Pareto curves} \add{approximations of Pareto curves,} while NSGA-II produces a dense but stochastic approximation of the Pareto front. In contrast, the lexicographic method yields only extreme solutions due to its strict priority structure. These results highlight the deterministic and structure-preserving nature of the homotopy approach relative to both classical scalarization techniques and evolutionary methods.
	
	
	\begin{figure}[htb]
		\hspace*{-0.5cm}
		\centering
		\includegraphics[scale=0.35]{Code20261002/Example1/figure2_six_methods_updated.png}
		\caption{Example~\ref{EX2}: Feasible region obtained by projecting $20{,}000$ randomly generated candidate points, sampled uniformly from the hypercube $[-5,5]^5$, onto the constraint set. \del{Each blue point represents a feasible solution in the $(f_{1}, f_{2})$ objective space.} Pareto fronts generated by the weighted sum, proposed homotopy, global criterion, genetic algorithm NSGA-II, Lexicographic, and $\epsilon-$constraint method are denoted in red color. For NSGA-II, the equality constraints are imposed using the feasibility tolerance $\tau_{\mathrm{eq}}=10^{-2}$; the sensitivity of the computed front to this tolerance is examined separately in Fig.~\ref{fig:nsga-tolerance}.} 
		\label{fig:pareto_plots_ex1}
	\end{figure}
	
	
	\begin{table}[h!]
		\centering
		\caption{Normalized Pareto-front quality measures for Example~1.
			The empirical reference front is formed from the nondominated union
			of the highly feasible solutions generated by the four deterministic
			front-generating methods.}
		\label{tab:front-quality}
		\begin{tabular}{lrrrrrr}
			\hline
			Method & ND pts. & IGD & Hausdorff & Spacing CV & Max gap & Coverage \\
			\hline
			Homotopy
			& 198 & 0.00653 & 0.07684 & 2.390 & 0.1553 & 0.999 \\
			
			Weighted Sum
			& 198 & 0.00754 & 0.09059 & 2.629 & 0.1857 & 0.999 \\
			
			Global Criterion
			& 199 & 0.01627 & 0.08024 & 0.378 & 0.0243 & 0.908--0.928 \\
			
			$\epsilon$-Constraint
			& 150 & 0.00227 & 0.02074 & 0.342 & 0.0418 & 1.000 \\
			\hline
		\end{tabular}
	\end{table}
	
	
	
	\add{The quantitative comparison confirms that the differences visible in
		Fig.~\ref{fig:pareto_plots_ex1} concern primarily the distribution of the
		computed points along the Pareto front. The proposed homotopy method
		and the weighted-sum method provide nearly complete objective-range
		coverage and exhibit similar distribution characteristics. For this
		example, the homotopy method gives slightly smaller IGD, Hausdorff
		distance, spacing variation, and maximum gap than the weighted-sum
		method.}
	
	\add{The global-criterion and $\epsilon$-constraint methods generate more
		uniformly distributed objective vectors. In particular, their spacing
		coefficients of variation are substantially smaller than those of the
		homotopy and weighted-sum methods. The global-criterion method,
		however, does not recover the extreme portions of the objective range
		as completely, whereas the $\epsilon$-constraint method provides full
		range coverage in this experiment. Of the 199 feasible solutions
		generated by the $\epsilon$-constraint method, 150 remain after
		nondominance filtering.}
	
	\add{These observations concern the distribution of a finite set of
		computed points and should not be interpreted as a general ranking of
		the optimization methods. In particular, uniform spacing of the
		continuation or scalarization parameter does not in general imply
		uniform spacing in objective space.}
	
	%=========================================================
	\medskip\noindent \add{\textbf{Sensitivity of NSGA-II to the equality-feasibility tolerance}}
	%=========================================================
	
	\add{For NSGA-II, the equality constraints cannot in general be expected to
		be satisfied exactly by the population generated during the
		evolutionary search. We therefore implement the equality constraints
		through the feasibility conditions
		\(|h_j(\bm{x})| \leq \tau_{\mathrm{eq}}, j=1,2, \)
		where $\tau_{\mathrm{eq}}>0$ is an equality-feasibility tolerance.
		This parameter should not be confused with the parameter used in the
		$\epsilon$-constraint method.}
	
	\add{To examine the sensitivity of NSGA-II to this choice, we repeated the
		calculation using
		$\tau_{\mathrm{eq}} = 10^{-1}$, $10^{-2}$, $10^{-3}$, $10^{-4}$,
		while keeping the remaining NSGA-II parameters fixed. In each run, a
		population size of 800 and 800 generations were used, corresponding to
		640,000 candidate evaluations.}
	
	\begin{figure}[htbp]
		\centering
		\includegraphics[width=0.8\textwidth]{Code20261002/Example1/nsga2_equality_tolerance_sensitivity.png}
		\caption{\add{Example~\ref{EX2}: Sensitivity of the NSGA-II solutions to the
				equality-feasibility tolerance $\tau_{\mathrm{eq}}$ for Example~1.
				The four calculations use
				$\tau_{\mathrm{eq}}=10^{-1},10^{-2},10^{-3},10^{-4}$,
				respectively, while all other NSGA-II settings are held fixed.
				Decreasing the tolerance improves equality-constraint satisfaction,
				but under the fixed computational budget the recovered solutions
				become increasingly concentrated in a smaller portion of the
				Pareto front.}}
		\label{fig:nsga-tolerance}
	\end{figure}
	
	\add{The results in Fig.~\ref{fig:nsga-tolerance} show that the apparent
		Pareto-front coverage obtained by NSGA-II depends strongly on the
		equality-feasibility tolerance. For
		$\tau_{\mathrm{eq}}=10^{-1}$, NSGA-II produces a broad distribution of
		solutions across almost the entire objective range. However, the
		maximum equality residual is approximately $10^{-1}$, so these points
		should be interpreted as solutions satisfying a relatively loose
		feasibility band rather than as high-accuracy solutions of the original
		equality-constrained problem.}
	
	\add{When the tolerance is reduced to
		$\tau_{\mathrm{eq}}=10^{-2}$, the maximum equality residual decreases
		accordingly, but the recovered objective-space coverage is substantially
		reduced. Further tightening to $10^{-3}$ and $10^{-4}$ causes the
		computed solutions to become increasingly concentrated in a small
		portion of the front under the same population size and number of
		generations. Thus, improved equality feasibility is accompanied by a
		loss of front coverage for the fixed computational budget used in this
		experiment.}
	
	\add{This behavior is also reflected in the normalized front-quality
		measures. Using the highly feasible nondominated solutions from the
		deterministic methods as an empirical reference front, the results are
		summarized in Table~\ref{tab:nsga-tolerance}.}
	
	\begin{table}[htbp]
		\centering
		\caption{Sensitivity of NSGA-II to the equality-feasibility tolerance
			for Example~1. Front-quality measures are computed in normalized
			objective space relative to the empirical reference front constructed
			from the highly feasible deterministic solutions.}
		\label{tab:nsga-tolerance}
		\begin{tabular}{c r r r r r}
			\hline
			$\tau_{\mathrm{eq}}$
			& Points
			& Max.\ $|h_j|$
			& IGD
			& Hausdorff
			& $f_1/f_2$ coverage \\
			\hline
			$10^{-1}$
			& 800
			& $9.999\times10^{-2}$
			& 0.00536
			& 0.0253
			& 0.988 / 0.978 \\
			
			$10^{-2}$
			& 800
			& $1.000\times10^{-2}$
			& 0.15242
			& 0.5983
			& 0.622 / 0.489 \\
			
			$10^{-3}$
			& 629
			& $1.000\times10^{-3}$
			& 0.42010
			& 1.1005
			& 0.0329 / 0.0236 \\
			
			$10^{-4}$
			& 736
			& $9.999\times10^{-5}$
			& 0.43785
			& 1.1195
			& 0.0018 / 0.0011 \\
			\hline
		\end{tabular}
	\end{table}
	
	\add{The first row of Table~\ref{tab:nsga-tolerance} illustrates why front
		coverage alone can be misleading for equality-constrained problems.
		With $\tau_{\mathrm{eq}}=10^{-1}$, NSGA-II achieves an IGD of
		approximately $5.36\times10^{-3}$ and covers approximately
		$98.8\%$ and $97.8\%$ of the reference ranges of $f_1$ and $f_2$,
		respectively. These apparently favorable front-quality measures are
		obtained, however, while permitting equality residuals of order
		$10^{-1}$.}
	
	\add{In contrast, the deterministic methods considered above satisfy the
		equality constraints to substantially higher numerical accuracy. For
		example, the maximum equality residuals in the present calculations are
		approximately $1.25\times10^{-12}$ for the proposed homotopy method,
		$9.57\times10^{-11}$ for the weighted-sum method,
		$9.39\times10^{-11}$ for the global-criterion method, and
		$2.18\times10^{-11}$ for the $\epsilon$-constraint method.}
	
	\add{Consequently, the NSGA-II results should not be compared with the
		deterministic methods solely on the basis of the number or visual
		distribution of the computed objective vectors. For the present
		equality-constrained example, there is a clear trade-off between
		equality-feasibility accuracy and Pareto-front coverage for NSGA-II
		under a fixed computational budget. The experiment is intended to
		illustrate this sensitivity rather than to provide a general ranking
		of NSGA-II relative to the deterministic methods.
	}
	
	
	\del{Fig.~\ref{fig:exa1fig1} summarizes the numerical results for Example~\ref{EX2} obtained using several MOO methods. The proposed homotopy method, the weighted-sum method, and the global criterion method converge to nearly identical Pareto-critical solutions with objective values clustered around \(f(\bm{x}^*) \approx (2.7,-0.41)\). The $\epsilon$-constraint method yields a nearby trade-off solution, while NSGA-II identifies a different compromise point with smaller objective values. However, as shown in Table~\ref{tab:constraints}, the NSGA-II solution exhibits significantly larger equality-constraint residuals, whereas the proposed homotopy method satisfies both equality constraints to near machine precision while maintaining strict feasibility of the inequality constraint. These results demonstrate that the proposed continuation framework effectively computes KKT-consistent Pareto-critical solutions for the constrained MOO problem.}
	
	% \begin{figure}[h]
		% \centering
		% \includegraphics[scale=0.65]{Code/Example2/exa1fig1}
		% \caption{Example~\ref{EX2}: Numerical results obtained by using five MOO methods. For the deterministic methods, the same initial guess $\bm{x}^0=(1,2,0,1,1)$ and weights $\textbf{w}=(0.4,0.6)$ are used when applicable.}
		% \label{fig:exa1fig1}
		% \end{figure}
	% %
	
	% \begin{table}\centering
		% \caption{Example~\ref{EX2}: Constraint evaluation at $\bm{x}^*$ obtained by using five MOO methods.}
		% \label{tab:constraints}
		% \vspace{.5em}
		% \begin{tabular}{lccc}
			% \toprule
			% Method & $h_1(\bm{x}^*)$ & $h_2(\bm{x}^*)$ & $g(\bm{x}^*)$ \\
			% \midrule
			% Proposed & & & \\
			% Homotopy & $-7.36\times10^{-6}$ & $5.00\times10^{-5}$ & $-9.52$ \\
			% Weighted Sum & $-1.18\times10^{-3}$ & $-2.30\times10^{-3}$ & $-9.52$ \\
			% Global Criterion & $-1.18\times10^{-3}$ & $-2.30\times10^{-3}$ & $-9.52$ \\
			% $\epsilon$-Constraint & $-8.73\times10^{-4}$ & $-1.55\times10^{-3}$ & $-9.52$ \\
			% NSGA-II & $-6.78\times10^{-2}$ & $-9.89\times10^{-2}$ & $-9.48$ \\
			% \botrule
			% \end{tabular}
		% \end{table}
	
	
	\begin{table}[h]
		\centering
		\caption{\del{Comparison of CPU time and function evaluations across methods for Example~\ref{EX2}}.\add{Computational results for the methods used to generate the Pareto-front approximations shown in Fig.~\ref{fig:pareto_plots_ex1}.}}
		\begin{tabular}{l c c c c }
			\toprule[1.5pt]
			\\[-1.5ex]
			\textbf{Method} & \textbf{CPU Time (s)} & \textbf{$f_1$} & \textbf{$f_2$} & \textbf{$(g, h_1, h_2)$} \\
			\\[-1.5ex]
			\midrule
			\del{Proposed Homotopy (Parallelized)} & \del{1.67} &  \del{$6{,}559$} & \del{$6{,}559$} & \del{$6{,}559$} \\
			\\[-1.5ex]
			Proposed Homotopy  & 2.10 & $6{,}559$ & $6{,}559$ & $6{,}559$ \\
			\\[-1.5ex]
			Weighted Sum Method  & 0.73 & $4{,}820$ & $4{,}820$ & $5{,}699$ \\
			\\[-1.5ex]
			$\epsilon$-Constraint Method & 1.26 & $5{,}911$ & $7{,}030$ & $7{,}052$ \\
			\\[-1.5ex]
			
			Global Criterion Method & 2.51 & $6{,}386$ & $6{,}408$ & $7{,}579$ \\
			\\[-1.5ex]
			
			Genetic Algorithm NSGA2 & 100.38 & $640{,}000$ & $640{,}000$ & $640{,}000$ \\
			\\[-1.5ex]
			
			\toprule[1.5pt]
		\end{tabular}
		\label{tab:cpu_feval_comparisonEX2}
	\end{table}
	
	\begin{table}[htbp]
		\centering
		\caption{Comparison of the six methods for Example~\ref{EX2}.
			Front-quality measures are computed in normalized objective space
			relative to the empirical reference front constructed from the highly
			feasible deterministic solutions. For NSGA-II,
			$\tau_{\rm eq}=10^{-2}$.}
		\label{tab:example1_all_methods}
		\small
		\begin{tabular}{lrrrrrrr}
			\toprule
			Method
			& $N_{\rm ND}$
			& $\max |h_j|$
			& $\max(g,0)$
			& IGD
			& Hausdorff
			& Spacing CV
			& $f_1/f_2$ coverage \\
			\midrule
			Proposed Homotopy
			& 198 & $1.25{\times}10^{-12}$ & $3.55{\times}10^{-15}$
			& 0.00653 & 0.07684 & 2.390 & 0.999/0.999 \\
			
			Weighted Sum
			& 198 & $9.57{\times}10^{-11}$ & $9.53{\times}10^{-11}$
			& 0.00754 & 0.09059 & 2.629 & 0.999/1.000 \\
			
			Global Criterion
			& 199 & $9.39{\times}10^{-11}$ & $9.30{\times}10^{-11}$
			& 0.01627 & 0.08024 & 0.378 & 0.928/0.908 \\
			
			$\epsilon$-Constraint
			& 150 & $2.18{\times}10^{-11}$ & $6.49{\times}10^{-11}$
			& 0.00227 & 0.02074 & 0.342 & 1.000/1.000 \\
			
			Lexicographical
			& 2 & $1.39{\times}10^{-16}$ & $0$
			& 0.24502 & 0.71372 & -- & 1.000/1.000 \\
			
			NSGA-II ($\tau_{\rm eq}=10^{-2}$)
			& 800 & $1.00{\times}10^{-2}$ & $0$
			& 0.15242 & 0.59832 & 9.419 & 0.622/0.489 \\
			\bottomrule
		\end{tabular}
	\end{table}
	
	%---------------------------------------------------------
	\medskip\noindent \textbf{Computational Performance and Efficiency}
	%---------------------------------------------------------
	
	Table\,\ref{tab:cpu_feval_comparisonEX2} compares the computational cost of the considered methods for Example\,\ref{EX2} in terms of CPU time and the total number of function evaluations. The proposed homotopy method computes Pareto-critical points with moderate computational cost, requiring only a few thousand evaluations of the objective and constraint functions. \del{When parallelization is employed, the CPU time is further reduced while maintaining the same number of function evaluations.}
	In contrast, the weighted-sum method is computationally inexpensive in terms of CPU time but requires a comparable number of objective evaluations, reflecting its reliance on repeated scalar optimization solves. The $\epsilon$-constraint and global criterion methods incur higher computational costs due to the additional constraint handling and penalty evaluations required during the optimization process. NSGA-II is significantly more expensive than the deterministic methods, both in CPU time and function evaluations, as it relies on population-based sampling and repeated fitness evaluations to approximate the Pareto front.
	Overall, the results demonstrate that the proposed homotopy approach provides a favorable balance between computational efficiency and solution quality, achieving KKT-consistent Pareto-critical points at a substantially lower cost than evolutionary MOO methods.
	
	\add{Table~\ref{tab:example1_all_methods} provides a quantitative comparison
		of the six methods using a common empirical reference front and common
		feasibility criteria. The proposed homotopy and weighted-sum methods
		produce very similar approximations of the Pareto front, with essentially
		complete objective-range coverage. The proposed homotopy method yields
		slightly smaller IGD and Hausdorff distance than the weighted-sum method,
		with values of $0.00653$ and $0.07684$, compared with $0.00754$ and
		$0.09059$, respectively. It also gives smaller equality- and
		inequality-constraint residuals. These differences are modest, however,
		and the results should be interpreted as showing comparable performance
		of the two methods for this example rather than numerical superiority of
		one method over the other.}
	
	\add{The $\epsilon$-constraint method gives the best front-quality measures in
		this experiment, attaining the smallest IGD ($0.00227$), Hausdorff
		distance ($0.02074$), and spacing coefficient of variation ($0.342$),
		together with complete objective-range coverage. The global-criterion
		method also produces a relatively uniform distribution of solutions, as
		indicated by its spacing coefficient of variation, but its objective-range
		coverage is lower than that of the homotopy, weighted-sum, and
		$\epsilon$-constraint methods.}
	
	\add{The lexicographical method returns only the two priority-order solutions.
		Consequently, its near-complete objective-range coverage reflects the fact
		that these solutions are located near the two objective extremes and
		should not be interpreted as coverage of the interior of the Pareto
		front. This is also reflected in its substantially larger IGD and
		Hausdorff distance. For this reason, a spacing coefficient of variation is
		not reported for the lexicographical method.}
	
	\add{For NSGA-II with $\tau_{\rm eq}=10^{-2}$, the equality residual reaches
		$10^{-2}$ and the resulting front has substantially larger IGD and
		Hausdorff distance and lower objective-range coverage than the highly
		feasible deterministic fronts. As discussed in the preceding sensitivity
		study, these results depend strongly on the equality-feasibility
		tolerance: tightening the tolerance improves equality feasibility but,
		for the fixed population size and generation budget used here,
		substantially reduces front coverage.}
	
	\add{Overall, the results do not indicate that the proposed homotopy method
		dominates the other methods on every numerical metric. Rather, for
		Example~\ref{EX2}, it produces a highly feasible, nearly complete
		approximation of the Pareto-critical front with front-quality measures
		comparable to those of established scalarization methods. Its principal
		distinction is that these solutions are generated through a continuation
		framework for tracing KKT-consistent Pareto-critical solutions, rather
		than through a sequence of independent scalar optimization problems.}
	
	
	%--------------------------------------------------------
	\medskip\noindent \textbf{Constraint Feasibility Validation}
	%--------------------------------------------------------
	
	Table \ref{tab:comparison_literature_example1} 
	assess the reliability of reported solutions in the literature - \cite{shang2011} and  
	\cite{zhao2012quasinormal}, together with the solution obtained by our proposed homotopy method: 
	
	
	\begin{itemize}
		\item 
		The solution reported in~\cite{shang2011},
		$(0.3077,\,0.5374,\,-0.2703,\,-0.1336,\,0.2804)$,
		yields objective values
		$(f_{1}(\bm{x}), f_{2}(\bm{x})) \approx (0.5530, 2.0873)$.
		The inequality constraint is strictly satisfied
		($g(\bm{x}) \approx -9.5256$),
		and the second equality constraint satisfies
		$h_{2}(\bm{x}) \approx 0$,
		while the first equality constraint exhibits a moderate violation
		($h_{1}(\bm{x}) \approx -0.1011$).
		Thus, the point is close to feasibility but does not satisfy all constraints simultaneously.
		
		\item 
		The solution reported in~\cite{zhao2012quasinormal},
		$(-1.3074,\,-2.8605,\,-1.0470,\,0.4103,\,0.4475)$,
		produces
		$(f_{1}(\bm{x}), f_{2}(\bm{x})) \approx (11.3566, -9.2942)$.
		Although $h_{1}(\bm{x}) \approx 1.08\times 10^{-4}$,
		the second equality constraint has a large violation
		($h_{2}(\bm{x}) \approx -7.7391$),
		and the inequality constraint is violated
		($g(\bm{x}) \approx 1.1563 > 0$).
		Hence, this point is infeasible.
		
		\item 
		The proposed homotopy method computes the solution
		$(0.3176,\,0.5067,\,-0.2797,\,-0.1403,\,0.3192)$
		with objective values
		$(f_{1}(\bm{x}), f_{2}(\bm{x})) \approx (0.5574, 2.0584)$.
		Both equality constraints are satisfied to near machine precision
		($h_{1}(\bm{x}) \approx -1.39\times10^{-17}$ and $h_{2}(\bm{x}) \approx -2.22\times10^{-16}$),
		while the inequality constraint remains strictly feasible
		($g(\bm{x}) \approx -9.5445$).
		\del{Among the compared methods, the proposed homotopy
			framework produces the only fully feasible
			Pareto-critical solution.}
		\add{Among the three candidate solutions reported
			in this comparison, only the solution obtained
			by the proposed homotopy method satisfies all
			constraints to the reported numerical accuracy.}
	\end{itemize}
	\begin{table}[htbp]
		\centering
		\caption{Comparison of candidate solutions for Example~\ref{EX2} and objective values and constraint residuals.
			The proposed homotopy method uses
			$\bm{x}^{0}=(1,2,0,1,1)$ and $\bm{w}=(0.96,0.04)$.
			Here,
			$\|H_{\bm{\omega}^*}\|
			:=\|H(\bm{\omega}^*,\bm{\omega}^{0},t_*)\|$.}
		\label{tab:comparison_literature_example1}
		
		\begin{tabular}{lrrrrrr}
			\toprule
			Method
			& $x_1^*$ & $x_2^*$ & $x_3^*$ & $x_4^*$ & $x_5^*$ & \\
			\midrule
			\cite{shang2011}
			& 0.3077 & 0.5374 & -0.2703 & -0.1336 & 0.2804 & \\
			\cite{zhao2012quasinormal}
			& -1.3074 & -2.8605 & -1.0470 & 0.4103 & 0.4475 & \\
			Proposed Homotopy
			& 0.3176 & 0.5067 & -0.2797 & -0.1403 & 0.3192 & \\
			
			\midrule
			Method
			& $f_1(\bm{x}^*)$
			& $f_2(\bm{x}^*)$
			& $h_1(\bm{x}^*)$
			& $h_2(\bm{x}^*)$
			& $g(\bm{x}^*)$
			& $\|H_{\bm{\omega}^*}\|$ \\
			\midrule
			\cite{shang2011}
			& 0.5530
			& 2.0873
			& -0.1011
			& $\add{0.0000}$
			& -9.5256
			& $6.6\times10^{-16}$ \\
			\cite{zhao2012quasinormal}
			& 11.3566
			& -9.2942
			& $1.08\times10^{-4}$
			& -7.7391
			& 1.1563
			& $1.0\times10^{-12}$ \\
			Proposed Homotopy
			& 0.557392
			& 2.0584
			& $-1.39\times10^{-17}$
			& $-2.22\times10^{-16}$
			& -9.5445
			& $7.1\times10^{-12}$ \\
			\bottomrule
		\end{tabular}
	\end{table}
	
	
	\begin{table}[h]
		\centering
		\caption{\del{Numerical results for Example \ref{EX2} with proposed homotopy algorithm \ref{alg:homotopy} and $\textbf{w} = (0.4, 0.6)$, add feasible Boolean, where $\|H_{\textbf{w}^*}\| = \|H (\textbf{w}^*, \textbf{w}^0, t_*)\|.$} \add{Sensitivity of the proposed homotopy method to the initial
				point for Example~\ref{EX2} with $\bm{w}=(0.4,0.6)$.
				Both feasible and infeasible initial points are considered.
				Here,
				$\|H_{\textbf{w}^*}\|
				:=\|H(\textbf{w}^*,\textbf{w}^0,t_*)\|$.}}
		\begin{tabular}{c c c c c}
			\toprule[1.5pt] \\
			$\bm{x}_0$ & $\bm{x}^0$ \add{Feasible?} & $\bm{x}^*$ & $f(\bm{x}^*)$ & $\|H_{\textbf{w}^*}\|$ \\
			\\
			\midrule
			$(1.0, 2.0, 0.0, 1.0, 1.0)$ & No & $ (-0.1390, -0.0518, -0.5309, -0.4189, 1.5023)$ & $(2.7363, -0.4147)$ & $2.8e^{-5}$ \\ 
			\\
			$(0.5, 0.5, 0.5, 0.5, 0.5)$ & No & $(-0.1375, -0.0488,  -0.5281,  -0.4152, 1.4993)$ & $( 2.7207, -0.4042)$ & $5.4e^{-5}$ \\ 
			\\
			$(-2.0, 0.0, 0.0, 0.0, 4.0)$ & Yes & $ (-0.1370, -0.0507, -0.5311, -0.4180,  1.4983)$ & $(2.7231, -0.4058)$ & $3.7e^{-5}$ \\ 
			\\
			$(0.4, 0.8, 0.0, 0.0, 0.0)$ & Yes & $ (-0.1389, -0.0513, -0.5306, -0.4179, 1.5020)$ & $(2.7341, -0.4132)$ & $5.0e^{-5}$ \\ 
			\\
			\toprule[1.5pt]
		\end{tabular}
		\label{tab:homotopyEX2}
	\end{table}
	
	
	%---------------------------------------------------------
	\medskip\noindent \textbf{Sensitivity to Initialization}
	%---------------------------------------------------------
	
	Table\,\ref{tab:homotopyEX2} reports the numerical results for
	Example\,\ref{EX2} obtained using the proposed homotopy algorithm with
	$\textbf{w}=(0.4,0.6)$, starting from both feasible and infeasible
	initial guesses. The feasibility of each initial point $\bm{x}^0$ is
	explicitly indicated. Regardless of whether $\bm{x}^0$ is feasible, the
	homotopy method consistently converges to a Pareto-critical point with
	objective values clustered around $f(\bm{x}^*)\approx(2.7,-0.41)$.
	\add{The computed solutions are close but not identical, indicating some
		dependence on the initialization.}
	Moreover, the small values of the KKT residual
	$\lVert H_{\bm{x}^*}\rVert$ confirm that the computed solutions satisfy
	the optimality conditions to high accuracy. These results
	\del{demonstrate the robustness of the proposed homotopy approach with
		respect to the choice of initial guess and}
	\add{indicate limited sensitivity to the initial guesses considered here,
		while demonstrating}
	its ability to drive infeasible starting points toward KKT-consistent
	Pareto-critical points.
	
	%%%%%%%%%%%%%%%%%%%%%%%%%%%%%%%%%%%%%%%%%%%%%%%%%%%%%%%%%%%%%%%
	%-------------------------------------------------------------------------
	
	This example highlights one of the key advantages of continuation-based methods: the ability to systematically trace disconnected and nonconvex Pareto structures through branch-following mechanisms embedded in the homotopy formulation.
	
	\begin{example}[\textbf{Degenerate and Disconnected Four-Objective Problem}]
		
		Consider the decision vector \(\bm{x}=(x_1,x_2,x_3,x_4)^T\). We study the constrained four-objective optimization problem
		\begin{equation}
			\label{eq:degenerate_disconnected_mop_eq}
			\begin{aligned}
				\min_{\bm{x}}\quad
				&F(\bm{x})=
				\bigl(f_1(\bm{x}),f_2(\bm{x}),f_3(\bm{x}),f_4(\bm{x})\bigr)\\
				\text{s.t.}\quad
				&x_4=0,\\
				&x_3^2-x_2^2=0,\\
				&(x_1^2-1)^2+x_2^2-0.08\le 0,\\
				&-2\le x_1\le 2,\qquad -1\le x_2,x_3,x_4\le 1,
			\end{aligned}
		\end{equation}
		where
		\begin{equation}
			\label{eq:degenerate_disconnected_objectives_eq}
			\begin{aligned}
				f_1(\bm{x})&=(x_1+1)^2+x_2^2+0.05\sin(8x_1),\\
				f_2(\bm{x})&=(x_1-1)^2+x_2^2+0.05\sin(8x_1),\\
				f_3(\bm{x})&=x_2^2+0.1(1-\cos(6x_1)),\\
				f_4(\bm{x})&=(x_1^2-1)^2+x_2^2.
			\end{aligned}
		\end{equation}
	\end{example}
	
	This example combines several challenging features for continuation-based many-objective optimization. The problem is many-objective because it contains four competing objectives. It is nonconvex due to the oscillatory sinusoidal terms and the quartic nonlinear constraint defining the feasible set. The inequality constraint \((x_1^2-1)^2+x_2^2\le0.08\) creates two disconnected feasible components centered near \(x_1=-1\) and \(x_1=1\). The equality constraint \(x_3^2-x_2^2=0\) \del{introduces a branched nonlinear manifold, while the constraint \(x_4=0\) fixes a redundant variable} \add{defines two intersecting branches, $x_3=x_2$ and $x_3=-x_2$, with a singular intersection at $x_2=x_3=0$, while the constraint $x_4=0$
		fixes one decision variable}. The problem is also degenerate because the objective functions do not explicitly depend on \(x_3\) or \(x_4\). Consequently, multiple feasible points may map to the same objective vector, and the Pareto-critical set in decision space may possess a higher-dimensional or branched structure relative to its image in objective space.
	
	\add{At the intersection $x_2=x_3=0$, the
		gradient of the nonlinear equality constraint
		vanishes, so the linear independence constraint
		qualification (Assumption A2) fails. Consequently, this example
		also tests the numerical continuation method
		outside the regularity assumptions used in the
		global convergence theorem.}
	
	%--------------------------------------------------------------
	\medskip \noindent \textbf{Continuation Setup and Branch Initialization}
	%--------------------------------------------------------------
	
	For this example, the decision dimension is \(n=4\), the number of objectives is \(p=4\), the number of equality constraints is \(s=2\), and the number of inequality constraints is \(m=9\), including the box constraints. The homotopy state vector is initialized as $\bm{\omega}^0=[\bm{x}^0,\bb{w}^0,\bm{u}^0,\bm{v}^0]^T,$
	with initial homotopy parameter \(t_0=1\), equality multipliers \(\bm{v}^0=\mathbf{0}\in\mathbb{R}^2\), and inequality multipliers \(\bm{u}^0=[0.05\,\mathbf{1}]\in\mathbb{R}^9\).
	
	The objective-weight vectors are selected from a positive discretization of the four-objective simplex using discretization level of 10 and lower bound \(0.03\). Thus each weight vector satisfies
	\[
	\textbf{w}^0=(\text{w}_1,\text{w}_2,\text{w}_3,\text{w}_4),\qquad
	\sum_{i=1}^4 \text{w}_i=1,\qquad
	\text{w}_i\ge0.03.
	\]
	
	To explore disconnected branches, the continuation is restarted from multiple feasible initial points near both disconnected feasible components. The restart points are centered near \(x_1=\pm1\), with
	\[
	x_2\in\{0.02,0.04,0.06,0.08,0.10,0.12,0.15,-0.04,-0.08,-0.12,-0.15\}.
	\]
	The equality manifold \(x_3^2-x_2^2=0\) is initialized using both branches \(x_3=x_2\) and \(x_3=-x_2\), with \(x_4=0\). Consequently, the restart points take the form
	$\bm{x}=(\pm1,x_2,\pm x_2,0)^T,$
	subject to the strict interior feasibility condition
	$(x_1^2-1)^2+x_2^2<0.08.$
	The continuation paths are advanced using a pseudo-arclength predictor--corrector strategy. The initial continuation step size is \(\alpha_0=0.06\), with bounds \(\alpha_{\min}=10^{-5}\) and \(\alpha_{\max}=0.10\). The stopping tolerance for the homotopy parameter is \(\epsilon_t=10^{-4}\), and the maximum number of continuation steps per path is \(k_{\max}=900\).
	
	At each continuation step, the tangent direction is computed from the null space of the homotopy Jacobian and oriented consistently using the previous tangent vector. Newton-type corrector iterations are then applied to the augmented pseudo-arclength system with maximum Newton iterations \(15\), Newton tolerance \(10^{-9}\), and a damped line search. A trial step is accepted only if the corrected point is finite, the homotopy residual is sufficiently small, and the continuation path progresses toward \(t=0\). The step size is increased following rapid Newton convergence and reduced following slow convergence or failed corrections.
	
	
	%--------------------------------------------------------------
	\medskip \noindent \textbf{Pareto-Critical Structure and Continuation Results}
	%--------------------------------------------------------------
	
	The following figures illustrate the geometry of the computed Pareto-critical solution set, the disconnected feasible components, the equality-constraint structure, and the objective-space behavior of the proposed homotopy continuation framework.
	%%%%%%%%%%%%%%%%%%%%%%%%%%%%%%%%%%%%%%%%%%%%%%%%%%%%%%%%%%%%%%%%%%%%%%%%%%%
	\begin{figure}[H]
		\centering
		\includegraphics[width=0.45\textwidth]{Code/DegeneratedExample/manyobj_3d_projection_colored_by_f4.png}
		\includegraphics[width=0.45\textwidth]{Code/DegeneratedExample/manyobj_nonconvex_cross_sections_x1.png}
		\caption{
			Left: Three-dimensional projection of the computed Pareto-critical set in the $(f_1,f_2,f_3)$ objective space, with $f_4$ encoded by color. Right: Cross-sections of the four objective functions along the reduced manifold $x_2=x_3=x_4=0$.
		}
		\label{fig:manyobj_3d_projection_colored}
	\end{figure}
	
	The left panel of Fig.~\ref{fig:manyobj_3d_projection_colored} illustrates the geometric structure of the computed Pareto-critical set in objective space. The disconnected vertical clusters correspond to continuation trajectories associated with the two disconnected feasible components centered near \(x_1\approx -1\) and \(x_1\approx 1\). The narrow branch-like structures indicate degeneracy in the objective mapping, where multiple feasible decision vectors collapse onto low-dimensional Pareto-critical structures in objective space. The color variation associated with \(f_4\) further reveals how the fourth objective changes along the continuation trajectories. 
	
	The right panel displays one-dimensional cross-sections of the objective functions along the reduced manifold \(x_2=x_3=x_4=0\). The oscillatory sinusoidal perturbations introduce nonconvexity and multiple competing local trade-off regions, producing a highly nonlinear many-objective optimization landscape. These competing objective structures contribute to the disconnected and branched Pareto-critical behavior observed in the continuation results.
	
	
	Fig.~\ref{fig:manyobj_pairwise_objective} illustrates the pairwise trade-off structure among the four competing objectives. The disconnected vertical branches visible in several projections reflect the presence of disconnected feasible components and multiple continuation branches generated by the homotopy framework. The color distributions indicate that the omitted objectives vary nonuniformly along the Pareto-critical trajectories, demonstrating the coupled many-objective nature of the problem. 
	
	The projections involving \(f_1\) and \(f_2\) exhibit strong separation between disconnected branches associated with the two feasible regions centered near \(x_1\approx -1\) and \(x_1\approx 1\). In contrast, the \((f_3,f_4)\) projection reveals a nearly one-dimensional structure, indicating degeneracy in the objective mapping and strong correlation between these objectives along the computed Pareto-critical set. Overall, the pairwise projections confirm that the proposed continuation framework successfully captures disconnected, nonlinear, and degenerate Pareto-critical structures in objective space.
	%%%%%%%%%%%%%%%%%%%%%%%%%%%%%%%%%%%%%%%%%%%%%%%%%%%%%%%%%%%%%
	
	Fig.~\ref{fig:manyobj_decision_space_3d} illustrates the geometric structure of the computed Pareto-critical set in decision space and its relationship to the disconnected feasible components. The left panel shows that the continuation trajectories collapse near a low-dimensional branched manifold in the \((x_1,x_2,x_3)\) subspace, reflecting the degeneracy introduced by the equality constraint \(x_3^2-x_2^2=0\). The two visible branches correspond to the manifolds \(x_3=x_2\) and \(x_3=-x_2\), while the concentration of solutions near \(x_2=x_3=0\) indicates partial collapse of the Pareto-critical structure in decision space.
	
	The right panel displays the disconnected feasible components induced by the nonlinear inequality constraint \((x_1^2-1)^2+x_2^2\le0.08\). The continuation framework successfully traces Pareto-critical branches associated with both disconnected regions near \(x_1\approx -1\) and \(x_1\approx 1\). These results demonstrate the capability of the proposed continuation strategy to explore disconnected feasible regions and branched Pareto-critical structures within a unified homotopy framework.
	%%%%%%%%%%%%%%%%%%%%%%%%%%%%%%%%%%%%%%%%%%%%%%%%%%%%%%%%%%%%%%%%%%%%%%%%%%%
	\begin{figure}[H]
		\centering
		\includegraphics[width=0.95\textwidth]{Code/DegeneratedExample/manyobj_pairwise_objective_projections_colored_by_omitted_objective.png}
		\caption{
			Pairwise objective-space projections of the computed Pareto-critical set. Each subplot displays a two-objective projection, while the omitted objective is represented by color.
		}
		\label{fig:manyobj_pairwise_objective}
	\end{figure}
	%%%%%%%%%%%%%%%%%%%%%%%%%%%%%%%%%%%%%%%%%%%%%%%%%%%%%%%%%%%%%%%%%%%%%%%%%%%
	\begin{figure}[H]
		\centering
		\includegraphics[width=0.45\textwidth]{Code/DegeneratedExample/manyobj_decision_space_x1_x2_x3.png}
		\includegraphics[width=0.45\textwidth]{Code/DegeneratedExample/manyobj_3d_projection_components.png}
		\caption{Left: Three-dimensional decision-space projection $(x_1,x_2,x_3)$ of the computed Pareto-critical set. Right: Disconnected continuation branches associated with the two disconnected feasible components near \(x_1\approx -1\) and \(x_1\approx 1\).}
		\label{fig:manyobj_decision_space_3d}
	\end{figure}
	%%%%%%%%%%%%%%%%%%%%%%%%%%%%%%%%%%%%%%%%%%%%%%%%%%%%%%%%%%%%%%%%%%%%%%%%%%%
	\begin{figure}[H]
		\centering
		\includegraphics[width=0.4\textwidth]{Code/DegeneratedExample/manyobj_decision_space_x1_x4_equality.png}
		\includegraphics[width=0.45\textwidth]{Code/DegeneratedExample/manyobj_constraint_feasibility_diagnostics.png}
		\caption{
			Left: Equality-constraint structure for \(h_1=x_4=0\). Right: Constraint feasibility diagnostics for the computed Pareto-critical solutions.
		}
		\label{fig:manyobj_feasibility}
	\end{figure}
	Fig.~\ref{fig:manyobj_feasibility} verifies the feasibility of the computed Pareto-critical solutions. The left panel shows that all computed points satisfy the equality constraint \(x_4=0\) to near machine precision, while remaining associated with the two disconnected feasible components near \(x_1\approx -1\) and \(x_1\approx 1\). This confirms that the continuation trajectories remain on the prescribed equality manifold.
	
	The right panel reports the equality residuals and inequality violations along the computed continuation solutions. Both equality residuals remain close to machine precision, and no positive inequality violation is observed. These diagnostics confirm that the computed Pareto-critical points are not merely approximate objective-space samples, but feasible KKT-consistent solutions of the constrained many-objective problem.
	%%%%%%%%%%%%%%%%%%%%%%%%%%%%%%%%%%%%%%%%%%%%%%%%%%%%%%%%%%%%%%%%%%%%%%%%%%%
	%-------------------------------------------------------------
	\medskip \noindent \textbf{Quality Indicators and Feasibility Diagnostics}
	%--------------------------------------------------------------
	
	Because the degenerate disconnected many-objective problem is specifically designed to test approximation quality and coverage of complicated Pareto-critical structures, additional set-based quality indicators are included for this example. 
	\add{We use the hypervolume indicator
		\cite{zitzler1999multiobjective}, generational
		distance (GD) \cite{vanveldhuizen1999multiobjective},
		inverted generational distance (IGD)
		\cite{ishibuchi2015modified}, and averaged
		Hausdorff distance \cite{schutze2012using}
		to assess the quality of the computed
		approximation set relative to a sampled
		reference set.}
	
	%-------------------------------------------------------------
	\begin{table}[h]
		\centering
		\caption{Quality indicators, feasibility diagnostics, and interpretation of the computed Pareto-critical approximation set.}
		\label{tab:quality_indicators_manyobj}
		\renewcommand{\arraystretch}{1.15}
		\begin{tabular}{p{5.2cm} p{2.2cm} p{5.5cm}}
			\toprule
			\textbf{Indicator} & \textbf{Value} & \textbf{Interpretation} \\
			\midrule
			Number of feasible successful solutions 
			& 129 
			& Moderate-size approximation of the disconnected Pareto-critical structure. \\\hline
			Normalized hypervolume 
			& 0.336215924
			& Measures dominated objective-space coverage of the computed approximation set. \\\hline
			Reference-set normalized hypervolume 
			& 0.5666067
			& Hypervolume of the densely sampled reference approximation set. \\\hline
			Hypervolume ratio 
			& 0.593385013
			& Decent but incomplete Pareto coverage for a disconnected nonconvex problem. \\\hline
			Averaged Hausdorff distance 
			& 0.246278236 
			& Moderate discrepancy between the computed approximation and the sampled reference structure. \\\hline
			Generational distance (GD) 
			& 0.019358087 
			& Small GD indicates strong local approximation accuracy. \\\hline
			Inverted generational distance (IGD) 
			& 0.246278236?????????????????
			& Larger IGD suggests incomplete global coverage of disconnected branches. \\\hline
			Maximum equality residual 
			& $7.27\times 10^{-20}$ 
			& Equality constraints are satisfied near machine precision. \\\hline
			Maximum inequality violation 
			& $0$ 
			& Feasibility maintained throughout the continuation process. \\\hline
			Median homotopy residual $\|H\|$ 
			& $7.43\times 10^{-14}$ 
			& Tiny residuals indicate numerically stable continuation and Newton correction. \\
			\bottomrule
		\end{tabular}
	\end{table}
	
	Table~\ref{tab:quality_indicators_manyobj} summarizes quality indicators and feasibility diagnostics for the computed Pareto-critical approximation set. The normalized hypervolume and distance-based indicators (GD, IGD, and averaged Hausdorff distance) suggest that the proposed continuation framework captures a representative approximation of the disconnected Pareto-critical structure while maintaining reasonable diversity across the computed branches. The relatively small generational distance indicates that the computed solutions closely approximate the sampled reference structure locally, while the larger inverted generational distance reflects the difficulty of achieving complete coverage of the disconnected Pareto-critical set through local continuation alone. Furthermore, since the reference set is generated through dense feasible sampling rather than from an exact analytic Pareto-front formula, the hypervolume ratio should be interpreted as a relative coverage indicator with respect to the sampled reference approximation.
	
	
	\add{These indicators characterize the quality
		of the computed approximation relative to
		the empirical reference set. Since no
		alternative optimization methods are evaluated
		on this particular benchmark, the reported
		values should not be interpreted as evidence
		of comparative superiority.}
	
	
	The feasibility diagnostics further demonstrate the numerical stability of the continuation framework. Equality residuals and inequality violations remain near machine precision throughout the continuation process, while the small homotopy residuals confirm accurate satisfaction of the KKT system along the computed trajectories. Because this example contains degeneracy and disconnected Pareto-critical structures, the reported hypervolume and distance metrics should be interpreted as quality measures of the computed approximation set rather than as guarantees of global Pareto optimality. \del{Furthermore, since the reference set is generated through dense feasible sampling rather than from an exact analytic Pareto-front formula, the hypervolume ratio should be interpreted as a relative coverage indicator with respect to the sampled reference approximation.}
	
	%===========================================================
	\begin{example}[\textbf{Degenerate and Disconnected Four-Objective Problem}]
		
		Consider the decision vector \(\bm{x}=(x_1,x_2,x_3,x_4)^T\). We study the constrained four-objective optimization problem
		\begin{equation}
			\label{eq:degenerate_disconnected_mop_eq}
			\begin{aligned}
				\min_{\bm{x}}\quad
				&F(\bm{x})=
				\bigl(f_1(\bm{x}),f_2(\bm{x}),f_3(\bm{x}),f_4(\bm{x})\bigr)\\
				\text{s.t.}\quad
				&x_4=0,\\
				&x_3^2-x_2^2=0,\\
				&(x_1^2-1)^2+x_2^2-0.08\le 0,\\
				&-2\le x_1\le 2,\qquad -1\le x_2,x_3,x_4\le 1,
			\end{aligned}
		\end{equation}
		where
		\begin{equation}
			\label{eq:degenerate_disconnected_objectives_eq}
			\begin{aligned}
				f_1(\bm{x})&=(x_1+1)^2+x_2^2+0.05\sin(8x_1),\\
				f_2(\bm{x})&=(x_1-1)^2+x_2^2+0.05\sin(8x_1),\\
				f_3(\bm{x})&=x_2^2+0.1(1-\cos(6x_1)),\\
				f_4(\bm{x})&=(x_1^2-1)^2+x_2^2.
			\end{aligned}
		\end{equation}
	\end{example}
	
	This example combines several challenging features for continuation-based many-objective optimization. The problem is many-objective because it contains four competing objectives. It is nonconvex due to the oscillatory sinusoidal terms and the quartic nonlinear constraint defining the feasible set. The inequality constraint \((x_1^2-1)^2+x_2^2\le0.08\) creates two disconnected feasible components centered near \(x_1=-1\) and \(x_1=1\). The equality constraint \(x_3^2-x_2^2=0\) \del{introduces a branched nonlinear manifold, while the constraint \(x_4=0\) fixes a redundant variable} \add{defines two intersecting branches, $x_3=x_2$ and $x_3=-x_2$, with a singular intersection at $x_2=x_3=0$, while the constraint $x_4=0$
		fixes one decision variable}. The problem is also degenerate because the objective functions do not explicitly depend on \(x_3\) or \(x_4\). Consequently, multiple feasible points may map to the same objective vector, and the Pareto-critical set in decision space may possess a higher-dimensional or branched structure relative to its image in objective space.
	
	\add{At the intersection $x_2=x_3=0$, the
		gradient of the nonlinear equality constraint
		vanishes, so the linear independence constraint
		qualification (Assumption A2) fails. Consequently, this example
		also tests the numerical continuation method
		outside the regularity assumptions used in the
		global convergence theorem.}
	
	%--------------------------------------------------------------
	\medskip \noindent \textbf{Continuation Setup and Branch Initialization}
	%--------------------------------------------------------------
	
	For this example, the decision dimension is \(n=4\), the number of objectives is \(p=4\), the number of equality constraints is \(s=2\), and the number of inequality constraints is \(m=9\), including the box constraints. The homotopy state vector is initialized as $\bm{\omega}^0=[\bm{x}^0,\bb{w}^0,\bm{u}^0,\bm{v}^0]^T,$
	with initial homotopy parameter \(t_0=1\), equality multipliers \(\bm{v}^0=\mathbf{0}\in\mathbb{R}^2\), and inequality multipliers \(\bm{u}^0=[0.05\,\mathbf{1}]\in\mathbb{R}^9\).
	
	The objective-weight vectors are selected from a positive discretization of the four-objective simplex using discretization level of 10 and lower bound \(0.03\). Thus each weight vector satisfies
	\[
	\textbf{w}^0=(\text{w}_1,\text{w}_2,\text{w}_3,\text{w}_4),\qquad
	\sum_{i=1}^4 \text{w}_i=1,\qquad
	\text{w}_i\ge0.03.
	\]
	
	To explore disconnected branches, the continuation is restarted from multiple feasible initial points near both disconnected feasible components. The restart points are centered near \(x_1=\pm1\), with
	\[
	x_2\in\{0.02,0.04,0.06,0.08,0.10,0.12,0.15,-0.04,-0.08,-0.12,-0.15\}.
	\]
	The equality manifold \(x_3^2-x_2^2=0\) is initialized using both branches \(x_3=x_2\) and \(x_3=-x_2\), with \(x_4=0\). Consequently, the restart points take the form
	$\bm{x}=(\pm1,x_2,\pm x_2,0)^T,$
	subject to the strict interior feasibility condition
	$(x_1^2-1)^2+x_2^2<0.08.$
	The continuation paths are advanced using a pseudo-arclength predictor--corrector strategy. The initial continuation step size is \(\alpha_0=0.06\), with bounds \(\alpha_{\min}=10^{-5}\) and \(\alpha_{\max}=0.10\). The stopping tolerance for the homotopy parameter is \(\epsilon_t=10^{-4}\), and the maximum number of continuation steps per path is \(k_{\max}=900\).
	
	At each continuation step, the tangent direction is computed from the null space of the homotopy Jacobian and oriented consistently using the previous tangent vector. Newton-type corrector iterations are then applied to the augmented pseudo-arclength system with maximum Newton iterations \(15\), Newton tolerance \(10^{-9}\), and a damped line search. A trial step is accepted only if the corrected point is finite, the homotopy residual is sufficiently small, and the continuation path progresses toward \(t=0\). The step size is increased following rapid Newton convergence and reduced following slow convergence or failed corrections.
	
	
	%--------------------------------------------------------------
	\medskip \noindent \textbf{Pareto-Critical Structure and Continuation Results}
	%--------------------------------------------------------------
	
	The following figures, retained from the original numerical study, illustrate the geometry of the computed Pareto-critical solution set, the disconnected feasible components, the equality-constraint structure, and the objective-space behavior of the proposed homotopy continuation framework.
	%%%%%%%%%%%%%%%%%%%%%%%%%%%%%%%%%%%%%%%%%%%%%%%%%%%%%%%%%%%%%%%%%%%%%%%%%%%
	\begin{figure}[H]
		\centering
		\includegraphics[width=0.45\textwidth]{Code/DegeneratedExample/manyobj_3d_projection_colored_by_f4.png}
		\includegraphics[width=0.45\textwidth]{Code/DegeneratedExample/manyobj_nonconvex_cross_sections_x1.png}
		\caption{
			Left: Three-dimensional projection of the computed Pareto-critical set in the $(f_1,f_2,f_3)$ objective space, with $f_4$ encoded by color. Right: Cross-sections of the four objective functions along the reduced manifold $x_2=x_3=x_4=0$.
		}
		\label{fig:manyobj_3d_projection_colored}
	\end{figure}
	
	The left panel of Fig.~\ref{fig:manyobj_3d_projection_colored} illustrates the geometric structure of the computed Pareto-critical set in objective space. The disconnected vertical clusters correspond to continuation trajectories associated with the two disconnected feasible components centered near \(x_1\approx -1\) and \(x_1\approx 1\). The narrow branch-like structures indicate degeneracy in the objective mapping, where multiple feasible decision vectors collapse onto low-dimensional Pareto-critical structures in objective space. The color variation associated with \(f_4\) further reveals how the fourth objective changes along the continuation trajectories. 
	
	The right panel displays one-dimensional cross-sections of the objective functions along the reduced manifold \(x_2=x_3=x_4=0\). The oscillatory sinusoidal perturbations introduce nonconvexity and multiple competing local trade-off regions, producing a highly nonlinear many-objective optimization landscape. These competing objective structures contribute to the disconnected and branched Pareto-critical behavior observed in the continuation results.
	
	
	Fig.~\ref{fig:manyobj_pairwise_objective} illustrates the pairwise trade-off structure among the four competing objectives. The disconnected vertical branches visible in several projections reflect the presence of disconnected feasible components and multiple continuation branches generated by the homotopy framework. The color distributions indicate that the omitted objectives vary nonuniformly along the Pareto-critical trajectories, demonstrating the coupled many-objective nature of the problem. 
	
	The projections involving \(f_1\) and \(f_2\) exhibit strong separation between disconnected branches associated with the two feasible regions centered near \(x_1\approx -1\) and \(x_1\approx 1\). In contrast, the \((f_3,f_4)\) projection reveals a nearly one-dimensional structure, indicating degeneracy in the objective mapping and strong correlation between these objectives along the computed Pareto-critical set. Overall, the pairwise projections confirm that the proposed continuation framework successfully captures disconnected, nonlinear, and degenerate Pareto-critical structures in objective space.
	%%%%%%%%%%%%%%%%%%%%%%%%%%%%%%%%%%%%%%%%%%%%%%%%%%%%%%%%%%%%%
	
	Fig.~\ref{fig:manyobj_decision_space_3d} illustrates the geometric structure of the computed Pareto-critical set in decision space and its relationship to the disconnected feasible components. The left panel shows that the continuation trajectories collapse near a low-dimensional branched manifold in the \((x_1,x_2,x_3)\) subspace, reflecting the degeneracy introduced by the equality constraint \(x_3^2-x_2^2=0\). The two visible branches correspond to the manifolds \(x_3=x_2\) and \(x_3=-x_2\), while the concentration of solutions near \(x_2=x_3=0\) indicates partial collapse of the Pareto-critical structure in decision space.
	
	The right panel displays the disconnected feasible components induced by the nonlinear inequality constraint \((x_1^2-1)^2+x_2^2\le0.08\). The continuation framework successfully traces Pareto-critical branches associated with both disconnected regions near \(x_1\approx -1\) and \(x_1\approx 1\). These results demonstrate the capability of the proposed continuation strategy to explore disconnected feasible regions and branched Pareto-critical structures within a unified homotopy framework.
	%%%%%%%%%%%%%%%%%%%%%%%%%%%%%%%%%%%%%%%%%%%%%%%%%%%%%%%%%%%%%%%%%%%%%%%%%%%
	\begin{figure}[H]
		\centering
		\includegraphics[width=0.95\textwidth]{Code/DegeneratedExample/manyobj_pairwise_objective_projections_colored_by_omitted_objective.png}
		\caption{
			Pairwise objective-space projections of the computed Pareto-critical set. Each subplot displays a two-objective projection, while the omitted objective is represented by color.
		}
		\label{fig:manyobj_pairwise_objective}
	\end{figure}
	%%%%%%%%%%%%%%%%%%%%%%%%%%%%%%%%%%%%%%%%%%%%%%%%%%%%%%%%%%%%%%%%%%%%%%%%%%%
	\begin{figure}[H]
		\centering
		\includegraphics[width=0.45\textwidth]{Code/DegeneratedExample/manyobj_decision_space_x1_x2_x3.png}
		\includegraphics[width=0.45\textwidth]{Code/DegeneratedExample/manyobj_3d_projection_components.png}
		\caption{Left: Three-dimensional decision-space projection $(x_1,x_2,x_3)$ of the computed Pareto-critical set. Right: Disconnected continuation branches associated with the two disconnected feasible components near \(x_1\approx -1\) and \(x_1\approx 1\).}
		\label{fig:manyobj_decision_space_3d}
	\end{figure}
	%%%%%%%%%%%%%%%%%%%%%%%%%%%%%%%%%%%%%%%%%%%%%%%%%%%%%%%%%%%%%%%%%%%%%%%%%%%
	\begin{figure}[H]
		\centering
		\includegraphics[width=0.4\textwidth]{Code/DegeneratedExample/manyobj_decision_space_x1_x4_equality.png}
		\includegraphics[width=0.45\textwidth]{Code/DegeneratedExample/manyobj_constraint_feasibility_diagnostics.png}
		\caption{
			Left: Equality-constraint structure for \(h_1=x_4=0\). Right: Constraint feasibility diagnostics for the computed Pareto-critical solutions.
		}
		\label{fig:manyobj_feasibility}
	\end{figure}
	Fig.~\ref{fig:manyobj_feasibility} verifies the feasibility of the computed Pareto-critical solutions. The left panel shows that all computed points satisfy the equality constraint \(x_4=0\) to near machine precision, while remaining associated with the two disconnected feasible components near \(x_1\approx -1\) and \(x_1\approx 1\). This confirms that the continuation trajectories remain on the prescribed equality manifold.
	
	The right panel reports the equality residuals and inequality violations along the computed continuation solutions. The plotted residuals illustrate constraint satisfaction in the original continuation run; the updated benchmark feasibility statistics are reported below. These diagnostics confirm that the computed Pareto-critical points are not merely approximate objective-space samples, but feasible KKT-consistent solutions of the constrained many-objective problem.
	%%%%%%%%%%%%%%%%%%%%%%%%%%%%%%%%%%%%%%%%%%%%%%%%%%%%%%%%%%%%%%%%%%%%%%%%%%%
	%-------------------------------------------------------------
	\medskip \noindent \textbf{Quality Indicators and Feasibility Diagnostics}
	%--------------------------------------------------------------
	
	Because the degenerate disconnected many-objective problem is specifically designed to test approximation quality and coverage of complicated Pareto-critical structures, additional set-based quality indicators are included for this example. 
	\add{We use the hypervolume indicator
		\cite{zitzler1999multiobjective}, generational
		distance (GD) \cite{vanveldhuizen1999multiobjective},
		inverted generational distance (IGD)
		\cite{ishibuchi2015modified}, and averaged
		Hausdorff distance \cite{schutze2012using}
		to assess the quality of the computed
		approximation set relative to a sampled
		reference set.}
	
	%-------------------------------------------------------------
	\begin{table}[h]
		\centering
		\caption{Quality indicators, feasibility diagnostics, and interpretation of the computed Pareto-critical approximation set.}
		\label{tab:quality_indicators_manyobj}
		\renewcommand{\arraystretch}{1.15}
		\begin{tabular}{p{5.2cm} p{2.2cm} p{5.5cm}}
			\toprule
			\textbf{Indicator} & \textbf{Value} & \textbf{Interpretation} \\
			\midrule
			Number of feasible successful solutions 
			& 129 
			& Moderate-size approximation of the disconnected Pareto-critical structure. \\\hline
			Normalized hypervolume 
			& 0.336809712
			& Measures dominated objective-space coverage of the computed approximation set. \\\hline
			Averaged Hausdorff distance 
			& 0.246220894 
			& Moderate discrepancy between the computed approximation and the sampled reference structure. \\\hline
			Generational distance (GD) 
			& 0.007366698 
			& Small GD indicates strong local approximation accuracy. \\\hline
			Inverted generational distance (IGD) 
			& 0.246220894?????????????????
			& Larger IGD suggests incomplete global coverage of disconnected branches. \\\hline
			Minimum reported constraint violation
			& $8.08\times10^{-28}$
			& Smallest aggregated constraint violation among the homotopy solutions. \\\hline
			Coverage of disconnected components
			& Both
			& Feasible solutions were found on both components. \\
			\bottomrule
		\end{tabular}
	\end{table}
	
	Table~\ref{tab:quality_indicators_manyobj} summarizes quality indicators and feasibility diagnostics for the computed Pareto-critical approximation set. The normalized hypervolume and distance-based indicators (GD, IGD, and averaged Hausdorff distance) suggest that the proposed continuation framework captures a representative approximation of the disconnected Pareto-critical structure while maintaining reasonable diversity across the computed branches. The relatively small generational distance indicates that the computed solutions closely approximate the sampled reference structure locally, while the larger inverted generational distance reflects the difficulty of achieving complete coverage of the disconnected Pareto-critical set through local continuation alone. Furthermore, since the reference set is generated through dense feasible sampling rather than from an exact analytic Pareto-front formula, the quality indicators should be interpreted relative to the common sampled reference approximation.
	
	
	\add{These indicators characterize the quality of the computed approximation relative to the common sampled reference set. We also compare the method with three constrained evolutionary algorithms, as reported below.}
	
	
	\medskip \noindent \textbf{Comparison with Constrained Evolutionary Algorithms}
	
	\add{To assess performance on the original constrained formulation, we compare the homotopy method with NSGA-II, NSGA-III, and C-TAEA. Each evolutionary algorithm is run independently 20 times, using a population size of 84 and 400 generations (33,600 evaluations per run). The homotopy method is run once with the continuation settings described above. A common feasibility tolerance of $10^{-7}$ and a common reference approximation are used for evaluating feasible nondominated sets. No analytical reduction of the constraints is supplied to the evolutionary algorithms.}
	
	\begin{table}[h]
		\centering
		\caption{Feasibility and computational comparison on the original four-objective constrained problem. Evolutionary-algorithm values are based on 20 independent runs; homotopy values are from one run.}
		\label{tab:manyobj_algorithm_comparison}
		\renewcommand{\arraystretch}{1.1}
		\begin{tabular}{lrrrr}
			\toprule
			Method & Feasible runs & Feasible points & Best violation & Mean time (s) \\
			\midrule
			Homotopy & $1/1$ & $129$ & $8.08\times10^{-28}$ & $1836.9$ \\
			NSGA-II & $0/20$ & $0$ & $1.75\times10^{-6}$ & $9.7$ \\
			NSGA-III & $0/20$ & $0$ & $2.10\times10^{-6}$ & $13.3$ \\
			C-TAEA & $0/20$ & $0$ & $9.13\times10^{-7}$ & $25.4$ \\
			\bottomrule
		\end{tabular}
	\end{table}
	
	\add{The homotopy method produced 129 feasible nondominated points, including points on both disconnected components. None of the evolutionary runs yielded feasible final-population points at the prescribed tolerance, so feasible-front hypervolume and distance indicators are not defined for these runs. Although homotopy achieves better feasibility under these settings, its computational time is substantially greater, and the different algorithmic evaluation budgets preclude an equal-cost efficiency claim. The experiment therefore demonstrates successful numerical handling of this degenerate example, not universal superiority over evolutionary optimization methods.}
	
	The feasibility diagnostics further demonstrate the numerical stability of the continuation framework. The updated benchmark reports a minimum aggregated constraint violation of $8.08\times10^{-28}$ for the homotopy solutions. This supports numerical feasibility; the homotopy residual and multiplier conditions should be examined separately when claiming KKT consistency. Because this example contains degeneracy and disconnected Pareto-critical structures, the reported hypervolume and distance metrics should be interpreted as quality measures of the computed approximation set rather than as guarantees of global Pareto optimality. \del{Furthermore, since the reference set is generated through dense feasible sampling rather than from an exact analytic Pareto-front formula, the quality indicators should be interpreted relative to the common sampled reference approximation.}
	%===========================================================
	\begin{example}[\textbf{Three-Objective IEEE 5-Bus Optimal Power Flow Problem}]
		\label{ieee5bus}
		
		\begin{figure}[H]
			\centering
			\includegraphics[scale=0.25]{Code/ieee5bus_new.png}
			\caption{IEEE 5-bus test system with five generators, demands of 300 MW at bus 3, 300 MW at bus 4, and 400 MW at bus 5, transmission-line reactances, and thermal limits on lines $(1,2)$ and $(4,5)$. Bus 4 is selected as the reference (slack) bus with $\theta_4=0$.}
			\label{fig:ieee5bus}
		\end{figure}
		%---------------------------------
		
		
		\del{This example considers the IEEE 5-bus transmission
			network for a multiobjective optimal power flow (MOPF)
			problem.}
		\add{This example considers a five-bus transmission
			network for a multiobjective optimal power flow
			(MOPF) problem, using the MATPOWER framework
			\cite{zimmerman2011matpower}. The three-objective
			formulation is constructed to investigate generation
			cost, dispatch tracking, and transmission-line stress.}
		The system consists of five buses connected by six transmission lines, where each bus is associated with a voltage phase angle \(\theta_i\), and each generator bus has a real-power generation variable \(P_i\) measured in megawatts (MW). The decision vector is
		\[
		\bm{x}=
		\begin{bmatrix}
			P_1 & P_2 & P_3 & P_4 & P_5 &
			\theta_1 & \theta_2 & \theta_3 & \theta_4 & \theta_5
		\end{bmatrix}^{T}.
		\]
		The network includes five dispatchable
		generators together with fixed loads of
		300 MW at \del{buses 2 and 3 and 400 MW at bus 4.}
		\add{buses 3 and 4 and 400 MW at bus 5.} 
		Figure~\ref{fig:ieee5bus} provides a compact but representative benchmark for studying the interaction among economic dispatch, preferred operating points, and transmission congestion. The resulting MOO problem is
		%
		\begin{equation}
			\label{eq:MOPF_compact}
			\begin{aligned}
				\min_{\bm{x}}\quad &F(\bm{x})=\bigl(f_1(\bm{x}),f_2(\bm{x}),f_3(\bm{x})\bigr)\\[2mm]
				\text{s.t.}\quad
				&P_1+P_2
				-100\left(
				224.7319\theta_1
				-35.5872\theta_2
				-32.8947\theta_4
				-156.25\theta_5
				\right)=0,\\
				&P_3-300
				-100\left(
				-35.5872\theta_1
				+128.1798\theta_2
				-92.5926\theta_3
				\right)=0,\\
				&P_4-300
				-100\left(
				-92.5926\theta_2
				+126.2626\theta_3
				-33.67\theta_4
				\right)=0,\\
				&P_5-400
				-100\left(
				-32.8947\theta_1
				-33.67\theta_3
				+100.2348\theta_4
				-33.67\theta_5
				\right)=0,\\
				&-100\left(
				-156.25\theta_1
				-33.67\theta_4
				+189.92\theta_5
				\right)=0,\\
				&-\pi\le \theta_i\le \pi,\qquad i=1,\dots,5,\\
				&\theta_4=0,\\
				&-400\le 100(35.5872\,\theta_1-35.5872\,\theta_2)\le 400,\\
				&-240\le 100(33.67\,\theta_4-33.67\,\theta_5)\le 240,\\
				&0\le P_1\le 40,\quad
				0\le P_2\le 170,\quad
				0\le P_3\le 520,\\
				&0\le P_4\le 200,\quad
				0\le P_5\le 600.
			\end{aligned}
		\end{equation}
		The objectives represent generation cost, dispatch tracking, and transmission-line stress:
		\begin{equation}
			\label{eq:objectives_vector}
			\begin{aligned}
				f_1(\bm{x})
				&=14P_1+15P_2+30P_3+40P_4+10P_5,\\
				f_2(\bm{x})
				&=(P_1-40)^2+(P_2-170)^2+(P_3-323.4948)^2\\
				&\quad +(P_4-0)^2+(P_5-466.5052)^2,\\
				f_3(\bm{x})
				&=\left(\frac{100(35.5872\,\theta_1-35.5872\,\theta_2)}{400}\right)^2
				+\left(\frac{100(33.67\,\theta_4-33.67\,\theta_5)}{240}\right)^2.
			\end{aligned}
		\end{equation}
		The first objective minimizes generation cost using the MATPOWER \texttt{GENCOST} coefficients~\cite{zimmerman2011matpower}. The second objective penalizes deviations from a preferred DC optimal power flow dispatch,
		\[
		(P_1^{\mathrm{ref}},P_2^{\mathrm{ref}},P_3^{\mathrm{ref}},P_4^{\mathrm{ref}},P_5^{\mathrm{ref}})
		=(40,170,323.4948,0,466.5052).
		\]
		The third objective penalizes normalized power flows on the thermally constrained lines \((1,2)\) and \((4,5)\), thereby acting as a convex proxy for network congestion.
	\end{example}
	
	
	The formulation is based on the DC power flow approximation, where real-power flows are modeled using voltage phase-angle differences and line reactances:
	\[
	P_{ij}=\frac{\theta_i-\theta_j}{x_{ij}}.
	\]
	Bus 4 is selected as the slack bus with \(\theta_4=0\), and the first five equality constraints enforce power balance at each bus. The coefficients appearing in the balance equations correspond to the scaled susceptance matrix derived from the line reactances shown in Fig.~\ref{fig:ieee5bus}. Two transmission lines are subject to explicit thermal limits: line \((1,2)\) is limited to 400 MW, while line \((4,5)\) is limited to 240 MW.
	
	%===========================================================
	Below we present the numerical results for the multi-objective optimal power flow problem on the IEEE 5-bus system using the proposed homotopy method. Since the problem involves three competing objectives: generation cost, dispatch tracking, and network stress, the most informative results are those that illustrate the tradeoffs among these objectives and the corresponding changes in the physical operating point of the system. In particular, how the Pareto-critical points are distributed in objective space, how the generator dispatch changes across the Pareto front, how the constrained transmission lines behave, and how reliably the proposed homotopy method traces the solution set are presented.
	
	%--------------------------------------------------------------
	\medskip\noindent \textbf{Pareto Front Geometry and Objective Tradeoffs}
	%--------------------------------------------------------------
	
	\begin{figure}[htbp]
		\centering
		\includegraphics[width=0.45\textwidth]{Code/ExampleIEEE5Bus/3D_Pareto_Front.png}
		\includegraphics[width=0.45\textwidth]{Code/ExampleIEEE5Bus/f1_vs_f2.png}
		
		\includegraphics[width=0.45\textwidth]{Code/ExampleIEEE5Bus/f1_vs_f3.png}
		\includegraphics[width=0.45\textwidth]{Code/ExampleIEEE5Bus/f2_vs_f3.png}
		\caption{Three-dimensional Pareto front and pairwise objective projections for the IEEE 5-bus multi-objective optimal power flow problem. 
			The upper-left panel shows the Pareto surface in the objective space $(f_1,f_2,f_3)$, where each point represents a nondominated operating solution obtained using the homotopy-based method. 
			The remaining panels illustrate the pairwise projections $(f_1,f_2)$, $(f_1,f_3)$, and $(f_2,f_3)$, providing additional insight into the trade-off structure between generation cost, dispatch tracking performance, and transmission-line stress. 
			Representative solutions with original objective values are reported in Table~\ref{tab:representative-solutions}.}
		\label{fig:pareto3d}
	\end{figure}
	
	\begin{table}[h!]
		\centering
		\caption{Representative Pareto solutions with original objective values.}
		\label{tab:representative-solutions}
		\begin{tabular}{lrrrrrrr}
			\toprule
			Solution & $\text{w}_1$ & $\text{w}_2$ & $\text{w}_3$ & $t$ & $f_1$ & $f_2$ & $f_3$ \\
			\midrule
			Minimum-cost & 0.450000 & 0.075000 & 0.475000 & 0.000925 & 16887.555309 & 1754.176041 & 1.012251 \\
			Best-tracking & 0.075000 & 0.825000 & 0.100000 & 0.000305 & 17470.956598 & 0.402755 & 0.914419 \\
			Minimum-stress & 0.075000 & 0.825000 & 0.100000 & 0.000305 & 17470.956598 & 0.402755 & 0.914419 \\
			Balanced compromise & 0.375000 & 0.150000 & 0.475000 & 0.000850 & 17232.425284 & 306.158722 & 0.953713 \\
			\bottomrule
		\end{tabular}
	\end{table}
	Figure~\ref{fig:pareto3d} illustrates the trade-off structure among the three competing objectives. 
	\del{The computed Pareto front appears smooth
		and continuous, indicating stable interactions
		among the objectives and consistent numerical
		behavior of the proposed homotopy framework.}
	\add{The computed Pareto-front approximation
		exhibits a smooth geometric structure in
		objective space, suggesting consistent
		trade-off behavior among the objectives
		over the sampled solutions.}
	
	The geometry of the surface reveals the compromise relationships between generation cost, dispatch tracking performance, and network stress. 
	In particular, \del{the knee region corresponds to operating points}
	\add{the apparent knee region corresponds to operating points} where small improvements in one objective produce relatively large degradations in the others, making these solutions especially attractive from a practical decision-making perspective.
	
	The pairwise projections provide a clearer interpretation
	of the interactions between individual objective pairs.
	\del{Strong agreement is observed between the homotopy solutions
		and the Weighted Sum solutions across all projections,
		providing additional numerical validation of the proposed method.}
	
	For the $(f_1,f_2)$ projection, a pronounced nonlinear trade-off is observed between generation cost and dispatch tracking performance. 
	In particular, reducing the generation cost below approximately $1.72\times10^4$ leads to a rapid increase in the dispatch-tracking error, indicating a strong conflict between economic operation and reference-tracking objectives.
	
	The $(f_1,f_3)$ projection exhibits an approximately monotone and nearly linear trade-off between generation cost and network stress. 
	This behavior suggests that economically favorable operating points tend to increase transmission loading and overall network stress.
	
	Finally, the $(f_2,f_3)$ projection reveals a positive correlation between dispatch tracking performance and network stress. 
	Solutions with improved dispatch tracking generally correspond to lower network stress levels, indicating partial compatibility between these two objectives over the computed Pareto set.
	
	
	\del{The homotopy solutions exhibit strong agreement with the
		Weighted Sum solutions across both the three-dimensional
		Pareto surface and the pairwise objective projections,
		providing additional numerical validation of the proposed
		continuation framework. Moreover, the homotopy method
		generates a denser and more uniform distribution of
		Pareto-critical solutions, particularly near the knee
		region where the trade-offs among the competing objectives
		are most significant. The resulting Pareto structure
		appears smooth and predominantly convex, indicating
		stable and well-behaved interactions among the objectives
		for this benchmark multiobjective optimal power flow
		problem.}
	
	\add{Figure~\ref{fig:pareto3d} illustrates the
		homotopy-generated approximation of the Pareto front.
		The comparison with the Weighted Sum method is reported
		separately through the numerical results in
		Appendix~A and Table~\ref{tab:constraints_feas_opf}.
		These results allow comparison of the computed solutions
		and their constraint feasibility, although they do not
		establish that the two methods produce identical
		decision vectors or Pareto-front approximations.}
	Moreover, the homotopy method generates a denser and more uniform distribution of Pareto-critical solutions, particularly near the knee region where the trade-offs among the competing objectives are most significant. The resulting Pareto structure appears smooth and predominantly convex, indicating stable and well-behaved interactions among the objectives for this benchmark multiobjective optimal power flow problem.
	
	
	%=-----------------------------------------------------------
	\noindent \textbf{Representative Pareto Operating Points}
	%=-----------------------------------------------------------
	To better understand the physical significance of the Pareto-optimal operating points, four representative solutions are selected from the computed Pareto front:
	\begin{itemize}
		\item the minimum-cost solution,
		\item the minimum-stress solution,
		\item the best-tracking solution,
		\item a balanced compromise solution.
	\end{itemize}
	
	The corresponding objective values are reported in Table~\ref{tab:representative-solutions}. 
	The minimum-cost solution achieves the smallest generation cost objective \(f_1\), but exhibits significantly larger dispatch-tracking and stress objectives. This behavior indicates that aggressively minimizing generation cost leads to deterioration in both tracking performance and network stress characteristics.
	
	In contrast, the best-tracking and minimum-stress solutions \del{are located very close to one another on the Pareto front, suggesting a strong correlation between the tracking objective \(f_2\) and the stress objective \(f_3\) for this test system}\add{coincide to the
		reported numerical precision, indicating that the same computed operating point attains the smallest values of both
		$f_2$ and $f_3$ among the reported solutions}. Improving dispatch tracking simultaneously reduces transmission-line stress, although at the expense of increased generation cost.
	
	The balanced compromise solution lies near the knee region of the Pareto front, where substantial improvements in tracking and stress performance can still be achieved with relatively moderate increases in generation cost. Beyond this region, additional improvements in \(f_2\) and \(f_3\) require disproportionately larger increases in \(f_1\), indicating diminishing returns in the tradeoff structure.
	
	The generator dispatches associated with the representative Pareto solutions are reported in Table~\ref{tab:dispatches}. For this test system, the minimum-stress and best-tracking solutions coincide numerically, indicating a strong correlation between the dispatch-tracking and network-stress objectives over the computed Pareto set. These dispatch values provide additional physical interpretation of how the competing objectives redistribute generation across the network.
	
	\begin{table}[htbp]
		\centering
		\caption{Generator dispatches for representative Pareto solutions of the IEEE 5-bus system.}
		\label{tab:dispatches}
		\begin{tabular}{lrrrrr}
			\toprule
			Solution & $P_1$ & $P_2$ & $P_3$ & $P_4$ & $P_5$ \\
			\midrule
			Minimum-cost & 39.9949 & 169.9939 & 293.8798 & 0.0004 & 496.1311 \\
			Best-tracking & 39.9905 & 169.9889 & 323.0514 & 0.0007 & 466.9685 \\
			Minimum-stress & 39.9905 & 169.9889 & 323.0514 & 0.0007 & 466.9685 \\
			Balanced compromise & 39.9944 & 169.9932 & 311.1235 & 0.0004 & 478.8885 \\
			\bottomrule
		\end{tabular}
	\end{table}
	
	For all representative solutions, generators \(P_1\) and \(P_2\) remain close to their \del{lower}\add{upper} operating limits, \del{indicating that the optimization framework consistently prefers these lower-cost generation units}\add{consistent
		with their relatively low generation costs}. In contrast, the primary redistribution of generation occurs between generators \(P_3\) and \(P_5\), which absorb most of the variation required to balance the competing objectives.
	
	The minimum-cost solution allocates substantially more power to generator \(P_5\) while reducing the output of generator \(P_3\). This operating strategy minimizes the economic objective \(f_1\), but results in significantly larger dispatch-tracking and stress objectives. In particular, the increased reliance on generator \(P_5\) produces larger transmission flows and greater deviation from the reference dispatch profile.
	
	The best-tracking and minimum-stress solutions exhibit nearly identical dispatch patterns, further confirming the strong correlation between the tracking objective \(f_2\) and the stress objective \(f_3\). Relative to the minimum-cost solution, these operating points increase the output of generator \(P_3\) while reducing the output of generator \(P_5\), thereby improving both dispatch-tracking performance and network stress characteristics.
	
	The balanced compromise solution lies between these two extremes. Its dispatch values demonstrate how moderate redistribution between generators \(P_3\) and \(P_5\) can simultaneously maintain reasonable economic performance while achieving substantial improvements in tracking and stress objectives. This behavior is consistent with the knee region observed on the Pareto front, where favorable tradeoffs among the competing objectives are obtained.
	
	%%%%%%%%%%%%%%%%%%%%%%%%%%%%%%%%%%%%%%%%%%%%%%%%%%%%%%%%%%
	\medskip \noindent \textbf{Constraint Feasibility and Numerical Accuracy}
	%%%%%%%%%%%%%%%%%%%%%%%%%%%%%%%%%%%%%%%%%%%%%%%%%%%%%%%%%%
	
	
	Additional numerical solution details, including the computed decision vectors and objective values for the compared methods, are provided in Appendix~A.
	Table~\ref{tab:constraints_feas_opf} summarizes the numerical feasibility and constraint satisfaction properties of the computed Pareto solutions. 
	\del{The reported equality residual corresponds
		to the maximum norm of the power-balance
		constraint residuals associated with the
		equality system $\mathbf{h}(\bm{x})=0$,}
	\add{The reported equality residual is the
		maximum absolute residual over all equality
		constraints, including the power-balance
		equations and the slack-bus condition,} while the inequality violation measures the maximum violation among all operational inequality constraints \(\mathbf{g}(\bm{x})\leq0\).
	\begin{table}[h!]
		\centering
		\caption{Constraint feasibility comparison.}
		\label{tab:constraints_feas_opf}
		\begin{tabular}{lcc}
			\toprule
			Method & Max Equality Residual & Max Inequality Violation \\
			\midrule
			Homotopy Continuation & $1.5\times10^{-10}$ & $0$ \\
			Weighted Sum & $4.5\times10^{-4}$ & $0$ \\
			\bottomrule
		\end{tabular}
	\end{table}
	
	\del{The proposed homotopy framework maintained equality
		residuals near machine precision throughout the
		continuation process while satisfying all inequality
		constraints. These results confirm that the computed
		Pareto solutions are numerically accurate and physically
		realizable within the operational limits of the
		IEEE 5-bus system.}
	\add{Table~\ref{tab:constraints_feas_opf} reports the
		maximum constraint residuals obtained by the two
		numerical implementations. The proposed homotopy
		method achieved a smaller equality residual under
		the computational settings used, while neither method exhibited a positive
		inequality-constraint violation at the reported accuracy. However, the difference in equality
		residuals may reflect solver tolerances and
		termination criteria. In particular, further tuning
		of the scalar optimization solver could improve
		the equality-constraint accuracy of the Weighted
		Sum method. Therefore, these results demonstrate
		the feasibility accuracy achieved by the present
		implementations, rather than an intrinsic numerical
		advantage of one optimization formulation over
		the other.}
	\add{These results indicate that the computed
		solutions satisfy the specified DC power-flow
		model and its operational constraints to
		high numerical accuracy.}
	
	
	%%%%%%%%%%%%%%%%%%%%%%%%%%%%%%%%%%%%%%%%%%%%%%%%%%%%%%%%%%%
	
	\begin{example}[\textbf{Three-Objective Problems with Analytically Known Pareto Fronts}]
		\label{EX4}
		
		\add{We consider two three-objective optimization problems with the same
			feasible region but different objective functions. These problems
			allow us to investigate the effects of Pareto-front geometry and
			initialization on the proposed homotopy continuation method}
		
		\add{Let $\bm{x}=(x_1,x_2)^T$ and define the common feasible set
			\begin{equation}
				\label{eq:example4_feasible}
				\Omega=
				\left\{
				\bm{x}\in\mathbb{R}^2:
				x_1\geq0,\;
				x_2\geq0,\;
				x_1+x_2\leq1
				\right\}.
			\end{equation}
			The feasible set $\Omega$ is a nonempty, compact triangle,
			so assumption (A) holds. Writing the constraints as
			$g_1(\bm{x})=-x_1\leq0$,
			$g_2(\bm{x})=-x_2\leq0$, and
			$g_3(\bm{x})=x_1+x_2-1\leq0$, their gradients are
			$(-1,0)^T$, $(0,-1)^T$, and $(1,1)^T$, respectively.
			At most two constraints are active at any feasible point,
			and the active gradients are linearly independent.
			Thus assumption (B) holds.}
		
		\add{Moreover, for the standard outward normal cone
			$\mathcal N_\Omega(\bm{x})$, any
			$\bm{x}^0\in\Omega^0$ satisfies assumption (C).
			Indeed, if
			$\bm{x}^0-\bm{x}\in\mathcal N_\Omega(\bm{x})$
			for some $\bm{x}\in\Omega$, the normal-cone inequality
			implies
			$\|\bm{x}^0-\bm{x}\|^2\leq0$,
			and hence $\bm{x}=\bm{x}^0$.
			This excludes boundary points and establishes
			(C)(i). Since
			$\mathcal N_\Omega(\bm{x})=\{\bm{0}\}$
			for $\bm{x}\in\Omega^0$, condition (C)(ii) follows.
			Therefore, both problems satisfy assumptions
			(A)--(C).}
		
		
		\add{\medskip
			\noindent\textbf{(a) Convex MOO problem.}
			Consider the three-objective problem
			\begin{equation}
				\label{eq:example4_convex}
				\min_{\bm{x}\in\Omega}
				F^+(\bm{x})
				=
				\bigl(x_1^2,\;x_2^2,\;(1-x_1-x_2)^2\bigr).
			\end{equation}
			Every feasible point is Pareto optimal. Indeed, suppose that two
			feasible points $\bm{x}$ and $\widehat{\bm{x}}$ satisfy
			$F^+(\bm{x})\leq F^+(\widehat{\bm{x}})$ componentwise.
			Since $x_1,x_2,\widehat{x}_1,\widehat{x}_2\geq0$, the first two
			objective inequalities imply
			$x_1\leq\widehat{x}_1$ and
			$x_2\leq\widehat{x}_2$.
			The third objective inequality, together with
			$1-x_1-x_2\geq0$ and
			$1-\widehat{x}_1-\widehat{x}_2\geq0$, gives
			$x_1+x_2\geq\widehat{x}_1+\widehat{x}_2$.
			These inequalities can hold simultaneously only if
			$\bm{x}=\widehat{\bm{x}}$. Hence the entire feasible triangle
			$\Omega$ is the Pareto set.}
		
		
		\add{Since $x_1,x_2$, and $1-x_1-x_2$ are nonnegative
			throughout $\Omega$, we have
			$\sqrt{f_1(\bm{x})}=x_1$,
			$\sqrt{f_2(\bm{x})}=x_2$, and
			$\sqrt{f_3(\bm{x})}=1-x_1-x_2$.
			Their sum is therefore one. As every feasible point
			is Pareto optimal, the exact Pareto front is
			\begin{equation}
				\label{eq:example4_convex_pf}
				\sqrt{f_1}+\sqrt{f_2}+\sqrt{f_3}=1,
				\qquad f_i\geq0.
			\end{equation}
			For any positive normalized weight vector
			$\bm{\lambda}=(\lambda_1,\lambda_2,\lambda_3)^T\in W^{++}$,
			the interior Pareto-criticality conditions are
			\[
			\lambda_1x_1
			=
			\lambda_3(1-x_1-x_2),
			\qquad
			\lambda_2x_2
			=
			\lambda_3(1-x_1-x_2).
			\]
			Using these equations and the identity
			$x_1+x_2+(1-x_1-x_2)=1$, the analytical interior
			Pareto-critical point is
			\begin{equation}
				\label{eq:example4_exact_kkt}
				x_1^*=
				\frac{\lambda_1^{-1}}
				{\sum_{j=1}^3\lambda_j^{-1}},
				\qquad
				x_2^*=
				\frac{\lambda_2^{-1}}
				{\sum_{j=1}^3\lambda_j^{-1}}.
			\end{equation}
			Furthermore, the Hessian of the corresponding positively weighted
			objective is
			\[
			\nabla_{\bm{x}}^2
			\left(\sum_{i=1}^3\lambda_i f_i\right)
			=
			2\begin{pmatrix}
				\lambda_1+\lambda_3 & \lambda_3\\
				\lambda_3 & \lambda_2+\lambda_3
			\end{pmatrix},
			\]
			which is positive definite for every $\bm{\lambda}\in W^{++}$.
			Consequently, $\bm{x}^*$ is the unique minimizer of the
			corresponding weighted objective. The analytical solution
			\eqref{eq:example4_exact_kkt} provides a reference for assessing
			the accuracy of the continuation method.}
		
		\medskip
		\noindent\textbf{(b) Nonconvex MOO problem.}
		
		Using the same feasible set $\Omega$, consider
		\begin{equation}
			\label{eq:example4_concave}
			\min_{\bm{x}\in\Omega}
			F^-(\bm{x})
			=
			\bigl(-x_1^2,\;-x_2^2,\;-(1-x_1-x_2)^2\bigr).
		\end{equation}
		Again, every feasible point is Pareto optimal. If
		$F^-(\bm{x})\leq F^-(\widehat{\bm{x}})$ componentwise,
		then the first two objective inequalities imply
		$x_1\geq\widehat{x}_1$ and
		$x_2\geq\widehat{x}_2$, while the third implies
		$x_1+x_2\leq\widehat{x}_1+\widehat{x}_2$.
		Therefore, $\bm{x}=\widehat{\bm{x}}$, and the Pareto set
		remains $\Omega$.
		
		Since $\sqrt{-f_1}=x_1$, $\sqrt{-f_2}=x_2$, and
		$\sqrt{-f_3}=1-x_1-x_2$, the exact Pareto front is
		\begin{equation}
			\label{eq:example4_concave_pf}
			\sqrt{-f_1}+\sqrt{-f_2}+\sqrt{-f_3}=1,
			\qquad f_i\leq0.
		\end{equation}
		The interior Pareto-criticality conditions again yield the
		analytical solution \eqref{eq:example4_exact_kkt}. However,
		the Hessian of the corresponding positively weighted objective is
		\[
		\nabla_{\bm{x}}^2
		\left(\sum_{i=1}^3\lambda_i f_i\right)
		=
		-2\begin{pmatrix}
			\lambda_1+\lambda_3 & \lambda_3\\
			\lambda_3 & \lambda_2+\lambda_3
		\end{pmatrix},
		\]
		which is negative definite for every $\bm{\lambda}\in W^{++}$.
		Thus, the analytical interior Pareto-critical point is a strict
		maximum of the corresponding weighted objective and cannot be
		recovered as a weighted-sum minimizer. Nevertheless, it satisfies
		the terminal Pareto-criticality conditions of the proposed
		homotopy formulation.
		
		\medskip
		\noindent\textbf{Numerical results: convex case.}
		
		For the numerical experiment, 21 positive weight vectors were selected
		from a simplex-lattice discretization of $W^{++}$, and four strictly
		feasible initial points were used for each weight, giving 84
		continuation paths. Algorithm~2 successfully followed all 84 paths to
		feasible terminal points. All 84 terminal points were within $10^{-3}$
		in objective space of their corresponding analytical interior
		Pareto-critical points, and no path terminated on the boundary.
		The median and maximum objective-space errors were
		\[
		9\times10^{-6}
		\qquad\text{and}\qquad
		1.14\times10^{-4},
		\]
		respectively. The stationarity and complementarity
		residuals are evaluated using the original
		KKT system at $t=0$, rather than the
		homotopy equations at the terminal
		continuation parameter. The maximum stationarity residual was
		$4.7\times10^{-5}$ and the maximum complementarity residual was
		$2\times10^{-6}$. Using a dense discretization of the exact Pareto
		front as the reference set,
		\[
		\mathrm{GD}=0.003195,
		\qquad
		\mathrm{IGD}=0.108137.
		\]
		The relatively larger IGD reflects the use of only 21 weight vectors
		and the nonlinear mapping from regularly distributed weights to
		Pareto-critical points; a regular weight discretization does not, in
		general, produce a uniform distribution on the Pareto front.
		
		
		\medskip
		\noindent\textbf{(b) Concave Pareto front.}
		Consider now
		\[
		\min_{\mathbf{x}}\;
		\mathbf{F}^{-}(\mathbf{x})
		=
		\begin{pmatrix}
			-z_1^2\\
			-z_2^2\\
			-z_3^2
		\end{pmatrix}
		\]
		over the same feasible region.
		
		Again, every feasible point is Pareto optimal. If $\mathbf{z}$
		dominated $\widehat{\mathbf{z}}$, then
		\[
		-z_i^2\leq-\widehat z_i^2,
		\]
		which implies $z_i\geq\widehat z_i$ for all $i$. The unit-sum
		constraint again forces equality in all three coordinates.
		
		The exact Pareto front is
		\[\sqrt{-f_1}+\sqrt{-f_2}+\sqrt{-f_3}=1, f_i\leq0.\]
		
		For each positive weight vector, the interior first-order
		Pareto-criticality equations yield
		\[z_i^*=\frac{\lambda_i^{-1}}{\displaystyle\sum_{j=1}^{3}\lambda_j^{-1}}, i=1,2,3.\]
		However,
		\[\nabla_{\mathbf{x}}^2
		\left(
		\sum_{i=1}^{3}\lambda_i f_i
		\right)
		=
		-2
		\begin{pmatrix}
			\lambda_1+\lambda_3 & \lambda_3\\
			\lambda_3 & \lambda_2+\lambda_3
		\end{pmatrix},
		\]
		which is negative definite. Thus the analytical interior
		Pareto-critical point is a strict local maximum, rather than a
		minimum, of the corresponding positively weighted objective. These
		interior Pareto points are therefore unsupported by the usual
		weighted-sum minimization.
		
		This distinction is important because the positive weights in the
		present homotopy formulation enter the Pareto-criticality/KKT system;
		the proposed method does not independently minimize a weighted sum for
		each weight. Hence unsupported Pareto-critical points are not excluded
		from the terminal homotopy system.
		
		Using the same 21 weights and four generic initial points as in part
		(a), all 84 continuation paths converged to feasible KKT points located
		on the boundary of the feasible triangle. The resulting
		\[
		\mathrm{GD}=4\times10^{-6}
		\]
		shows that the computed points lie essentially on the exact Pareto
		front, whereas
		\[
		\mathrm{IGD}=0.295819
		\]
		indicates poor coverage of the complete Pareto surface. Thus the issue
		is not convergence to points away from the Pareto front, but selection
		of boundary KKT branches rather than the unsupported interior branch.
		
		To examine this behavior more closely, five representative positive
		weight vectors were selected:
		\[
		\begin{split}
			&(1/3,1/3,1/3),\\
			&(0.25,0.375,0.375),\qquad
			(0.375,0.25,0.375),\\
			&(0.375,0.375,0.25),\qquad
			(0.25,0.25,0.50).
		\end{split}
		\]
		For each weight, the analytical interior Pareto-critical point
		$\mathbf{x}^*(\boldsymbol{\lambda})$ was first used as the initial
		point. Algorithm~2 recovered the corresponding interior
		Pareto-critical point for all five weights, with objective-space errors
		on the order of machine precision.
		
		The initial point was then perturbed according to
		\[
		\mathbf{x}^{0}
		=
		\mathbf{x}^{*}
		+
		\varepsilon\mathbf{d},
		\]
		where $\mathbf{d}$ was selected from coordinate, diagonal, and
		anti-diagonal directions and
		\[
		\varepsilon
		\in
		\left\{
		10^{-6},10^{-5},10^{-4},10^{-3},10^{-2},
		5\times10^{-2}
		\right\}.
		\]
		Forty perturbed continuation paths were considered for each value of
		$\varepsilon$. The results are summarized in
		Table~\ref{tab:example4_concave_perturbation}.
		
		\begin{table}[htbp]
			\centering
			\caption{Initialization sensitivity for the concave three-objective
				problem in Example~4.}
			\label{tab:example4_concave_perturbation}
			\begin{tabular}{c|c|c|c}
				\hline
				$\varepsilon$
				& Number of paths
				& Interior branch
				& Boundary branch\\
				\hline
				$10^{-6}$ & 40 & 40 & 0\\
				$10^{-5}$ & 40 & 40 & 0\\
				$10^{-4}$ & 40 & 37 & 3\\
				$10^{-3}$ & 40 & 22 & 18\\
				$10^{-2}$ & 40 & 0 & 40\\
				$5\times10^{-2}$ & 40 & 0 & 40\\
				\hline
			\end{tabular}
		\end{table}
		
		For the perturbations and initializations
		considered, the results reveal a clear
		transition in branch selection. All
		perturbations of magnitude $10^{-6}$ and $10^{-5}$ remained on the
		analytical interior branch. At $\varepsilon=10^{-4}$, 37 of the 40
		paths still reached the interior branch, whereas at
		$\varepsilon=10^{-3}$ only 22 did so. For perturbations of magnitude
		$10^{-2}$ or larger, all paths transitioned to boundary KKT branches.
		
		Thus the unsupported interior Pareto-critical branch is present in the
		homotopy zero set and can be accurately followed by Algorithm~2.
		However, the terminal branch selected by continuation depends on the
		initialization. This distinguishes the ability of the homotopy
		formulation to represent and trace unsupported Pareto-critical points
		from the ability of a finite collection of initializations to provide
		uniform coverage of a nonconvex Pareto front.
		
		Taken together, parts (a) and (b) also illustrate that the proposed
		continuation method is directly applicable to problems with more than
		two objectives. For each selected weight and initialization,
		Algorithm~2 follows a one-dimensional homotopy path in the enlarged
		continuation space toward a terminal Pareto-critical point. The fact
		that the three-objective Pareto front is two-dimensional does not
		require an individual continuation path to be two-dimensional.
		
	\end{example}
	%%%%%%%%%%%%%%%%%%%%%%%%%%%%%%%%%%%%%%%%%%%%%%%%%%%%%%%%%%%
	\section{Conclusion and Future Work} \label{sec:conclusion}
	This work developed a homotopy continuation framework for constrained MOO problems by embedding the Karush--Kuhn--Tucker (KKT) conditions into a parameterized continuation structure. The proposed homotopy map continuously deforms an auxiliary system into the KKT system of the original MOO problem, yielding a deterministic and theoretically grounded continuation framework for computing Pareto-critical solutions. 
	\del{Under suitable regularity assumptions,
		global convergence of the continuation
		trajectories was established.}
	\add{Under the stated regularity assumptions,
		the theoretical analysis establishes the
		existence of a smooth continuation path
		with a nonempty set of accumulation points
		satisfying the KKT conditions as the
		homotopy parameter approaches zero.}
	
	
	The proposed framework was investigated through benchmark optimization problems, disconnected nonconvex many-objective examples, and a multiobjective IEEE 5-bus optimal power flow problem. The numerical experiments demonstrated that the continuation approach can successfully compute feasible Pareto-critical solution branches while maintaining strong agreement with classical scalarization methods and competitive performance relative to NSGA-II. 
	\del{In addition, the framework exhibited robust
		behavior for disconnected, degenerate, and
		nonconvex Pareto-critical structures through
		adaptive continuation and branch-following
		strategies.}
	\add{The numerical experiments also illustrate
		the ability of adaptive continuation and
		multiple branch initializations to explore
		disconnected, degenerate, and nonconvex
		Pareto-critical structures, while revealing
		limitations in global coverage and
		sensitivity to initialization.} 
	\add{The analytically solvable three-objective
		examples further demonstrate that the
		homotopy formulation can recover unsupported
		interior Pareto-critical points that cannot
		be obtained through positive weighted-sum
		minimization. However, the numerical results
		also show that reaching these branches may
		depend strongly on the continuation
		initialization.}
	
	
	
	The results also highlight several practical challenges associated with numerical continuation methods, including step-size control, path stability near singular regions, and sensitivity to initialization. To address these issues, adaptive predictor--corrector strategies and continuation enhancements were incorporated to improve robustness and numerical stability.
	
	Several directions for future research remain. Extending the framework to large-scale and high-dimensional MOO problems is of particular interest, especially in settings where Jacobian evaluations dominate computational cost. Additional improvements may be obtained through advanced continuation techniques, parallel path tracking, and machine-learning-assisted parameter selection. Further theoretical developments involving weaker regularity assumptions and refined convergence analysis would also broaden the applicability of the proposed framework.
	
	
	\del{Overall, the present study demonstrates
		that homotopy continuation provides an
		effective and flexible deterministic
		framework for constrained MOO, particularly
		for problems involving nonconvexity,
		degeneracy, and disconnected Pareto-critical
		structures.}
	\add{Overall, the theoretical and numerical
		results support homotopy continuation as a
		flexible deterministic framework for
		computing Pareto-critical solutions of
		constrained MOO problems, including selected
		problems with nonconvex, degenerate, and
		disconnected solution structures.}
	\add{Future work will include comparisons with
		more recent constrained and many-objective
		evolutionary algorithms, including implementations
		available through platforms such as pymoo,
		jMetal, and PlatEMO.}
	
	
	\section*{Acknowledgements}
	The authors would like to thank the anonymous reviewers for their careful reading of the manuscript and their valuable comments and suggestions, which significantly improved the clarity, organization, and quality of this work.
	
	%%%%%%%%%%%%%%%%%%%%%%%%%%%%%%%%%%%%%%%%%%%%%%%%%%%%%%%%%%%
	\appendix  %This command ends the counting of sections.
	%%%%%%%%%%%%%%%%%%%%%%%%%%%%%%%%%%%%%%%%%%%%%%%%%%%%%%%%%%
	\begin{appendices}
		\section*{Appendices}
		\addcontentsline{toc}{section}{Appendices}
		The following appendices summarize the classical scalarization and evolutionary algorithms used for comparison in the numerical experiments.
		%%%%%%%%%%%%%%%%%%%%%%%%%%%%%%%%%%%%%%%%%%%%%%
		\section{Classical Scalarization Methods}
		\label{appendix:A}
		%%%%%%%%%%%%%%%%%%%%%%%%%%%%%%%%%%%%%%%%%%%%%%
		\subsection{The Weighted Sum Method (WSM)}
		The Weighted Sum Method (WSM) is arguably the most straightforward scalarization technique, transforming MOO into a single objective function by taking a linear convex combination of the objectives \cite{marler2004survey, van2012evaluation,gass1955computational}. 
		Given the $p$ objectives $f_i(\bm{x})$, the WSM formulation is defined by minimizing the aggregated objective function $F(\bm{x},\bb{w})$:
		%
		\begin{mini!}{\bm{x} \in \Omega}{F(\bm{x},\bb{w}) = \sum_{i=1}^{p} w_i f_i(\bm{x})}{}{}\label{wsm}\tag{5}
		\end{mini!}
		%
		The weight vector $\mathbf{w} = [w_1, w_2, \dots, w_p]^T$ is chosen by the decision-maker, where $w_i \geq 0$ for all $i$ and typically $\sum_{i=1}^{p} w_i = 1$. When all weights are strictly positive, any minimizer of~\eqref{wsm} is Pareto optimal. By varying the weight vector, different Pareto-optimal solutions can be generated; however, the Weighted Sum Method may fail to identify Pareto-optimal solutions located on nonconvex regions of the Pareto front. The resulting single-objective optimization problem is solved using the Sequential Least Squares Programming (SLSQP) algorithm\,\cite{nocedal2006numerical}.
		\begin{algorithm}[h!]
			\caption{Weighted Sum Method (WSM) for Pareto Front Generation}
			\label{alg:wsm}
			\begin{algorithmic}[1]
				\item[]  \hspace{-2em} \textbf{Require:} Objective functions $\{f_i(\bm{x})\}_{i=1}^p$, feasible set $\Omega$, 
				weight set $\mathcal{W} \subset \{\bm{w} \in \mathbb{R}^p \mid w_i \ge 0, \ \sum_{i=1}^p w_i = 1\}$
				\item[] \hspace{-2em} \textbf{Ensure:} Approximate Pareto set $\mathcal{P}_{\text{WSM}}$ and Pareto front $\mathcal{F}_{\text{WSM}}$
				\STATE $\mathcal{P}_{\text{WSM}} \gets \emptyset$ \{\emph{Set of decision vectors}\}
				\STATE $\mathcal{F}_{\text{WSM}} \gets \emptyset$ \{\emph{Set of objective vectors}\}
				\FOR{$\bm{w} \in \mathcal{W}$}
				\STATE Define scalarized objective:
				$F(\bm{x},\bm{w}) = \sum_{i=1}^{p} w_i f_i(\bm{x})$
				\STATE Solve the single-objective problem
				$\min_{\bm{x} \in \Omega} F(\bm{x},\bm{w})$
				and obtain a solution $\bm{x}^{\star}(\bm{w})$
				\STATE Compute objective vector 
				$\bm{f}^{\star}(\bm{w}) = \bigl(f_1(\bm{x}^{\star}(\bm{w})),\dots,f_p(\bm{x}^{\star}(\bm{w}))\bigr)$
				\STATE Add solution to candidate sets:
				$\mathcal{P}_{\text{WSM}} \gets \mathcal{P}_{\text{WSM}} \cup \{\bm{x}^{\star}(\bm{w})\}, \quad
				\mathcal{F}_{\text{WSM}} \gets \mathcal{F}_{\text{WSM}} \cup \{\bm{f}^{\star}(\bm{w})\}$
				\ENDFOR
				%\STATE Remove dominated points from $\mathcal{P}_{\text{WSM}}$ and $\mathcal{F}_{\text{WSM}}$ to retain only non-dominated solutions
				\STATE \textbf{return} $\mathcal{P}_{\text{WSM}}, \mathcal{F}_{\text{WSM}}$
			\end{algorithmic}
		\end{algorithm}
		%%%%%%%%%%%%%%%%%%%%%%%%%%%%%%%%%%%%%%%%%%%%%%%%%%%%%%%%%%%%%%
		\subsection{The $\epsilon$-Constraints Method (ECM)}
		The $\epsilon$-Constraints Method (ECM) overcomes the convexity limitation of the WSM by maintaining one objective function as the primary target for minimization while converting all remaining $p-1$ objectives into inequality constraints\,\cite{haimes1971bicriterion}.
		Assuming objective $f_m(\bm{x})$ is selected as the primary objective for minimization, the ECM is formulated as a single-objective problem subject to bounds defined by $\epsilon$:
		%
		\begin{mini!}{\bm{x} \in \Omega}{f_m(\bm{x})}{\label{eq:epsilon-constraint}}{}\tag{6}
			\addConstraint{f_k(\bm{x}) \leq \epsilon_k, \quad \text{for } k=1, \dots, p, \quad k \neq m}{}.\tag{7}
		\end{mini!}
		%
		Here, $\epsilon_k$ represents the maximum acceptable value for the constrained objective $f_k$. It can  be shown that a solution of \eqref{eq:epsilon-constraint} is weakly Pareto optimal. Furthermore, a solution $\bm{x}^*$ is Pareto optimal if and only if it is a solution of \eqref{eq:epsilon-constraint} for every $m = 1, \cdots, p$, where $\epsilon_k = f_k(\bm{x}^*)$ for $k=1, \cdots, p, k\neq m$. By iteratively minimizing $f_m(\mathbf{x})$ over a range of systematically varied $\epsilon_k$ values, the full Pareto front can be traced, including non-convex sections.
		%%%%%%%%%%%%%%%%%%%%%%%%%%%%%%%%%%%%%%%%%%%%%%%%%%%%%%%%%%%%
		\begin{algorithm}[h!]
			\caption{$\epsilon$-Constraint Method (ECM) for Pareto Front Generation}
			\label{alg:ecm}
			\begin{algorithmic}[1]
				\item[]  \hspace{-2em} \textbf{Require:} Objective functions $\{f_i(\bm{x})\}_{i=1}^p$, feasible set $\Omega$, index $m$ of the primary objective $f_m$, \\
				\hspace{2.8em} grid of $\epsilon$-vectors $\{\boldsymbol{\epsilon}^{(j)}\}_{j=1}^{N}$ 
				where 
				$\boldsymbol{\epsilon}^{(j)} = (\epsilon_1^{(j)},\dots,\epsilon_p^{(j)})$ 
				and $\epsilon_m^{(j)}$ is unused.
				% \item[] \hspace{-2em} \textbf{Require:} 
				%    Objective functions $\{f_i(\mathbf{x})\}_{i=1}^p$, feasible set $\Omega$, 
				%    primary objective index $m$, 
				%    grid of $\epsilon$-vectors $\{\boldsymbol{\epsilon}^{(s)}\}_{s=1}^{N_\epsilon}$ 
				%    where 
				%    $\boldsymbol{\epsilon}^{(s)} = (\epsilon_1^{(s)},\dots,\epsilon_p^{(s)})$ 
				%    and $\epsilon_m^{(s)}$ is unused.   
				\item[] \hspace{-2em} \textbf{Ensure:} Approximate Pareto set $\mathcal{P}_{\epsilon}$ and Pareto front $\mathcal{F}_{\epsilon}$
				\STATE $\mathcal{P}_{\epsilon} \gets \emptyset$ \{\emph{Set of decision vectors}\}
				\STATE $\mathcal{F}_{\epsilon} \gets \emptyset$ \{\emph{Set of objective vectors}\}
				% \FOR{combinations of bounds $\{\epsilon_k^{(j_k)}\}_{k \neq m}$}
				\FOR{each $\boldsymbol{\epsilon}^{(j)}$ in the grid}
				\STATE Define constrained problem:
				$\min_{\bm{x} \in \Omega} \quad & f_m(\bm{x})$                s.t. $f_k(\bm{x}) \le \epsilon_k^{(j)}, \quad k = 1,\dots,p, \ k \neq m$
				\STATE Solve the above single-objective problem and obtain a solution $\bm{x}^{\star}(\bm{\epsilon}^{(j)})$
				\STATE Compute objective vector 
				$\bm{f}^{\star}(\bm{\epsilon}^{(j)}) = \bigl(f_1(\bm{x}^{\star}),\dots,f_p(\bm{x}^{\star})\bigr)$
				\STATE Update candidate sets:
				$\mathcal{P}_{\epsilon} \gets \mathcal{P}_{\epsilon} \cup \{\bm{x}^{\star}(\bm{\epsilon}^{(j)})\}, \quad
				\mathcal{F}_{\epsilon} \gets \mathcal{F}_{\epsilon} \cup \{\bm{f}^{\star}(\bm{\epsilon}^{(j)})\}$
				\ENDFOR
				%\STATE Remove dominated points from $\mathcal{P}_{\epsilon}$ and $\mathcal{F}_{\epsilon}$ to retain only non-dominated solutions
				\STATE \textbf{return} $\mathcal{P}_{\epsilon}, \mathcal{F}_{\epsilon}$
			\end{algorithmic}
		\end{algorithm}
		%%%%%%%%%%%%%%%%%%%%%%%%%%%%%%%%%%%%%%%%%%%%%%%%%%%%%%%%%%%%%
		\subsection{The Global Criterion Method (GCM)}
		
		\del{The Global Criterion Method (GCM), also known as the
			Goal Programming or Compromise Programming method,
			transforms MOO into minimizing the distance from an
			unattainable ``Ideal Point'' ($\bm{f}^*$).
			The ideal point $\bm{f}^*$ is found by individually
			optimizing each objective function $f_i(\bm{x})$.}
		\add{The Global Criterion Method (GCM) used here is
			a particular form of the Method of Weighted Metrics,
			also known as compromise programming
			\cite{miettinen1999nonlinear}.
			It minimizes the distance between the objective
			vector and the ideal point $\bm{f}^*$, whose
			components are obtained by individually minimizing
			each objective function over the feasible set.}
		The GCM minimizes the deviation from the ideal point $\bm{f}^*$ using an $L_p$-norm distance metric:
		%
		\begin{mini!}{\bm{x} \in \Omega}{D_p(\bm{x}) = \left( \sum_{i=1}^{M} \left| \frac{f_i(\bm{x}) - f_i^*}{s_i} \right|^q \right)^{1/q}}{}{}\tag{8}
		\end{mini!}
		%
		where $f_i^*$ is the ideal value (minimum) for objective $i$, and $s_i$ is a normalization term, typically the range of the objective function ($s_i = f_i^{\max} - f_i^{\min}$) \cite{marler2004survey}. The use of $p=2$, corresponding to the $L_2$-norm (or Euclidean norm), is the most common approach, minimizing the geometric distance in the objective space. 
		\del{Despite its popularity and wide range
			of applications, there is no guarantee
			that it provides a Pareto optimal solution.}
		\add{For strictly positive scaling factors $s_i$ and a
			finite norm exponent $p\geq 1$, every global minimizer
			of this ideal-point scalarization is Pareto optimal,
			provided that $f_i^*$ is a lower bound on $f_i(\bm{x})$
			over the feasible set for every objective
			\cite{miettinen1999nonlinear}.
			This guarantee applies to global minimizers of the
			scalarized problem and does not necessarily extend
			to local stationary points returned by a numerical
			optimization solver.}
		
		%
		\begin{algorithm}[H]
			\caption{Global Criterion Method (GCM) Based on Ideal Point Deviation}
			\label{alg:gcm}
			\begin{algorithmic}[1]
				\item[] \hspace{-2em} \textbf{Require:} Objective functions $\{f_i(\bm{x})\}_{i=1}^p$, feasible set $\Omega$, norm parameter $\add{q} \ge 1$
				\item[] \hspace{-2em} \textbf{Ensure:} Approximate Pareto solution $\bm{x}^{\star}$ and objective vector $\bm{f}^{\star}$
				\item[] \hspace{-1.6em} \textbf{Compute Ideal Values}    
				\FOR{$i = 1,\dots,p$}
				\STATE Solve
				$f_i^* = \min_{\bm{x} \in \Omega} f_i(\bm{x})$
				and store the solution $\bm{x}_i^*$
				\ENDFOR
				\STATE Estimate scaling factors $s_i$ (e.g., $s_i = f_i^{\max} - f_i^{\min}$ using bounds or sampled values)
				% \STATE \textbf{Step 2: Define Global Criterion}
				\item[] \hspace{-1.6em} \textbf{Define Global Criterion}
				\STATE Define the scalar deviation measure
				$D_p(\bm{x}) = 
				\left( 
				\sum_{i=1}^{p} 
				\left| 
				\frac{f_i(\bm{x}) - f_i^*}{s_i} 
				\right|^q 
				\right)^{1/q}$
				\item[] \hspace{-1.6em} \textbf{Solve the Global Criterion Problem}
				\STATE Solve
				$\min_{\bm{x} \in \Omega} D_p(\bm{x})$
				to obtain $\bm{x}^{\star}$
				\STATE Compute the corresponding objective vector
				$\bm{f}^{\star} = \bigl(f_1(\bm{x}^{\star}),\dots,f_p(\bm{x}^{\star})\bigr)$
				\STATE \textbf{return:} $\bm{x}^{\star},\; \bm{f}^{\star}$
			\end{algorithmic}
		\end{algorithm}
		%%%%%%%%%%%%%%%%%%%%%%%%%%%%%%%%%%%%%%%%%%%%%%%%%%%%%%%%%%%%%%%%%
		\subsection{The Lexicographic Method}
		The Lexicographic Method is an \textit{a priori} approach characterized by the decision-maker assigning an absolute priority ranking to the objectives. This ranking dictates that the most important objective must be optimized first, followed sequentially by the next most important, subject to maintaining the optimal result of all higher-ranked objectives \cite{isermann1982linear, miettinen1999nonlinear, ehrgott2005multicriteria}.
		A lexicographic minimization problem defines an ordered set of functions $\bm{f}=[f_1,f_2,\dots,f_p]$, where $f_1$ is deemed the most important and $f_p$ the least important. The problem is concisely written as:
		%
		\begin{mini!}{\bm{x} \in \Omega}{\text{lex min } [f_1(\bm{x}), f_2(\bm{x}), \dots, f_p(\bm{x})]}{}{}\tag{9}
		\end{mini!}
		%
		This notation signifies that $f_1$ has overwhelming priority over $f_2$ ($f_1 \gg f_2$), and so forth down the hierarchy.
		The solution methodology for the lexicographic problem involves solving a strict sequence of $p$ single-objective optimization problems:
		%
		\begin{enumerate}[label=\textbf{Step \arabic*:}, leftmargin=*]
			\item \textbf{$k=1$:} Solve the unconstrained minimization problem for the most important objective:
			$$f_1^* = \min_{\bm{x} \in \Omega} f_1(\bm{x})$$
			%
			\item \textbf{$k=2$:} Solve the problem for the second objective, constrained by the optimal value obtained in the previous step. Define a new feasible set $\Omega_2$:
			$$\Omega_2 = \{ \bm{x} \in \Omega \mid f_1(\bm{x}) = f_1^* \}$$
			The problem becomes: $\displaystyle f_2^* = \min_{\bm{x} \in \Omega_2} f_2(\bm{x})$.
			%
			\item \textbf{$k=3, \dots, p$:} Continue this iterative process. At step $k$, the feasible set $\Omega_k$ is constrained by the optimal values $f_1^*, f_2^*, \dots, f_{k-1}^*$, ensuring the previous optimalities are maintained as hard equality constraints.
		\end{enumerate}
		%
		If any sequential problem is infeasible or unbounded, the process halts, and no solution exists for the full MOO under the specified priority structure.
		%
		\begin{algorithm}[H]
			\caption{Lexicographic Sequential Solution Algorithm}
			\begin{algorithmic}[1]
				\STATE \textbf{Input:} Objective functions $\{f_1, f_2, \dots, f_p\}$, feasible region $\Omega$
				\STATE \textbf{Initialize:} Set $k \gets 1$, $\Omega_1 \gets \Omega$
				\WHILE{$k \leq p$}
				\STATE \textbf{Step 1: Solve the $k^{th}$ Optimization Problem}
				\STATE \hspace{1em} Find the optimal value: 
				\STATE \hspace{2em} $\displaystyle f_k^* = \min_{\bm{x} \in \Omega_k} f_k(\bm{x})$
				\STATE \textbf{Step 2: Update Feasible Set}
				\STATE \hspace{1em} Define the next feasible region:
				\STATE \hspace{2em} $\Omega_{k+1} = \{ \bm{x} \in \Omega_k \mid f_k(\bm{x}) = f_k^* \}$
				\STATE \textbf{Step 3: Check Feasibility}
				\STATE \hspace{1em} If $\Omega_{k+1}$ is empty or unbounded:
				\STATE \hspace{2em} \textbf{Terminate:} No feasible solution exists under the current priority structure.
				\STATE \textbf{Step 4: Increment Priority Level}
				\STATE \hspace{1em} $k \gets k + 1$
				\ENDWHILE
				\STATE \textbf{Step 5: Termination}
				\STATE \hspace{1em} Output the solution $\bm{x}^*$ satisfying:
				\STATE \hspace{2em} $f_1(\bm{x}^*) = f_1^*, ~ f_2(\bm{x}^*) = f_2^*, ~ \dots, ~ f_p(\bm{x}^*) = f_p^*$
				\STATE \textbf{Output:} Lexicographically optimal solution $\bm{x}^*$
			\end{algorithmic}
		\end{algorithm}
		%%%%%%%%%%%%%%%%%%%%%%%%%%%%%%%%%%%%%%%%%%%%%%%%%%%%%%%%%%%%%%%%%%%%%%
		\section{Evolutionary Algorithm}
		\label{appendix:B}
		\subsection{The Non-Dominated Sorting Genetic Algorithm II (NSGA-II)}
		The Non-Dominated Sorting Genetic Algorithm II (NSGA-II) is a fast elitist multiobjective evolutionary algorithm introduced by Deb et al.~\cite{deb2002fast}. The Non-Dominated Sorting Genetic Algorithm II (NSGA-II) is an advanced, fast elitist Multi-Objective Genetic Algorithm (MOGA). It was developed to overcome the computational complexity and lack of effective elitism found in earlier non-domination-based evolutionary algorithms. NSGA-II is a highly robust meta-heuristic approach often applied to complex, large-scale problems such as optimizing resource allocation in supply chains, where it often yields better operational outcomes than traditional methods\,\cite{acerce2025application}. 
		NSGA-II utilizes an elitist approach, ensuring that the best solutions found across generations are retained. The algorithm proceeds through a standard generational cycle:
		%
		\begin{enumerate}[leftmargin=*]
			\item A parent population ($P_t$) generates an offspring population ($Q_t$) via genetic operators (selection, crossover, and mutation) \cite{acerce2025application}.
			\item The parent and offspring populations are combined into a larger population $R_t = P_t \cup Q_t$.
			\item $R_t$ is sorted based on non-domination into different Pareto fronts ($F_1, F_2, F_3, \dots$).
			\item The next generation population ($P_{t+1}$) is formed by selecting individuals sequentially, starting from the best front ($F_1$) and proceeding to subsequent fronts until the population limit $N$ is met.
		\end{enumerate}
		%
		\begin{algorithm}[H]
			\caption{Non-dominated Sorting Genetic Algorithm II (NSGA-II)}
			\begin{algorithmic}[1]
				\STATE \textbf{Input:} Population size $N$, number of generations $T$, genetic operators (selection, crossover, mutation)
				\STATE \textbf{Initialize:} Generate an initial parent population $P_0$ of size $N$
				\FOR{$t = 0$ to $T-1$}
				\STATE \textbf{Step 1: Generate Offspring Population}
				\STATE \hspace{1em} Apply genetic operators (selection, crossover, and mutation) on $P_t$ to produce offspring $Q_t$
				\STATE \textbf{Step 2: Combine Parent and Offspring Populations}
				\STATE \hspace{1em} Form the combined population $R_t = P_t \cup Q_t$
				\STATE \textbf{Step 3: Non-dominated Sorting}
				\STATE \hspace{1em} Sort $R_t$ into Pareto fronts: $F_1, F_2, F_3, \dots$
				\STATE \textbf{Step 4: Create the Next Generation}
				\STATE \hspace{1em} Initialize $P_{t+1} = \emptyset$
				\STATE \hspace{1em} Sequentially add individuals from $F_1, F_2, \dots$ into $P_{t+1}$ until $|P_{t+1}| = N$
				\STATE \textbf{Step 5: Apply Crowding Distance (if needed)}
				\STATE \hspace{1em} If the last front $F_k$ exceeds capacity $N$, select individuals based on crowding distance to preserve diversity
				\ENDFOR
				\STATE \textbf{Output:} Final Pareto-optimal set $P_T$
			\end{algorithmic}
		\end{algorithm}
		%
		The key to its effectiveness is a hierarchical selection strategy that relies on two primary mechanisms: the Fast Non-Dominated Sorting (FNDS) procedure for convergence, and the Crowding Distance calculation for diversity preservation.
		
	\end{appendices}
	%%%%%%%%%%%%%%%%%%%%%%%%%%%%%%%%%%%%%%%%%%%%%%
	\bibliography{references}   % <-- your .bib filename
\end{document}